\documentclass[11pt,a4paper]{article}

\usepackage[T1]{fontenc}
\usepackage[utf8]{inputenc}
\usepackage{lmodern}
\usepackage[margin=1in]{geometry}
\usepackage{amsmath,amssymb,amsthm}
\usepackage{booktabs}
\usepackage{array}
\newcolumntype{L}{>{\raggedright\arraybackslash}p{0.62\textwidth}}
\newcolumntype{R}[1]{>{\raggedright\arraybackslash}p{#1}}
\usepackage{longtable}
\usepackage{listings}
\usepackage{microtype}
\usepackage[table]{xcolor}
\usepackage{placeins}
\usepackage[hidelinks]{hyperref}

% ------------------------------------------------------------------ floats --
% Large tables in a long document jam the float queue: once one cannot be
% placed, every later float waits behind it, and they all end up dumped after
% the bibliography.  Relax the placement fractions so a big table may occupy
% most of a page, and put a barrier at every section so nothing drifts out of
% the section that discusses it.
\renewcommand{\topfraction}{0.92}
\renewcommand{\bottomfraction}{0.85}
\renewcommand{\textfraction}{0.06}
\renewcommand{\floatpagefraction}{0.72}
\setcounter{topnumber}{4}
\setcounter{bottomnumber}{3}
\setcounter{totalnumber}{6}
\let\mvoldsection\section
\renewcommand{\section}{\FloatBarrier\mvoldsection}

\hypersetup{
  pdftitle={The Generalised Collatz Multiverse},
  pdfauthor={RobinCodes},
  pdfsubject={Number theory, discrete dynamical systems},
  pdfkeywords={Collatz conjecture, 3x+1 problem, generalised Collatz maps,
    qx+1, 3x+d, cycles, normal form},
  colorlinks=true, linkcolor=black, citecolor=black, urlcolor=blue!55!black
}

% ---------------------------------------------------------------- theorems --
\theoremstyle{plain}
\newtheorem{theorem}{Theorem}[section]
\newtheorem{lemma}[theorem]{Lemma}
\newtheorem{proposition}[theorem]{Proposition}
\newtheorem{corollary}[theorem]{Corollary}
\theoremstyle{definition}
\newtheorem{definition}[theorem]{Definition}
\newtheorem{remark}[theorem]{Remark}
\newtheorem{example}[theorem]{Example}
\newtheorem{heuristic}[theorem]{Heuristic}
\newtheorem{conjecture}[theorem]{Conjecture}

% ---------------------------------------------------------------- colours ---
\definecolor{gCon} {HTML}{DCE9F7}   % a = 0
\definecolor{gShp} {HTML}{DFF3E0}   % b = 0, power of two
\definecolor{gRun} {HTML}{F6DDDD}   % runaway
\definecolor{gHal} {HTML}{EFEAF7}   % halving-reducible
\definecolor{gAtt} {HTML}{FBEEDA}   % attraction-reducible
\definecolor{gDes} {HTML}{E4F2F4}   % a = 1 descent
\definecolor{gOpen}{HTML}{FFF6A8}   % open
\newcommand{\cell}[2]{\cellcolor{#1}\textbf{\footnotesize #2}}
\newcommand{\solved}[1]{\textcolor{green!45!black}{\textsc{solved} --- #1}}
\newcommand{\reduc}[1]{\textcolor{orange!70!black}{\textsc{reduces} --- #1}}
\newcommand{\openv}[1]{\textcolor{red!70!black}{\textsc{#1}}}
\newcommand{\ok}{\textcolor{green!45!black}{\textbf{pass}}}
\newcommand{\bad}{\textcolor{red!70!black}{\textbf{fail}}}
\newcommand{\rn}{{\color{red!60!black}\textsc{run}}}

% -------------------------------------------------------------- listings ----
\lstdefinestyle{py}{
  basicstyle=\footnotesize\ttfamily, keywordstyle=\color{blue!60!black},
  commentstyle=\color{gray}\itshape, stringstyle=\color{green!40!black},
  numbers=left, numberstyle=\tiny\color{gray}, numbersep=8pt,
  frame=leftline, framesep=6pt, rulecolor=\color{gray!40},
  showstringspaces=false, breaklines=true, language=Python, xleftmargin=1.2em
}

% ---------------------------------------------------------------- macros ----
\newcommand{\N}{\mathbb{N}}
\newcommand{\Np}{\mathbb{N}_{\geq 1}}
\newcommand{\Z}{\mathbb{Z}}
\newcommand{\Tab}{T_{a,b}}
\newcommand{\vp}{v_p}
\newcommand{\vtwo}{v_2}
\newcommand{\Cyc}{\mathrm{Cyc}}
\newcommand{\PCyc}{\mathrm{Cyc}^{\ast}}
\newcommand{\oddpart}{\mathrm{odd}}
\newcommand{\normalform}{\mathrm{nf}}

\begin{document}

% =============================================================== title ======
\title{\bf The Generalised Collatz Multiverse\\[2mm]
\large Symmetries, reductions and fates of the maps
$n \mapsto n/2$ (even), $n \mapsto an+b$ (odd)}

\author{RobinCodes\thanks{This paper, the accompanying software and the
verification suite were produced in collaboration with an artificial-intelligence
system; see the note on the use of artificial intelligence in
Section~\ref{sec:ai}.}\\[1mm]
\small\href{https://github.com/RobinCodes/Collatz}{\texttt{github.com/RobinCodes/Collatz}}}

\date{18 August 2026}

\maketitle

\begin{abstract}
\noindent
For integers $a,b \geq 0$ let $\Tab$ send $n$ to $n/2$ when $n$ is even and to
$an+b$ when $n$ is odd, so that $T_{3,1}$ is the Collatz map. We give a
structure theory for this two-parameter family, over $\Z$.

Two symmetries and one dichotomy account for all of it. The pair $(2a,2b)$ is
a time change of $(a,b)$; the map $n \mapsto dn$ embeds $(a,b)$ into $(a,db)$
for odd $d$; and for each odd prime $p$ the quantity $\min(v_{p}(n),v_{p}(b))$
either climbs to $v_{p}(b)$ and sticks, when $p \mid a$, or is exactly
conserved, when $p \nmid a$. We show these are the \emph{only} affine
equivalences, so the multiplier $a$ can never change; that they form a
terminating, confluent rewriting system whose normal form is always coprime;
and that after reduction every universe is one of four families we solve
completely, or else has $a \geq 3$ odd, $b$ odd and $\gcd(a,b)=1$, where the
problem is open. Exactly $1/3$ of all universes are open, and $2/\pi^{2}$ are
open and irreducible; the rest are either settled or are a smaller universe in
disguise.

The conserved valuations partition a coprime universe into $\tau(b)$
mutually unreachable strata indexed by the divisors of $b$, which decomposes
its cycle set exactly. We also determine when $1$ itself is periodic, and
sharpen the row $a=1$ enough to make its cycle census exact rather than a
lower bound. Two further chapters ask what survives modulo $q$ rather than
modulo $2$, and what survives when the branch is chosen by a property that is
not a congruence at all; in the second case none of this paper's machinery
survives, and we argue that this, rather than any difficulty of degree, is
what separates $3x+1$ from its neighbours.

The framing and many of the phenomena are due to carykh and to Artem, who
found them experimentally; we attribute each result and correct three
conjectures, with counterexamples in $2x+1$, $5x+1$ and $3x+7$. Every theorem
is accompanied by a machine-run check, over $3 \times 10^{6}$ cases in total.
\end{abstract}

% =============================================================== intro ======
\section{Introduction}

The Collatz conjecture concerns the map that halves even numbers and sends an
odd $n$ to $3n+1$. Its resistance to proof is famous, and Erd\H{o}s'
assessment --- that mathematics is not ready for such problems --- is still
quoted. Less often asked is what happens when the constants $3$ and $1$ are
changed. Writing
\[
  \Tab(n) \;=\;
  \begin{cases}
    n/2 & n \text{ even},\\
    a n + b & n \text{ odd},
  \end{cases}
  \qquad a \in \Z_{\geq 0},\ b \in \Z,
\]
one obtains a two-parameter family --- a \emph{multiverse}, in the coinage of
the YouTube mathematician carykh~\cite{carykh}, whose video animates the
$6 \times 6$ grid $0 \leq a, b \leq 5$ and observes that most of its cells are
in some way degenerate. A follow-up ramble~\cite{lazykh} and a reply video by
Artem~\cite{artem} push the observations further, in particular into the
row $a = 3$.

Those three sources are experimental. They exhibit striking regularities ---
that even-indexed cells duplicate their halves, that $3x+3$ is a scaled copy
of $3x+1$ but $3x+5$ is not, that $3x+9$ takes two odd steps to become a copy
where $3x+3$ takes one, that composite offsets contain a copy of every divisor
--- but they state them as patterns observed rather than as theorems, and they
leave the boundary between what is solved and what is open unmapped. This
paper draws that boundary.

\paragraph{What is proved here.}
\begin{enumerate}\itemsep3pt
\item \textbf{Rigidity} (Theorem~\ref{thm:rigid}). The only affine maps
      intertwining two members of the family are $n \mapsto dn$ with $d$ an
      odd integer, and they satisfy $a' = a$, $b' = db$. So the multiplier $a$
      is a complete invariant of a row: no affine change of variable will ever
      relate $3x+1$ to $5x+1$. In particular the relation between $(2a,2b)$
      and $(a,b)$ --- carykh's ``even parallel universes'' --- is \emph{not} a
      conjugacy; it is a time change, and we give its exact form
      (Proposition~\ref{prop:halve}).
\item \textbf{$p$-adic stratification} (Theorem~\ref{thm:strat}). For an odd
      prime $p$, the behaviour of $\vp$ along an orbit is governed entirely by
      whether $p$ divides $a$. If it does, $\min(\vp(n),\vp(b))$ strictly
      increases at every odd step until it reaches $\vp(b)$, so
      $p^{\vp(b)}\Z$ is a global attractor; if it does not,
      $\min(\vp(n),\vp(b))$ is exactly conserved, so the state space splits.
      This is the theorem behind Artem's connected/disconnected dichotomy and
      behind carykh's ``factor loyalty''.
\item \textbf{Normal form} (Theorem~\ref{thm:nf}). Halving and attraction form
      a terminating, confluent rewriting system; the normal form
      $\normalform(a,b)$ is unique and always satisfies $\gcd = 1$. Both rules
      preserve the fate of every orbit and act bijectively on cycles.
\item \textbf{Master classification} (Theorem~\ref{thm:master}). Every
      universe is exactly one of: constant ($a=0$), pure ($b=0$), parity
      runaway ($a+b$ odd), descent ($a=1$), or primitive. The first four are
      solved here, completely and with explicit cycle data. The fifth is open.
\item \textbf{Densities} (Theorem~\ref{thm:density}). The open universes have
      density exactly $1/3$; the open \emph{irreducible} ones have density
      $2/\pi^{2}$.
\item \textbf{Stratification and cycle decomposition}
      (Theorem~\ref{thm:strata}, Corollary~\ref{cor:cycdec}). A coprime
      universe partitions into $\tau(b)$ invariant strata indexed by the
      divisors of $b$, and its cycle set is the corresponding disjoint union
      $\bigsqcup_{d \mid b} d \cdot \PCyc(T_{a,b/d})$.
\item \textbf{Over $\Z$} (Proposition~\ref{prop:Zsplit}). The positive
      integers are invariant exactly when $a+b \geq 1$ and the negative ones
      exactly when $b \leq a-1$; sign changes only at odd $n$ with
      $|n| \leq |b|/a$, so every orbit changes sign at most once. The negative
      half of $3x+1$ is the positive half of $3x-1$.
\item \textbf{The cycle through $1$} (Theorem~\ref{thm:one}). It is periodic
      exactly when its forward orbit meets a power of two; $a+b = 2^{k}$ is
      the one-step case, and accounts for $93$ of the $147$ instances we
      found.
\item \textbf{Two chapters outward} (Sections~\ref{sec:qmulti}
      and~\ref{sec:beyond}). Modulo $q$: both symmetries survive, the parity
      obstruction becomes a residue graph, the threshold $a<4$ becomes
      $\prod a_{i} < q^{q}$, and rigidity fails. Beyond congruences: the
      theory applies to eventually periodic branch predicates and to nothing
      else, which we argue is the real reason $3x+1$ is tractable at all.
\end{enumerate}

\paragraph{Software.}
Everything computational here is produced by the \texttt{Multiverse}
package. Its modules are \texttt{multiverse.py} for the two-parameter
family, \texttt{qary.py} for Sections~\ref{sec:qmulti}
and~\ref{sec:beyond}, and \texttt{literature.py} for
Appendix~\ref{app:lit}; \texttt{generate\_data.py} drives them, and
\texttt{Collatz\_Program.py} is the interactive front end. All of it lives in
the repository~\cite{repo}.
Section~\ref{sec:software} describes them and Section~\ref{sec:verif} reports
the verification suite. No number in this paper was transcribed by hand.

\subsection{Related work}
\label{sec:related}

The $3x+1$ problem itself has an enormous literature, surveyed by
Lagarias~\cite{lagarias,lagariasbook} and catalogued in his annotated
bibliographies~\cite{lagbib1,lagbib2,lagbib3}. We use only a few landmarks:
Terras~\cite{terras} showed almost every $n$ has finite stopping time;
Eliahou~\cite{eliahou} constrained the possible periods of a nontrivial
cycle; Tao~\cite{tao} showed almost all orbits attain almost bounded values;
Hercher~\cite{hercher} ruled out $m$-cycles for $m \leq 91$; and
Barina~\cite{barina} pushed numerical verification past $2^{71}$.
Gon\c{c}alves, Greenfeld and Madrid~\cite{ggm} extended Tao's method to a
class of generalized maps.

Generalisation of the constants is older than it looks.
Crandall~\cite{crandall} introduced the $qx+1$ maps in 1978 and devoted a
section to the general $qx+r$; his conjecture --- that for odd $q > 3$ the map
$qx+1$ has an orbit not reaching $1$ --- was proved for almost all $q$ in the
sense of asymptotic density by Franco and Pomerance~\cite{fp}.
Santos~\cite{santos} treats $qn+1$ extensively, establishing finiteness of the
periodic orbits and settling $q = 2^{k}-1$.
Belaga and Mignotte~\cite{bm} study $3x+d$, which is the row $a = 3$ of our
grid. Matthews~\cite{matthews} develops a different generalisation, with
several branches and a Markov-chain/ergodic analysis, closer to the $q$-ary
map of Section~\ref{sec:qary} than to the two-parameter family here.

Two facts from that literature bound what this paper can hope to do. First,
Conway~\cite{conway} showed that generalized Collatz functions are universal,
and Kurtz and Simon~\cite{kurtzsimon} that the associated reachability problem
is $\Pi^{0}_{2}$-complete. Their setting is broader than ours --- Conway
allows rational multipliers subject to integrality, which the fixed-modulus
integer shape of Section~\ref{sec:qary} does not --- so this does not show
that our $q$-ary maps are universal. What it does show is that no
classification theorem can cover generalized Collatz functions in the wide
sense, which is the ceiling any programme of this kind runs into. Second, Franco and Pomerance proved Crandall's conjecture --- that $qx+1$
admits an orbit never reaching $1$ --- for almost all odd $q$. So along the
column $b = 1$ the rows $a \geq 5$ are not merely unproven: for almost all
such $a$ the Collatz-style statement is known to be \emph{false}. Our contribution is orthogonal to both: it is a structure
theory for the two-parameter slice, saying which of its members are the same
problem and which are genuinely distinct.

Finally, the multiverse framing that prompted this paper is due to
carykh~\cite{carykh,lazykh} and Artem~\cite{artem}, and the community site
collatzresearch.org~\cite{cresearch} collects related experimental work.
Section~\ref{sec:attrib} attributes each individual result.

\subsection{Prior work on the generalised maps}
\label{sec:priorgen}

Because this paper claims a structure theory for a family that has been
studied for fifty years, we owe the reader a precise account of what was
already known. We surveyed the corpus catalogued by
ccchallenge.org~\cite{ccc} --- $375$ papers at the time of writing --- and
classified it (Appendix~\ref{app:lit}); $49$ of those papers concern
generalised maps. The following are the ones that overlap this paper.

\paragraph{The column $b = 1$.}
Crandall~\cite{crandall} posed the $qx+1$ problem;
Steiner~\cite{steiner1981a,steiner1981b} studied $Qx+1$ for odd $Q$ and
Franco~\cite{franco1990} connected the cycle problem to Diophantine
approximation, which is the analytic content behind our
Theorem~\ref{thm:cyceq}: a cycle needs $2^{m}/a^{k}$ close to $1$, hence
needs a good rational approximation to $\log_{2}a$. Franco and
Pomerance~\cite{fp} then proved Crandall's conjecture for almost all $q$ in
the sense of asymptotic density, and Santos~\cite{santos} established
finiteness of the cycle set for each $q$ and settled $q = 2^{k}-1$. Our
row $b=1$ (Table~\ref{tab:colb1}) is a census inside their setting, not a
competitor to it.

\paragraph{The row $a = 3$.}
This is the closest prior work, and it anticipates part of
Section~\ref{sec:nf}. Belaga and Mignotte~\cite{bm} embed $3x+1$ in the
$3x+d$ family and --- crucially --- restrict throughout to $d$ coprime to
$6$, that is, $d$ odd and not divisible by $3$. That restriction is exactly
our normal form specialised to $a = 3$: by Theorem~\ref{thm:nf} the even $d$
reduce by halving and the multiples of $3$ reduce by attraction, so nothing
is lost. They state the \emph{$3x+d$ conjecture} --- that each such $d$ admits
only finitely many cycles --- and their Proposition~3.1, that the fixed
points of the accelerated map are $n = d/(2^{k}-3)$, is the $k=1$ case of the
cycle equation. Belaga~\cite{belaga2003} gives polynomial upper bounds on the
number and size of $(3x+d)$-cycles of given odd length, and Belaga and
Mignotte~\cite{bm2006} tabulate the \emph{primitive} cycles of $3x+d$ ---
their term, and their notion agrees with ours: a cycle no element of which
shares a factor with $d$. Simons~\cite{simons2008} bounds $m$-cycles for
generalised Syracuse sequences.

We have not found the divisor decomposition (Corollary~\ref{cor:cycdec})
stated as a general theorem anywhere, though it is visible in the examples
of~\cite{bm}: their nontrivial fixed points of $S_{65}$ are $5$ and $13$,
which in our language are $5$ times the fixed point of $S_{13}$ and $13$
times that of $S_{5}$. What we add is that these copies \emph{exhaust} the
cycle set, and that the same decomposition holds for the whole phase space
(Theorem~\ref{thm:strata}).

\paragraph{The $q$-ary maps.}
Section~\ref{sec:qmulti} is the elementary layer of a subject with a
substantial theory. Matthews and Watts~\cite{mw1984,mw1985} generalised
Hasse's version of the Syracuse algorithm to exactly the shape of
Section~\ref{sec:qary} and analysed it with Markov chains;
Leigh~\cite{leigh1986} and Buttsworth and Matthews~\cite{bmatthews1990}
developed the Markov-matrix machinery, and Matthews~\cite{matthews,matthews2005}
surveys it. Venturini~\cite{venturini1992,venturini1997} studies the same
maps as \emph{periodically linear} functions, which is the terminology of our
Proposition~\ref{prop:congpred}. Mignosi~\cite{mignosi1995} and
Brocco~\cite{brocco1995} treat a related generalisation, and Rouet and
Feix~\cite{rouetfeix} construct cycles and give a stochastic model.
Matthews and Leigh~\cite{mleigh1987} and Hicks et al.~\cite{hicks2008} move
to polynomial rings, where the analogue of the conjecture is provable. Our
contribution at $q \geq 3$ is only Propositions~\ref{prop:qscale}
and~\ref{prop:qtwist} and Theorem~\ref{thm:resobs}, which we have not found
stated in this form, and the drift threshold $q^{q}$.

\paragraph{Undecidability.}
Conway~\cite{conway} and Kurtz and Simon~\cite{kurtzsimon} bound the whole
programme from above, as discussed above and in Section~\ref{sec:beyond}.

% ============================================================ definitions ===
\section{The family}
\label{sec:family}

\begin{definition}\label{def:map}
For $a \in \Z_{\geq 0}$ and $b \in \Z$ define $\Tab : \Z \to \Z$ by
$\Tab(n) = n/2$ for $n$ even and $\Tab(n) = an+b$ for $n$ odd. We call the
pair $(a,b)$ a \emph{universe} and write it $ax+b$. The phase space is $\Z$;
Proposition~\ref{prop:Zsplit} says when each half-line is invariant, and
Proposition~\ref{prop:neg} reduces $b < 0$ to $b > 0$, so the classification
of Section~\ref{sec:solved} is stated for $b \geq 0$ without loss.
\end{definition}

\subsection{The general \texorpdfstring{$q$}{q}-ary map}
\label{sec:qary}

Definition~\ref{def:map} is the case $q = 2$ of a wider family. Fix a modulus
$q \geq 1$ and coefficients $a_{1},\dots,a_{q-1}$, $b_{1},\dots,b_{q-1}$, and
define
\[
  F(n) \;=\;
  \begin{cases}
    a_{i}n + b_{i} & n \equiv i \pmod q,\quad 1 \leq i \leq q-1,\\[1mm]
    n/q            & n \equiv 0 \pmod q .
  \end{cases}
\]
For $q = 1$ every $n$ is divisible by $q$ and $F$ is the identity, so the
family is trivial. Sections~\ref{sec:sym}--\ref{sec:open} study $q = 2$, the
case with a single affine branch; we write $\Tab$ for it and call $a$ the
\emph{multiplier} and $b$ the \emph{offset}.

Most of that material is specific to $q = 2$, because it uses the fact that
the parity of $an+b$ is forced by those of $a$, $b$ and $n$. Some of it is
not, and Section~\ref{sec:qmulti} returns to general $q$ to say exactly which
parts survive: both symmetries do, with $2$ replaced by $q$ and ``odd''
replaced by ``coprime to $q$''; the parity obstruction becomes a statement
about a finite directed graph on $\Z/q$; and rigidity fails outright.

\subsection{The phase space over \texorpdfstring{$\Z$}{Z}}
\label{sec:overZ}

It is natural to take $\Np$ as the phase space, and much of the literature
does. But $\Tab$ is defined on all of $\Z$, and restricting to the positive
integers hides structure --- most conspicuously the three known cycles of
$3x-1$, which are the negative-side cycles of the Collatz map itself. We
therefore work over $\Z$ and record exactly when the sign is preserved.

\begin{proposition}[Structure of $\Z$]\label{prop:Zsplit}
Let $a \geq 1$ and $b \in \Z$.
\begin{enumerate}\itemsep3pt
\item[(i)] $\Tab(0) = 0$, and $\Tab(n) = 0$ for $n \neq 0$ exactly when $n$ is
      odd with $n = -b/a$. So $\{0\}$ is a fixed point with at most one other
      preimage.
\item[(ii)] $\Z_{\geq 1}$ is $\Tab$-invariant if and only if $a + b \geq 1$.
\item[(iii)] $\Z_{\leq -1}$ is $\Tab$-invariant if and only if $b \leq a - 1$.
\item[(iv)] In all cases the sign can change only at an \emph{odd} $n$ with
      $|n| \leq |b|/a$. Hence at most $\lceil |b|/a \rceil$ integers of each
      sign ever cross, and every orbit is eventually confined to one of
      $\Z_{\geq 1}$, $\{0\}$, $\Z_{\leq -1}$.
\end{enumerate}
\end{proposition}

\begin{proof}
(i) $0$ is even and $0/2 = 0$. If $n \neq 0$ maps to $0$ then $n$ cannot be
even (as $n/2 = 0$ forces $n=0$), so $n$ is odd and $an + b = 0$.

(ii) Halving preserves $\Z_{\geq 1}$, so invariance can fail only on the odd
branch. For odd $n \geq 1$ the quantity $an + b$ is increasing in $n$, so its
minimum over odd $n \geq 1$ is $a + b$, attained at $n = 1$. Thus
$an + b \geq 1$ for all odd $n \geq 1$ iff $a + b \geq 1$.

(iii) Apply (ii) to $T_{a,-b}$ and use Proposition~\ref{prop:neg}: the
negatives of $\Tab$ are the positives of $T_{a,-b}$, which are invariant iff
$a - b \geq 1$.

(iv) Halving never changes sign, so a change of sign happens at an odd $n$.
If $n \geq 1$ is odd and $an + b \leq 0$ then $an \leq -b$, so $b < 0$ and
$n \leq |b|/a$; symmetrically for $n \leq -1$. That bounds the crossing set,
which is therefore finite.

For the confinement, note that in each case at least one half-line is already
invariant, and crossings can only run \emph{into} it. If $b > 0$ then
$a + b \geq 1$, so $\Z_{\geq 1}$ is invariant by (ii) and every crossing goes
from negative to positive; if $b < 0$ then $b \leq a-1$, so $\Z_{\leq -1}$ is
invariant by (iii) and every crossing goes the other way; and if $b = 0$ then
$an$ has the sign of $n$ and there are no crossings at all. In each case an
orbit changes sign at most once, after which it is confined.
\end{proof}

\begin{remark}
Proposition~\ref{prop:Zsplit} is why the paper can state its theorems over
$\Z$ without a case explosion. For $b \geq 0$ --- the normalisation we use
after Proposition~\ref{prop:neg} --- part (ii) holds automatically whenever
$a \geq 1$, so the positive integers are invariant; part (iii) then says the
negatives are invariant precisely when $b < a$, and otherwise finitely many
negative integers leak into $\Z_{\geq 0}$. The classical case $b = 1 < 3 = a$
has both half-lines invariant, which is why $3x+1$ and $3x-1$ are usually
treated as separate problems. They are the two halves of one map.
\end{remark}

\begin{definition}[Fates]\label{def:fate}
An orbit is \emph{eventually periodic} if it repeats a value, and
\emph{divergent} if $|{\cdot}|$ is unbounded along it. Fix an invariant set
$S \subseteq \Z$ --- in practice $\Np$ or $\Z_{\leq -1}$, by
Proposition~\ref{prop:Zsplit}. The \emph{fate of $\Tab$ on $S$} is
\textsc{cyclic} if every orbit from $S$ is eventually periodic,
\textsc{divergent} if every orbit from $S$ is divergent, and \textsc{open}
otherwise or when unknown. An unqualified ``fate'' means the fate on $\Np$;
by Proposition~\ref{prop:neg} the fate on $\Z_{\leq -1}$ is the fate of
$T_{a,-b}$ on $\Np$, so the two half-lines are covered by classifying
$(a,b)$ and $(a,-b)$. We write $\Cyc(\Tab)$ for the set of cycles in $S$,
each normalised to begin at its least element.
\end{definition}

Two trivialities are used constantly and recorded once.

\begin{lemma}\label{lem:reachodd}
Let $n \in \Z$ with $n \neq 0$. Then the $\Tab$-orbit of $n$ contains
infinitely many odd terms, or reaches $0$. If moreover $a \geq 1$ and
$a + b \geq 1$ and $n \geq 1$, then $\Tab^{k}(n) \geq 1$ for all $k$.
\end{lemma}

\begin{proof}
Halving strictly decreases $|n|$ for $n \neq 0$, so no orbit can consist of
even nonzero terms from some point on: after at most $\vtwo(n)$ halvings an
odd value or $0$ appears, and the same applies to every tail. The positivity
claim is Proposition~\ref{prop:Zsplit}(ii).
\end{proof}

\begin{lemma}[Valuations]\label{lem:val}
Let $p$ be an odd prime and $n \neq 0$. Then $\vp(n/2) = \vp(n)$ whenever $n$
is even, and
\[
  \vp(an+b) \;\geq\; \min\bigl(\vp(a) + \vp(n),\; \vp(b)\bigr),
\]
with equality whenever $\vp(a)+\vp(n) \neq \vp(b)$. \textup{(}Here
$\vp(0) := +\infty$.\textup{)}
\end{lemma}

\begin{proof}
The first claim holds because $p$ is odd. The second is the ultrametric
inequality for $\vp$ together with $\vp(an) = \vp(a)+\vp(n)$.
\end{proof}

% ============================================================== symmetries ==
\section{The symmetries}
\label{sec:sym}

\subsection{Halving: carykh's even parallel universes}

\begin{proposition}[Halving]\label{prop:halve}
Let $U = T_{2a,2b}$ and $T = \Tab$, with $(a,b) \neq (0,0)$. Then for every
$n$,
\[
  n \text{ even} \;\Longrightarrow\; U(n) = T(n),
  \qquad
  n \text{ odd} \;\Longrightarrow\; U(n) = 2\,T(n) \text{ is even and }
  U^{2}(n) = T(n).
\]
Consequently the $U$-orbit of $n$ is obtained from the $T$-orbit of $n$ by
inserting the single extra term $2\,T(x)$ immediately after each odd term $x$.
In particular:
\begin{enumerate}\itemsep2pt
\item[(i)] $n$ is $U$-eventually-periodic iff it is $T$-eventually-periodic,
      and $U$-divergent iff $T$-divergent; the two universes have the same
      fate.
\item[(ii)] $C \mapsto C \cup \{\,2\,T(x) : x \in C \text{ odd}\,\}$ is a
      bijection $\Cyc(T) \to \Cyc(U)$, and the corresponding periods satisfy
      $|C_{U}| = |C_{T}| + k(C_{T})$, where $k$ counts the odd terms of $C_T$.
\end{enumerate}
\end{proposition}

\begin{proof}
For even $n$ both maps halve. For odd $n$, $U(n) = 2an+2b = 2(an+b) = 2T(n)$,
which is even, so $U$ halves it next and $U^{2}(n) = an+b = T(n)$. The
description of the orbit follows by induction, and (i) is immediate because
the two orbits have the same set of values --- the inserted terms are
determined by their predecessors --- and are simultaneously finite or
unbounded. For (ii), every $U$-cycle contains an odd term by
Lemma~\ref{lem:reachodd}, and starting from such a term the displayed
identities read off the $T$-cycle; conversely each $T$-cycle inflates to a
$U$-cycle. For the period: starting at $x \in C_{T}$, the insertion
description returns to $x$ after exactly $|C_{T}| + k(C_{T})$ steps of $U$,
and no sooner --- deleting the inserted terms from a shorter $U$-return would
give a shorter $T$-return, contradicting the minimality of $|C_{T}|$. Since a
periodic orbit of a map has exactly as many distinct elements as its period,
$|C_{U}| = |C_{T}| + k(C_{T})$, and in particular the inserted terms are
automatically distinct from the elements of $C_{T}$.
\end{proof}

Proposition~\ref{prop:halve} is carykh's observation that $4x+2$ ``is just
$2x+1$ with more steps''~\cite{carykh}, made exact. Note what it is
\emph{not}: the two maps are not conjugate, because $U$ takes strictly more
steps. Theorem~\ref{thm:rigid} below shows that no conjugacy could exist.

\subsection{Odd scaling: carykh's odd parallel universes}

\begin{proposition}[Odd scaling]\label{prop:scale}
Let $d$ be an odd integer. Then for every $n$,
\[
  T_{a,\,db}(dn) \;=\; d \cdot \Tab(n).
\]
Hence $d\Z$ is $T_{a,db}$-invariant and $n \mapsto dn$ is an isomorphism of
dynamical systems from $\Tab$ onto $T_{a,db}$ restricted to $d\Z$.
\end{proposition}

\begin{proof}
Since $d$ is odd, $dn$ and $n$ have the same parity. If $n$ is even,
$T_{a,db}(dn) = dn/2 = d(n/2)$. If $n$ is odd,
$T_{a,db}(dn) = a\,dn + db = d(an+b)$.
\end{proof}

\begin{example}
$T_{5,5}$ restricted to $5\Z$ is a copy of $T_{5,1}$. carykh observed that the
two have the identical ``cycle fingerprint'' of periods $7, 10, 10$; by
Proposition~\ref{prop:scale} the cycles of $T_{5,5}$ that lie in $5\Z$ are
exactly $5$ times those of $T_{5,1}$, so this is forced. Our census confirms
that \emph{all three} cycles of $T_{5,5}$ lie in $5\Z$, which is a separate
fact --- it is Theorem~\ref{thm:strat}(i), since $5 \mid \gcd(5,5)$.
\end{example}

\subsection{Negation}

\begin{proposition}[Negation duality]\label{prop:neg}
For all $a, b$ and all $n$, $\;T_{a,-b}(n) = -\,\Tab(-n)$. Hence $n \mapsto -n$
conjugates $\Tab$ on $\Z_{\leq -1}$ to $T_{a,-b}$ on $\Np$.
\end{proposition}

\begin{proof}
$n$ and $-n$ have the same parity. For $n$ even both sides are $n/2$; for $n$
odd, $-\Tab(-n) = -(-an+b) = an - b = T_{a,-b}(n)$.
\end{proof}

So the negative half of the classical Collatz map is the positive half of
$3x-1$, whose three known cycles begin at $1$, $5$ and $17$. Everything below
is stated for $b \geq 0$; Proposition~\ref{prop:neg} transfers it.

\subsection{Rigidity: there is nothing else}

\begin{theorem}[Rigidity]\label{thm:rigid}
Let $\varphi(n) = \lambda n + \mu$ with $\lambda \neq 0$, let $\varphi$ map
$\Z$ into $\Z$, and suppose
$\varphi \circ \Tab = T_{a',b'} \circ \varphi$ on $\Z$, where
$a,b,a',b' \in \Z$ and $a' \neq 0$. Then
\[
  \mu = 0, \qquad \lambda \text{ is an odd integer}, \qquad
  a' = a, \qquad b' = \lambda b .
\]
\end{theorem}

\begin{proof}
Taking $n = 0$ and $n = 1$ shows $\mu, \lambda + \mu \in \Z$, so
$\lambda, \mu \in \Z$. We rule out the bad parities in turn. Throughout,
``$\varphi(n)$ even'' selects the halving branch of $T_{a',b'}$.

\emph{Case $\lambda$ even.} Then $\varphi(n) \equiv \mu \pmod 2$ for every
$n$, so $T_{a',b'} \circ \varphi$ uses one fixed branch.
If $\mu$ is even, evaluating the intertwining relation at even $n$ gives
$\lambda n/2 + \mu = (\lambda n + \mu)/2$, hence $\mu = 0$; then at odd $n$ it
gives $\lambda(an+b) = \lambda n / 2$ for all odd $n$, forcing $a = 1/2$, which
is not an integer. If $\mu$ is odd, evaluating at even $n$ gives
$\lambda n/2 + \mu = a'(\lambda n + \mu) + b'$ for all even $n$; comparing the
coefficients of $n$ yields $\lambda/2 = a'\lambda$, i.e.\ $a' = 1/2$,
again impossible.

\emph{Case $\lambda$ odd, $\mu$ odd.} Then $\varphi$ reverses parity. At even
$n$ the relation reads $\lambda n /2 + \mu = a'(\lambda n + \mu) + b'$, and
comparing coefficients of $n$ gives $a' = 1/2$: impossible.

\emph{Case $\lambda$ odd, $\mu$ even.} Then $\varphi$ preserves parity. At
even $n$: $\lambda n/2 + \mu = (\lambda n + \mu)/2$, so $\mu = 0$. At odd $n$:
$\lambda(an+b) = a'\lambda n + b'$ for all odd $n$; comparing coefficients
gives $a' = a$ and $b' = \lambda b$.
\end{proof}

\begin{corollary}\label{cor:rigid}
\begin{enumerate}\itemsep2pt
\item[(i)] The only bijective affine conjugacies within the family are the
      identity and $n \mapsto -n$, which gives $\Tab \cong T_{a,-b}$.
\item[(ii)] The multiplier $a$ is an affine invariant: no affine change of
      variable relates two universes in different rows. In particular no such
      change of variable can reduce $5x+1$ to $3x+1$.
\item[(iii)] Within a row, two universes are affinely related only if their
      offsets differ by an odd factor.
\end{enumerate}
\end{corollary}

\begin{proof}
All three are immediate from Theorem~\ref{thm:rigid}, which gives $\mu = 0$,
$\lambda$ odd, $a' = a$ and $b' = \lambda b$. For (i), a bijection of $\Z$ of
the form $n \mapsto \lambda n$ forces $\lambda = \pm 1$, and $\lambda = -1$
gives $b' = -b$.
\end{proof}

Corollary~\ref{cor:rigid}(ii) is worth emphasising, because it is the precise
sense in which the rows of the multiverse are independent problems. Whatever
resolves $3x+1$ need not touch $5x+1$.

\subsection{Modified maps: what generalises, and what does not}
\label{sec:modified}

Propositions~\ref{prop:halve} and~\ref{prop:scale} both have the shape
``post-compose the odd branch with a map $g$ and see what survives''. It is
natural to ask how far that goes. Let $g : \Z \to \Z$ be arbitrary and set
\[
  T_{g}(n) \;=\;
  \begin{cases}
    n/2 & n \text{ even},\\
    g(an+b) & n \text{ odd}.
  \end{cases}
\]

\begin{proposition}[Post-composition]\label{prop:postcomp}
For every odd $n$, $\;T_{g}(n) = g\bigl(\Tab(n)\bigr)$.
\end{proposition}

\begin{proof}
Immediate from the definitions: both sides are $g(an+b)$.
\end{proof}

Proposition~\ref{prop:postcomp} is genuine but shallow --- it says nothing
about the even branch, and it is the even branch that carries the dynamics.
Taking $g(y) = 2y$ recovers $T_{g} = T_{2a,2b}$ and hence
Proposition~\ref{prop:halve}; taking $g(y) = cy$ gives
$T_{g} = T_{ca,\,cb}$.

It is tempting to go further and claim that $g$ itself conjugates $\Tab$ to
$T_{g}$ whenever $g$ preserves $2$-adic valuations --- in particular that for
odd $c$, every cycle $C$ of $\Tab$ lifts to the cycle $cC$ of $T_{ca,cb}$.
That is false, and it fails immediately.

\begin{proposition}[No such conjugacy]\label{prop:noconj}
Let $c$ be odd, $c > 1$. Then $g(y) = cy$ does not in general carry cycles of
$\Tab$ to cycles of $T_{ca,cb}$. For example $C = \{1,4,2\}$ is a cycle of
$T_{3,1}$ and $\vtwo(3x) = \vtwo(x)$ for every $x$, yet $3C = \{3,12,6\}$ is
not a cycle of $T_{9,3}$: indeed $T_{9,3}(3) = 9 \cdot 3 + 3 = 30 \neq 12$.
More generally, a map $\varphi(n) = \lambda n$ intertwines $\Tab$ with
$T_{a',b'}$ only when $a' = a$ \textup{(}Theorem~\ref{thm:rigid}\textup{)},
whereas $T_{ca,cb}$ has multiplier $ca \neq a$.
\end{proposition}

\begin{proof}
The displayed computation is the counterexample. For the general statement,
Theorem~\ref{thm:rigid} forces $a' = a$ for any affine intertwining, and
$ca = a$ only if $c = 1$.
\end{proof}

An exhaustive check over the cycles of $T_{a,b}$ for
$a \in \{3,5,7\}$, $b \in \{1,3,5\}$ and $c \in \{3,5,7,9\}$ finds
\emph{no} case in which $cC$ is a cycle of $T_{ca,cb}$ --- $0$ successes out
of $104$ --- while the same check against $T_{a,cb}$, which scales the offset
only, succeeds in all $104$. The correct statement is
Proposition~\ref{prop:scale}: it is the offset that may be scaled, not the
pair. Rigidity (Theorem~\ref{thm:rigid}) explains why no repair is possible:
$a$ is an affine invariant, so nothing that changes the multiplier can be a
conjugacy.

% ======================================================== stratification ====
\section{The \texorpdfstring{$p$}{p}-adic stratification}
\label{sec:strat}

This section proves the dichotomy that organises the whole odd part of the
multiverse. Fix an odd prime $p$ and write $\alpha = \vp(a)$,
$\beta = \vp(b)$, with the convention $\vp(0) = +\infty$.

\begin{theorem}[Stratification]\label{thm:strat}
Let $a \geq 1$, $b \geq 1$, and let $p$ be an odd prime. Put
$\Phi(n) = \min(\vp(n), \beta)$ for $n \neq 0$. \textup{(}Nothing below uses
the sign of $n$, so the statement holds on all of $\Z \setminus \{0\}$, not
merely on $\Np$.\textup{)}
\begin{enumerate}\itemsep3pt
\item[(i)] \textup{(Attraction.)} If $\alpha \geq 1$ then $\Phi$ is
      non-decreasing along every orbit, and strictly increases at every odd
      step at which $\Phi < \beta$. Consequently every orbit satisfies
      $\Phi = \beta$ --- that is, lies in $p^{\beta}\Z$ --- after at most
      $\beta$ odd steps, and remains there forever.
\item[(ii)] \textup{(Repulsion.)} If $\alpha = 0$ then $\Phi$ is
      \emph{exactly} conserved: $\Phi(\Tab(n)) = \Phi(n)$ for every
      $n \geq 1$.
\end{enumerate}
\end{theorem}

\begin{proof}
Halving does not change $\vp$ (Lemma~\ref{lem:val}), hence does not change
$\Phi$; so only odd steps matter. Let $n$ be odd with $\vp(n) = k$.

(i) Let $\alpha \geq 1$. By Lemma~\ref{lem:val},
$\vp(an+b) \geq \min(\alpha + k, \beta) \geq \min(k+1, \beta)$, so
\[
  \Phi(\Tab(n)) \;\geq\; \min\bigl(\min(k+1,\beta),\,\beta\bigr)
  \;=\; \min(k+1, \beta) .
\]
If $\Phi(n) = \min(k,\beta) < \beta$ then $k < \beta$, so
$\min(k+1,\beta) > k = \Phi(n)$ and $\Phi$ strictly increases; otherwise
$\Phi(n) = \beta$ and $\Phi(\Tab(n)) \geq \min(k+1,\beta) = \beta$, so $\Phi$
stays at $\beta$. Since $\Phi$ takes values in $\{0,1,\dots,\beta\}$ and every
orbit meets odd numbers infinitely often (Lemma~\ref{lem:reachodd}), the value
$\beta$ is reached within $\beta$ odd steps and never left.

(ii) Let $\alpha = 0$, so $\vp(an) = k$. If $k < \beta$ then
$k \neq \beta$, so Lemma~\ref{lem:val} gives $\vp(an+b) = k$ exactly, whence
$\Phi(\Tab(n)) = k = \Phi(n)$. If $k \geq \beta$ then
$\vp(an+b) \geq \min(k,\beta) = \beta$, so $\Phi(\Tab(n)) = \beta = \Phi(n)$.
\end{proof}

\begin{corollary}[Connected versus disconnected]\label{cor:conn}
Let $p$ be an odd prime dividing $b$, and let $\beta = \vp(b) \geq 1$.
\begin{itemize}\itemsep2pt
\item If $p \mid a$, the set $p^{\beta}\Z \cap \Np$ is a global attractor:
      every orbit enters it and stays. On it, $\Tab$ is a copy of
      $T_{a,\,b/p^{\beta}}$ by Proposition~\ref{prop:scale}. The functional
      graph is that copy with a forest of transients attached.
\item If $p \nmid a$, then $\Np$ splits into the $\beta+1$ pairwise
      unreachable invariant sets $\Phi^{-1}(0), \dots, \Phi^{-1}(\beta)$. The
      functional graph has at least $\beta+1$ connected components.
\end{itemize}
\end{corollary}

\begin{proof}
The first case is Theorem~\ref{thm:strat}(i), which puts every orbit into
$p^{\beta}\Z$, together with Proposition~\ref{prop:scale} to identify the
dynamics there. The second is Theorem~\ref{thm:strat}(ii): $\Phi$ is constant
along orbits, so its level sets are invariant and no orbit joins two of them;
each is non-empty because $\Phi(p^{j}) = \min(j,\beta)$ realises every value.
\end{proof}

This is precisely Artem's connected/disconnected dichotomy~\cite{artem}, which
he states for $a = 3$ in the form ``$b$ a multiple of $3$ attracts, otherwise
it repels''. Theorem~\ref{thm:strat} shows the criterion is $p \mid a$ for
each odd prime $p \mid b$ separately, and supplies the number of odd steps
needed --- $\vp(b)$ --- which is Artem's observation that $3x+3$ takes one
odd step and $3x+9$ takes two.

\begin{remark}[``Factor loyalty'']
carykh notes that in $5x+5$, once an orbit acquires a factor of $5$ it never
loses it, and that this happens on the whole diagonal $a = b$ when $a$ is odd.
That is Theorem~\ref{thm:strat}(i) with $\beta = 1$: for odd $a = b$, every
odd prime $p \mid a$ divides both parameters, so $p\Z$ is entered after one
odd step. He also notes correctly that the phenomenon fails on the even
diagonal, e.g.\ $4x+4$; there $p = 2$ is the only common prime, and
Lemma~\ref{lem:val} does not apply to it --- halving strips powers of two, so
$\vtwo$ is not monotone. The even diagonal instead reduces by
Proposition~\ref{prop:halve}.
\end{remark}

% =========================================================== normal form ====
\section{Reduction and normal form}
\label{sec:nf}

\begin{definition}[Rewriting rules]\label{def:rules}
On pairs $(a,b)$ with $a \geq 1$, $b \geq 1$ define
\begin{align*}
  \textbf{(H)} && (a,b) &\longrightarrow (a/2,\; b/2)
     && \text{if } 2 \mid a \text{ and } 2 \mid b, \\
  \textbf{(A)} && (a,b) &\longrightarrow (a,\; b/p^{\vp(b)})
     && \text{if } p \text{ is an odd prime with } p \mid a \text{ and } p \mid b .
\end{align*}
\end{definition}

\begin{theorem}[Normal form]\label{thm:nf}
The rewriting system $\{\textbf{(H)}, \textbf{(A)}\}$ is terminating and
confluent on $\{(a,b) : a \geq 1,\ b \geq 1\}$. Its normal forms are exactly
the pairs with $\gcd(a,b) = 1$, so every such universe has a unique normal
form $\normalform(a,b)$, and
\[
  \normalform(a,b) \;=\; \Bigl(\frac{a}{2^{s}},\ \frac{b}{2^{s} D}\Bigr),
  \qquad
  s = \min(\vtwo(a), \vtwo(b)),
  \quad
  D = \prod_{\substack{p \mid a,\ p \neq 2}} p^{\vp(b)} ,
\]
the product being finite because $a \geq 1$ has finitely many prime divisors.
\end{theorem}

\begin{remark}[Why $a \geq 1$]\label{rem:azero}
The restriction is not cosmetic. Every prime divides $0$, so at $a = 0$ the
product defining $D$ is infinite and rule \textbf{(A)} is applicable for every
odd prime at once. One could adopt a convention making $(0,b)$ reduce to
$(0,1)$, but there is no reason to: $\gcd(0,b) = b$, so the coprime normal
forms in that row would be only $(0,1)$, and the row $a = 0$ is already solved
outright by Proposition~\ref{prop:const} without any reduction. We therefore
exclude it, and Theorem~\ref{thm:master} treats $a' = 0$ as its own case.
\end{remark}

\begin{proof}
\emph{Termination.} Rule \textbf{(H)} decreases $\min(\vtwo(a),\vtwo(b))$ by
one and leaves $\vp$ unchanged for every odd $p$; rule \textbf{(A)} removes
one odd prime from $\gcd(a,b)$ --- after it, $\vp(b) = 0$ --- and changes no
power of two, since $p^{\vp(b)}$ is odd. So the pair
$\bigl(\min(\vtwo(a),\vtwo(b)),\ \#\{p \text{ odd prime} : p \mid \gcd(a,b)\}\bigr)$
strictly decreases in the product order at every step, and both coordinates
are non-negative integers.

\emph{Local confluence.} Two applicable rules act on different primes, and
each rule changes only the valuation at its own prime; so any two distinct
applicable steps commute. With termination, Newman's lemma gives confluence
and uniqueness of normal forms.

\emph{Normal forms are the coprime pairs.} If $\gcd(a,b) > 1$ then some prime
divides both, and either $p = 2$, enabling \textbf{(H)}, or $p$ is odd,
enabling \textbf{(A)}; so a normal form is coprime. Conversely neither rule
applies to a coprime pair. Finally, the displayed formula is the result of
applying \textbf{(H)} exactly $s$ times and \textbf{(A)} once for each odd
prime $p \mid a$ with $p \mid b$; the result is coprime, because for such a
$p$ we get $\vp(b/(2^{s}D)) = \vp(b) - \vp(b) = 0$, and
$\min(\vtwo(a/2^{s}), \vtwo(b/2^{s})) = 0$ by choice of $s$.
\end{proof}

\begin{theorem}[Reduction preserves the fate]\label{thm:redfate}
Both rules preserve the fate of a universe, and act bijectively on cycles:
\begin{itemize}\itemsep2pt
\item Under \textbf{(H)}, cycles correspond as in
      Proposition~\ref{prop:halve}\,(ii), with periods increased by the number
      of odd terms.
\item Under \textbf{(A)} with $D = p^{\vp(b)}$, the map $C \mapsto C/D$ is a
      bijection $\Cyc(\Tab) \to \Cyc(T_{a,b/D})$ preserving periods.
\end{itemize}
Consequently $\Tab$ and $T_{\normalform(a,b)}$ have the same fate.
\end{theorem}

\begin{proof}
The halving case is Proposition~\ref{prop:halve}. For \textbf{(A)}: by
Theorem~\ref{thm:strat}(i) every orbit --- in particular every cycle --- lies
eventually, hence entirely in the case of a cycle, inside $D\Z$. By
Proposition~\ref{prop:scale}, $n \mapsto n/D$ carries $\Tab|_{D\Z}$
isomorphically onto $T_{a,b/D}$, so it carries cycles to cycles bijectively
and preserves periods. For fates: given $n \geq 1$, its $\Tab$-orbit has a
tail inside $D\Z$, and dividing that tail by $D$ gives a $T_{a,b/D}$-orbit;
the tail is eventually periodic (resp.\ unbounded) iff the whole orbit is, and
iff its image is. Every $m \geq 1$ arises as such an image, namely from
$n = Dm$, so the universal quantifiers match.
\end{proof}

\begin{example}
$3x+9 \xrightarrow{\ \textbf{(A)},\,p=3\ } 3x+1$: the offset loses $3^{2}$ at
once, and Theorem~\ref{thm:strat}(i) says the absorption into $9\Z$ takes two
odd steps. $3x+15 \xrightarrow{\ \textbf{(A)},\,p=3\ } 3x+5$, since
$v_{3}(15) = 1$; the prime $5$ is repelling and survives. $4x+4
\xrightarrow{\textbf{(H)}} 2x+2 \xrightarrow{\textbf{(H)}} 1x+1$. All three
are observations of~\cite{artem} and~\cite{carykh}, here derived.
\end{example}

\subsection{One-step absorption and the radical criterion}
\label{sec:radical}

Rule \textbf{(A)} says the offset is stripped one prime at a time, over
$\vp(b)$ odd steps. There is a cruder condition under which the whole offset
is stripped in a \emph{single} step, and it is worth isolating because it is
the one that is easy to test.

\begin{proposition}[One-step absorption]\label{prop:onestep}
Let $a, b \geq 1$. Then $an + b \equiv 0 \pmod{b}$ for every $n$ if and only
if $b \mid a$. In that case every orbit lies in $b\Z$ after a single odd step,
and $\Tab$ restricted to $b\Z$ is a copy of $T_{a,1}$.
\end{proposition}

\begin{proof}
$an + b \equiv an \pmod b$, so the condition is $an \equiv 0 \pmod b$ for all
$n$, which at $n = 1$ gives $b \mid a$; conversely if $a = kb$ then
$an + b = b(kn + 1)$. The last claim is Proposition~\ref{prop:scale} with
$d = b$.
\end{proof}

\begin{theorem}[Radical criterion]\label{thm:radical}
Let $a \geq 1$ and let $b \geq 1$ be odd. Then $\normalform(a,b) = (a,1)$ ---
equivalently, $\Tab$ is, on a global attractor, a scaled copy of $T_{a,1}$ ---
if and only if
\[
  \operatorname{rad}(b) \mid a,
\]
where $\operatorname{rad}(b) = \prod_{p \mid b} p$ is the radical of $b$.
\end{theorem}

\begin{proof}
By Theorem~\ref{thm:nf}, $\normalform(a,b) = (a, b/D)$ with
$D = \prod_{p \mid a,\, p \neq 2} p^{\vp(b)}$. Since $b$ is odd, every prime
of $b$ is odd, so $b/D = 1$ iff $\vp(D) = \vp(b)$ for every $p \mid b$, iff
every prime of $b$ divides $a$, iff $\operatorname{rad}(b) \mid a$.
\end{proof}

\begin{remark}
The two conditions differ, and the difference is exactly the phenomenon that
$3x+9$ exhibits. Proposition~\ref{prop:onestep} requires $b \mid a$, which
fails for $(a,b) = (3,9)$; Theorem~\ref{thm:radical} requires only
$\operatorname{rad}(9) = 3 \mid 3$, which holds. The absorption into $9\Z$
takes two odd steps rather than one, but it happens, and $3x+9$ is a scaled
copy of $3x+1$. The same applies to $3x+27$, $3x+81$, $5x+25$ and $9x+27$.
So $b \mid a$ is sufficient but not necessary for equivalence; the radical is
the right invariant.
\end{remark}

% ============================================================ solved =======
\section{The solved families}
\label{sec:solved}

We now dispose of every universe that is not primitive. Throughout,
$a, b \geq 0$.

\begin{proposition}[$a = 0$: carykh's hub-and-spoke]\label{prop:const}
Let $a = 0$.
If $b = 0$, every orbit reaches the fixed point $0$.
If $b \geq 1$, the universe is \textsc{cyclic} and has exactly one cycle,
namely $b \to b/2 \to \cdots \to \oddpart(b) \to b$, of period
$\vtwo(b) + 1$; every orbit enters it.
\end{proposition}

\begin{proof}
For $b = 0$ every odd number maps to $0$ and every even number halves toward
$0$. For $b \geq 1$: every orbit reaches an odd number
(Lemma~\ref{lem:reachodd}), which maps to $b$; from $b$ the orbit halves to
$\oddpart(b)$ in $\vtwo(b)$ steps and returns to $b$. Any cycle must contain
an odd term, and every odd term maps to $b$, so every cycle contains $b$;
hence there is exactly one.
\end{proof}

\begin{proposition}[$b = 0$: carykh's shapes]\label{prop:shape}
Let $b = 0$ and $a \geq 1$, and write $a = 2^{t}u$ with $u$ odd. If $u = 1$
the universe is \textsc{cyclic}: every odd $n$ lies on the cycle
$n \to 2^{t}n \to \cdots \to n$ of period $t+1 = \log_{2}(a)+1$, and there are
infinitely many cycles. If $u \geq 3$ the universe is \textsc{divergent}.
\end{proposition}

\begin{proof}
For odd $n$, $T(n) = 2^{t}un$ and the next $t$ steps halve it to $un$. So on
odd numbers the induced map is $n \mapsto un$. If $u = 1$ this is the identity
and the displayed cycle is immediate; the period is $t+1$ because one odd step
is followed by exactly $t$ halvings. If $u \geq 3$ the odd terms form the
strictly increasing sequence $n, un, u^{2}n, \dots \to \infty$.
\end{proof}

Proposition~\ref{prop:shape} is carykh's ``all powers of two create clean-cut
shapes with $\log_{2}(n)+1$ sides'', with the accompanying observation that
$3x$ and $5x$ run away and that $6x$ runs away at the reduced rate of its odd
part.

\begin{theorem}[Parity obstruction: carykh's runaways]\label{thm:parity}
Let $a \geq 1$ and $b \geq 1$ with $a + b$ odd. Then every orbit from $\Np$
diverges; indeed every orbit is eventually a strictly increasing sequence of
odd numbers. The growth is linear if $a = 1$ and exponential if $a \geq 2$.
\end{theorem}

\begin{proof}
If $n$ is odd then $an+b$ is odd, because exactly one of $an$, $b$ is even.
So once an orbit reaches an odd number --- which it does, by
Lemma~\ref{lem:reachodd} --- it never again reaches an even one, and thereafter
$n \mapsto an+b$ with $a \geq 1$, $b \geq 1$, which strictly increases. The
growth rates are those of the recursion $n_{j+1} = an_{j}+b$.
\end{proof}

Theorem~\ref{thm:parity} explains the checkerboard of runaway cells that
dominates carykh's grid, and his ``finite state machine'' picture: the
$\text{odd} \to \text{even}$ transition simply does not exist. His four
subtypes (linear, exponential, and their halving-images) are
Theorem~\ref{thm:parity} together with Proposition~\ref{prop:halve}.

\begin{theorem}[$a=1$: carykh's black holes]\label{thm:descent}
Let $a = 1$ and let $b \geq 1$ be odd. Then:
\begin{enumerate}\itemsep2pt
\item[(i)] every $n > b$ strictly decreases within at most two steps;
\item[(ii)] every orbit from $\Np$ is eventually periodic;
\item[(iii)] every cycle has least element at most $b$, and is contained
      entirely in $[1, 2b]$;
\item[(iv)] $\{b, 2b\}$ is a cycle, and the bound $2b$ in (iii) is attained by
      it, so (iii) is sharp;
\item[(v)] over $\Z$: every orbit from $\Z_{\leq -1}$ enters $\Z_{\geq 0}$ in
      finitely many steps, so \emph{every} orbit of $\Tab$ on $\Z$ is
      eventually periodic.
\end{enumerate}
In particular the number of cycles is finite and effectively computable: it
suffices to run the seeds $1, \dots, b$.
\end{theorem}

\begin{proof}
(i) If $n$ is even then $T(n) = n/2 < n$. If $n$ is odd and $n > b$ then
$n + b$ is even and $T^{2}(n) = (n+b)/2 < n$.

(ii) By (i) the orbit strictly decreases until it enters the finite set
$[1,b]$, so it must repeat.

(iii) Let $C$ be a cycle and $m = \min C$. By (i), $m > b$ would make $m$
strictly decrease, contradicting minimality; so $m \leq b$.

For the upper bound, list the odd elements of $C$ in cyclic order as
$y_{1} \to y_{2} \to \cdots \to y_{k} \to y_{1}$, where $y_{i+1}$ is reached
from $y_{i}$ by one odd step followed by $d_{i} \geq 1$ halvings; the list is
non-empty by Lemma~\ref{lem:reachodd}. Then
\[
  y_{i+1} \;=\; \frac{y_{i} + b}{2^{d_{i}}} \;\leq\; \frac{y_{i}+b}{2}.
\]
First, some $y_{i} \leq b$: otherwise $y_{i} > b$ for every $i$, and then
$y_{i+1} \leq (y_{i}+b)/2 < y_{i}$ for every $i$, so the $y_{i}$ strictly
decrease around a finite cycle, which is impossible. Second, if $y_{i} \leq b$
then $y_{i+1} \leq (b+b)/2 = b$. By induction around the cycle, therefore,
$y_{i} \leq b$ for \emph{all} $i$.

Now every element of $C$ is either odd, hence $\leq b$, or even, hence of the
form $(y_{i} + b)/2^{j}$ for some $i$ and some $0 \leq j < d_{i}$, hence at
most $y_{i} + b \leq 2b$. Thus $C \subseteq [1, 2b]$.

(iv) $T(b) = b + b = 2b$ and $T(2b) = b$, and this cycle contains $2b$.

(v) Let $n \leq -1$. If $n$ is even, $|n|$ halves. If $n$ is odd, $T(n) =
n + b > n$, and $|T(n)| < |n|$ whenever $n < -b$. So $|n|$ is strictly
decreasing until the orbit reaches $[-b, -1]$, from which one odd step lands
in $[0, b-1] \subseteq \Z_{\geq 0}$. Then (ii) applies. \textup{(}This is
Proposition~\ref{prop:Zsplit}(iii) in concrete form: $b \leq a-1 = 0$ fails,
so $\Z_{\leq -1}$ is not invariant.\textup{)}
\end{proof}

\begin{remark}
Part (iii) is due to the present authors' earlier notes and is strictly
stronger than the bound on $\min C$ alone: it confines the whole cycle, which
is what makes the census in Table~\ref{tab:rowa1} \emph{exact} rather than a
lower bound. A search over all odd $b \leq 159$ finds
$\max(C)/b \leq 2$ with equality attained, confirming sharpness.
\end{remark}

\begin{remark}[The $a=1$ census is complete, and there are usually more than
two cycles]\label{rem:a1}
Theorem~\ref{thm:descent} makes this the one open-ended family whose cycle
count is not merely a lower bound: since every cycle meets $[1,b]$, testing the
seeds $1, \dots, b$ finds all of them. Table~\ref{tab:rowa1} is therefore
exact.

It also corrects a natural guess. The orbit of $1$ and the two-cycle
$\{b,2b\}$ are always present, and for small $b$ they are all there is --- so
one might expect exactly two cycles for every odd $b > 1$. That fails as soon
as $b = 7$, which has three: $\{1,8,4,2\}$, $\{3,10,5,12,6\}$ and
$\{7,14\}$. The count is $2$ for $b = 11,13,19,29,37,53,59,61$ but reaches
$13$ at $b = 63$. At $b=1$ the two named cycles coincide, so the count is
$1$.
\end{remark}

\begin{table}[!htbp]
\centering\footnotesize
\setlength{\tabcolsep}{4pt}
\begin{tabular}{@{}rrr@{\hskip 1.6em}rrr@{\hskip 1.6em}rrr@{\hskip 1.6em}rrr@{}}
\toprule
$b$ & all & prim & $b$ & all & prim & $b$ & all & prim & $b$ & all & prim \\
\midrule
% >>>BEGIN-DATA row_a1
1 & 1 & 1 & 21 & 6 & 2 & 41 & 3 & 2 & 61 & 2 & 1 \\
3 & 2 & 1 & 23 & 3 & 2 & 43 & 4 & 3 & 63 & 13 & 6 \\
5 & 2 & 1 & 25 & 3 & 1 & 45 & 8 & 2 & 65 & 7 & 4 \\
7 & 3 & 2 & 27 & 4 & 1 & 47 & 3 & 2 & 67 & 2 & 1 \\
9 & 3 & 1 & 29 & 2 & 1 & 49 & 5 & 2 & 69 & 6 & 2 \\
11 & 2 & 1 & 31 & 7 & 6 & 51 & 8 & 4 & 71 & 3 & 2 \\
13 & 2 & 1 & 33 & 5 & 2 & 53 & 2 & 1 & 73 & 9 & 8 \\
15 & 5 & 2 & 35 & 6 & 2 & 55 & 5 & 2 & 75 & 8 & 2 \\
17 & 3 & 2 & 37 & 2 & 1 & 57 & 5 & 2 & 77 & 6 & 2 \\
19 & 2 & 1 & 39 & 5 & 2 & 59 & 2 & 1 & 79 & 3 & 2 \\
% <<<END-DATA row_a1
\bottomrule
\end{tabular}
\caption{Cycles of $x+b$ for odd $b \leq 79$, with the primitive count of
Corollary~\ref{cor:cycdec}. Unlike every other census in this paper these
counts are \emph{exact}, by Remark~\ref{rem:a1}.}
\label{tab:rowa1}
\end{table}

Theorem~\ref{thm:descent} is the sharpening of carykh's remark that the
$a=1$ row has ``crushing decreases while barely giving any chance to grow''.
It is worth stating what it costs: the same two-step argument fails for
$a = 3$ at exactly the point where $(3n+b)/2 < n$ would require $n < -b$. The
$a = 1$ row is solved and the $a = 3$ row is not, and this is the whole of
the difference.

\subsection{The cycle through \texorpdfstring{$1$}{1}}
\label{sec:one}

A separate and very concrete question is when $1$ itself lies on a cycle. It
has a complete answer, and the answer explains a family of universes that
would otherwise look accidental.

\begin{theorem}[Power-of-two criterion]\label{thm:one}
Let $a \geq 1$ and $b \geq 1$. Then $1$ lies on a cycle of $\Tab$ if and only
if the forward orbit of $1$ contains a power of $2$.
\end{theorem}

\begin{proof}
($\Leftarrow$) If $T^{j}(1) = 2^{k}$ then $k$ further halvings give
$T^{j+k}(1) = 1$, so $1$ is periodic.

($\Rightarrow$) Suppose $1$ lies on a cycle $C$ and let $y \in C$ be its
predecessor. If $y$ is even then $y = 2$, a power of $2$, and $y$ lies on the
orbit. If $y$ is odd then $ay + b = 1$ with $y \geq 1$, forcing
$a + b \leq ay + b = 1$, impossible for $a, b \geq 1$.
\end{proof}

\begin{corollary}[$a + b$ a power of two]\label{cor:onestep}
If $a + b = 2^{k}$ then $1$ lies on a cycle of $\Tab$, of period exactly
$k+1$, namely $1 \to 2^{k} \to 2^{k-1} \to \cdots \to 2 \to 1$.
\end{corollary}

\begin{proof}
$T(1) = a + b = 2^{k}$, and the remaining $k$ steps are halvings.
\end{proof}

Corollary~\ref{cor:onestep} is the ``$a = 2^{k} - b$'' observation of the
authors' earlier notes. Theorem~\ref{thm:one} shows it is the one-step case of
a criterion that is both necessary and sufficient, and the gap is
substantial: among the $147$ pairs with $a$ odd, $a < 200$, $b$ odd, $b < 60$
for which $1$ lies on a cycle, exactly $93$ have $a+b$ a power of two, while
the other $54$ reach a power of two only after a longer excursion. For
instance in $3x+11$ the orbit of $1$ runs
$1 \to 14 \to 7 \to 32 \to 16 \to 8 \to 4 \to 2 \to 1$, meeting $2^{5}$ at the
fourth step although $3 + 11 = 14$ is not a power of two.

\begin{remark}
Corollary~\ref{cor:onestep} generalises verbatim to the $q$-ary map of
Section~\ref{sec:qary}: if $a + b = q^{k}$ then $1$ lies on a cycle of period
$k+1$. Theorem~\ref{thm:one} generalises too, with ``power of $2$'' replaced
by ``power of $q$''.
\end{remark}

We can now state the classification.

\begin{theorem}[Master classification]\label{thm:master}
Let $a \geq 0$ and $b \geq 1$, and let $(a',b') = \normalform(a,b)$. Then
$\Tab$ has the same fate as $T_{a',b'}$ \textup{(}Theorem~\ref{thm:redfate}\textup{)},
and exactly one of the following holds.
\begin{enumerate}\itemsep3pt
\item[\textup{(1)}] $a' = 0$. The fate is \textsc{cyclic}, with a single cycle
      of period $\vtwo(b')+1$ \textup{(}Proposition~\ref{prop:const}\textup{)}.
\item[\textup{(2)}] $a' \geq 1$ and $a'+b'$ is odd. The fate is
      \textsc{divergent} \textup{(}Theorem~\ref{thm:parity}\textup{)}.
\item[\textup{(3)}] $a' = 1$ and $b'$ is odd. The fate is \textsc{cyclic},
      with every cycle contained in $[1, 2b']$, so the cycle set is finite and
      computable \textup{(}Theorem~\ref{thm:descent}\textup{)}.
\item[\textup{(4)}] $a' \geq 3$ is odd, $b'$ is odd, and $\gcd(a',b')=1$. We
      call such a universe \emph{primitive}, and for these this paper proves
      no fate at all.
\end{enumerate}
The case $b = 0$ is separate and is settled by
Proposition~\ref{prop:shape}: the fate is \textsc{cyclic} if $a$ is a power of
two \textup{(}including $a=1$\textup{)}, \textsc{divergent} otherwise, and
$a = b = 0$ is terminal.
\end{theorem}

\begin{proof}
Exhaustiveness. By Theorem~\ref{thm:nf}, $\gcd(a',b')=1$, so $a'$ and $b'$ are
not both even. If $a' = 0$ we are in (1). Otherwise $a' \geq 1$; if
$a'+b'$ is odd we are in (2). If $a'+b'$ is even and they are not both even,
they are both odd; then $a' = 1$ gives (3) and $a' \geq 3$ gives (4).
Mutual exclusivity is clear. The cited results supply the fates.
\end{proof}

\begin{remark}[What case (4) does and does not say]
Cases (1)--(3) are proved outright and are statements about the dynamics.
Case (4) is not: ``primitive'' is a condition on $(a',b')$, and the
accompanying claim is about \emph{us}, namely that we prove nothing about
those maps. \textsc{open} is an epistemic label --- Definition~\ref{def:fate}
assigns it ``otherwise or when unknown'' --- so the honest reading of the
theorem is:

\begin{quote}
every universe reduces either to one of four families whose fate is proved
here, or to a primitive family for which this paper establishes no fate.
\end{quote}

It is emphatically not a claim that the primitive universes share some
dynamical property. What is expected of them is heuristic and is discussed in
Section~\ref{sec:open}; some of them, by Franco and Pomerance~\cite{fp}, are
in fact known to fail the Collatz-style statement.
\end{remark}

% =============================================================== table ======
\section{The table of universes}
\label{sec:table}

Table~\ref{tab:grid} classifies every universe with $0 \leq a, b \leq 12$
according to Theorem~\ref{thm:master}. The codes are
\textbf{C} (constant, $a=0$), \textbf{S} (shape, $b=0$ with $a$ a power of
two), \textbf{R} (runaway --- all orbits diverge), \textbf{H}
(halving-reducible), \textbf{A} (attraction-reducible), \textbf{D}
($a=1$ descent) and \textbf{?} (primitive and open).

\begin{table}[!htbp]
\centering
% >>>BEGIN-DATA grid
% generated by generate_data.py -- do not edit
\begin{tabular}{@{}rccccccccccccc@{}}
\toprule
$a\backslash b$ & 0 & 1 & 2 & 3 & 4 & 5 & 6 & 7 & 8 & 9 & 10 & 11 & 12 \\ \midrule
$0$ & \cell{gCon}{C} & \cell{gCon}{C} & \cell{gCon}{C} & \cell{gCon}{C} & \cell{gCon}{C} & \cell{gCon}{C} & \cell{gCon}{C} & \cell{gCon}{C} & \cell{gCon}{C} & \cell{gCon}{C} & \cell{gCon}{C} & \cell{gCon}{C} & \cell{gCon}{C} \\
$1$ & \cell{gShp}{S} & \cell{gDes}{D} & \cell{gRun}{R} & \cell{gDes}{D} & \cell{gRun}{R} & \cell{gDes}{D} & \cell{gRun}{R} & \cell{gDes}{D} & \cell{gRun}{R} & \cell{gDes}{D} & \cell{gRun}{R} & \cell{gDes}{D} & \cell{gRun}{R} \\
$2$ & \cell{gShp}{S} & \cell{gRun}{R} & \cell{gHal}{H} & \cell{gRun}{R} & \cell{gHal}{H} & \cell{gRun}{R} & \cell{gHal}{H} & \cell{gRun}{R} & \cell{gHal}{H} & \cell{gRun}{R} & \cell{gHal}{H} & \cell{gRun}{R} & \cell{gHal}{H} \\
$3$ & \cell{gRun}{R} & \cell{gOpen}{?} & \cell{gRun}{R} & \cell{gAtt}{A} & \cell{gRun}{R} & \cell{gOpen}{?} & \cell{gRun}{R} & \cell{gOpen}{?} & \cell{gRun}{R} & \cell{gAtt}{A} & \cell{gRun}{R} & \cell{gOpen}{?} & \cell{gRun}{R} \\
$4$ & \cell{gShp}{S} & \cell{gRun}{R} & \cell{gHal}{H} & \cell{gRun}{R} & \cell{gHal}{H} & \cell{gRun}{R} & \cell{gHal}{H} & \cell{gRun}{R} & \cell{gHal}{H} & \cell{gRun}{R} & \cell{gHal}{H} & \cell{gRun}{R} & \cell{gHal}{H} \\
$5$ & \cell{gRun}{R} & \cell{gOpen}{?} & \cell{gRun}{R} & \cell{gOpen}{?} & \cell{gRun}{R} & \cell{gAtt}{A} & \cell{gRun}{R} & \cell{gOpen}{?} & \cell{gRun}{R} & \cell{gOpen}{?} & \cell{gRun}{R} & \cell{gOpen}{?} & \cell{gRun}{R} \\
$6$ & \cell{gRun}{R} & \cell{gRun}{R} & \cell{gHal}{H} & \cell{gRun}{R} & \cell{gHal}{H} & \cell{gRun}{R} & \cell{gHal}{H} & \cell{gRun}{R} & \cell{gHal}{H} & \cell{gRun}{R} & \cell{gHal}{H} & \cell{gRun}{R} & \cell{gHal}{H} \\
$7$ & \cell{gRun}{R} & \cell{gOpen}{?} & \cell{gRun}{R} & \cell{gOpen}{?} & \cell{gRun}{R} & \cell{gOpen}{?} & \cell{gRun}{R} & \cell{gAtt}{A} & \cell{gRun}{R} & \cell{gOpen}{?} & \cell{gRun}{R} & \cell{gOpen}{?} & \cell{gRun}{R} \\
$8$ & \cell{gShp}{S} & \cell{gRun}{R} & \cell{gHal}{H} & \cell{gRun}{R} & \cell{gHal}{H} & \cell{gRun}{R} & \cell{gHal}{H} & \cell{gRun}{R} & \cell{gHal}{H} & \cell{gRun}{R} & \cell{gHal}{H} & \cell{gRun}{R} & \cell{gHal}{H} \\
$9$ & \cell{gRun}{R} & \cell{gOpen}{?} & \cell{gRun}{R} & \cell{gAtt}{A} & \cell{gRun}{R} & \cell{gOpen}{?} & \cell{gRun}{R} & \cell{gOpen}{?} & \cell{gRun}{R} & \cell{gAtt}{A} & \cell{gRun}{R} & \cell{gOpen}{?} & \cell{gRun}{R} \\
$10$ & \cell{gRun}{R} & \cell{gRun}{R} & \cell{gHal}{H} & \cell{gRun}{R} & \cell{gHal}{H} & \cell{gRun}{R} & \cell{gHal}{H} & \cell{gRun}{R} & \cell{gHal}{H} & \cell{gRun}{R} & \cell{gHal}{H} & \cell{gRun}{R} & \cell{gHal}{H} \\
$11$ & \cell{gRun}{R} & \cell{gOpen}{?} & \cell{gRun}{R} & \cell{gOpen}{?} & \cell{gRun}{R} & \cell{gOpen}{?} & \cell{gRun}{R} & \cell{gOpen}{?} & \cell{gRun}{R} & \cell{gOpen}{?} & \cell{gRun}{R} & \cell{gAtt}{A} & \cell{gRun}{R} \\
$12$ & \cell{gRun}{R} & \cell{gRun}{R} & \cell{gHal}{H} & \cell{gRun}{R} & \cell{gHal}{H} & \cell{gRun}{R} & \cell{gHal}{H} & \cell{gRun}{R} & \cell{gHal}{H} & \cell{gRun}{R} & \cell{gHal}{H} & \cell{gRun}{R} & \cell{gHal}{H} \\
\bottomrule
\end{tabular}
% <<<END-DATA grid
\caption{The multiverse for $0 \leq a,b \leq 12$. Only the cells marked
\textbf{?} are open, and only those; every other cell is settled by
Theorem~\ref{thm:master}. The cells marked \textbf{H} and \textbf{A} are open
or solved according to the fate of their normal form.}
\label{tab:grid}
\end{table}

The visible regularities all have proofs above. The row $a=0$ and the column
$b=0$ are Propositions~\ref{prop:const} and~\ref{prop:shape}. The
checkerboard of \textbf{R} is Theorem~\ref{thm:parity}. The quarter-lattice of
\textbf{H} at even $(a,b)$ is Proposition~\ref{prop:halve}. The \textbf{A}
cells lie where an odd prime divides both coordinates --- so the diagonal
$a = b$ with $a$ odd, and more generally the sublattices $p \mid a$, $p \mid b$.
What is left is the odd--odd sublattice with $\gcd = 1$.

\subsection{carykh's nine}

carykh singles out nine cells of his $6 \times 6$ grid as having ``no obvious
indicators'' --- these are exactly the cells with $a, b \in \{1,3,5\}$, i.e.\
the odd--odd sublattice, which is the correct set to single out.
Theorem~\ref{thm:master} resolves five of the nine.

\begin{theorem}\label{thm:nine}
Of carykh's nine universes, three are solved, two reduce, and four are open
and pairwise affinely inequivalent:
\begin{center}\small
\begin{tabular}{@{}lll@{}}
\toprule
universe & family & status \\
\midrule
% >>>BEGIN-DATA carykh_nine
$1x+1$ & descent & \solved{every orbit cycles} \\
$1x+3$ & descent & \solved{every orbit cycles} \\
$1x+5$ & descent & \solved{every orbit cycles} \\
$3x+1$ & primitive & \openv{open} \\
$3x+3$ & attracting-reducible & \reduc{equivalent to $3x+1$} \\
$3x+5$ & primitive & \openv{open} \\
$5x+1$ & primitive & \openv{open} \\
$5x+3$ & primitive & \openv{open} \\
$5x+5$ & attracting-reducible & \reduc{equivalent to $5x+1$} \\
% <<<END-DATA carykh_nine
\bottomrule
\end{tabular}
\end{center}
The four open universes are $3x+1$, $3x+5$, $5x+1$ and $5x+3$.
\end{theorem}

\begin{proof}
$1x+1$, $1x+3$, $1x+5$ are Theorem~\ref{thm:descent}. $3x+3$ and $5x+5$ have
$3 \mid \gcd(3,3)$ and $5 \mid \gcd(5,5)$, so rule \textbf{(A)} applies and
Theorem~\ref{thm:redfate} identifies their fates with those of $3x+1$ and
$5x+1$. The remaining four are primitive by inspection of $\gcd$, hence case
(4) of Theorem~\ref{thm:master}. For pairwise inequivalence, suppose an affine
$\varphi$ intertwines two of them. By Theorem~\ref{thm:rigid} it has
$a' = a$ and $b' = \lambda b$ with $\lambda$ an odd integer, and if $\varphi$
is a bijection of $\Z$ then $\lambda = \pm 1$. Different multipliers rule out
every pair with $a = 3$ against $a = 5$; and within a row, $b' = \pm b$ fails
for $(1,5)$ and for $(1,3)$. So no two of the four are affinely conjugate.
\end{proof}

\begin{remark}
The four are inequivalent \emph{as members of this family under affine change
of variable}, which is what Theorem~\ref{thm:rigid} controls. That is weaker
than being non-isomorphic as abstract dynamical systems, and we do not claim
the stronger statement: $n \mapsto 5n$ does embed $T_{3,1}$ into $T_{3,5}$
(Proposition~\ref{prop:scale}), just not onto it.
\end{remark}

\begin{remark}
$3x+3$ deserves a word, because carykh treats it as a distinct universe with
its own character, noting that its graph has many extra branches that
$3x+1$ lacks. Both observations are right and compatible:
Theorem~\ref{thm:redfate} says the two have identical fates and
period-preserving matching cycles, while Corollary~\ref{cor:conn} says the
copy of $3x+1$ sits inside $3x+3$ as an attractor with a transient forest
--- the branches --- hanging off it. Dynamically they are the same problem;
as labelled graphs they are not isomorphic.
\end{remark}

\subsection{How much of the multiverse is open}

\begin{theorem}[Densities]\label{thm:density}
As $N \to \infty$,
\begin{align*}
  \frac{\#\{(a,b) \in [1,N]^{2} : \Tab \text{ open}\}}{N^{2}}
    &\;\longrightarrow\; \frac{1}{3},
  \\[1mm]
  \frac{\#\{(a,b) \in [1,N]^{2} : \Tab \text{ open and primitive}\}}{N^{2}}
    &\;\longrightarrow\; \frac{2}{\pi^{2}} .
\end{align*}
\end{theorem}

\begin{proof}
By Theorem~\ref{thm:master}, $\Tab$ is open iff its normal form is primitive,
i.e.\ iff $a' \geq 3$ is odd. Rule \textbf{(A)} changes neither $a$ nor the
parity of $b$, and rule \textbf{(H)} is applied exactly
$s = \min(\vtwo(a),\vtwo(b))$ times, so $a' = a/2^{s}$ and
$b' = b/(2^{s}D)$ with $D$ odd. Both $a'$ and $b'$ are odd exactly when
$\vtwo(a) = \vtwo(b)$. Hence
\[
  \Tab \text{ open}
  \iff \vtwo(a) = \vtwo(b) \ \text{ and } \ a/2^{\vtwo(a)} \geq 3 .
\]
The second condition fails only when $a$ is a power of two, a set of density
zero. The natural density of $\{a : \vtwo(a) = s\}$ is $2^{-s-1}$, and the two
coordinates are independent, so
$\sum_{s \geq 0} 4^{-s-1} = \tfrac14 \cdot \tfrac{4}{3} = \tfrac13$.

For the second limit: a universe is open and primitive iff $a \geq 3$ is odd,
$b$ is odd and $\gcd(a,b)=1$. Conditioning on both being odd (density $1/4$),
the density of coprimality among odd pairs is
\[
  \prod_{p \text{ odd prime}} \Bigl(1 - \frac{1}{p^{2}}\Bigr)
  \;=\; \frac{\zeta(2)^{-1}}{1 - 2^{-2}}
  \;=\; \frac{6/\pi^{2}}{3/4}
  \;=\; \frac{8}{\pi^{2}},
\]
and removing the density-zero set of powers of two leaves
$\tfrac14 \cdot \tfrac{8}{\pi^{2}} = 2/\pi^{2}$.
\end{proof}

\begin{table}[!htbp]
\centering\small
\begin{tabular}{@{}rrrrr@{}}
\toprule
$N$ & open & open$/N^{2}$ & open primitive & primitive$/N^{2}$ \\
\midrule
% >>>BEGIN-DATA density
$10$ & $26$ & $0.260000$ & $14$ & $0.140000$ \\
$50$ & $794$ & $0.317600$ & $486$ & $0.194400$ \\
$100$ & $3244$ & $0.324400$ & $1957$ & $0.195700$ \\
$250$ & $20646$ & $0.330336$ & $12508$ & $0.200128$ \\
$500$ & $82896$ & $0.331584$ & $50451$ & $0.201804$ \\
$1000$ & $332396$ & $0.332396$ & $202161$ & $0.202161$ \\
$2000$ & $1331396$ & $0.332849$ & $809865$ & $0.202466$ \\
% <<<END-DATA density
\bottomrule
\end{tabular}
\caption{Convergence to the constants of Theorem~\ref{thm:density}:
$1/3 = 0.3333\ldots$ and $2/\pi^{2} = 0.202642\ldots$}
\label{tab:density}
\end{table}

Table~\ref{tab:density} shows the convergence.
So two thirds of the multiverse is settled outright, and of the remaining
third, about $61\%$ ($= 3 \cdot 2/\pi^{2}$) is irreducibly new while the rest
reduces to a smaller open universe. carykh's intuition that ``Collatz picked
the only rule set that wasn't obvious in some way'' is a little too strong:
the open region has positive density. What is true is that he picked the
smallest interesting one --- $3x+1$ is the least primitive universe in the
lexicographic order on $(a,b)$ with $a \geq 3$.

\subsection{The knowledge ladder}
\label{sec:ladder}

Theorem~\ref{thm:master} sorts universes by \emph{how} they are settled.
For deciding what is worth studying, a coarser grading is more useful, and it
separates two things the family classification conflates. Consider $2x+2$ and
$6x+2$: both are halving-reducible, so the classification of
Section~\ref{sec:table} gives them the same label. But $2x+2$ reduces to
$1x+1$, which is solved, while $6x+2$ reduces to $3x+1$ and is therefore
exactly as open as the Collatz conjecture. Any presentation that colours the
two alike is misleading.

\begin{definition}[Knowledge tiers]\label{def:tiers}
Write $(a',b') = \normalform(a,b)$. The universe $\Tab$ is
\begin{itemize}\itemsep2pt
\item \textsc{closed} if $T_{a',b'}$ has a proved fate and its complete cycle
      set is given in closed form --- that is, $a' = 0$, or $b' = 0$, or
      $a'+b'$ is odd;
\item \textsc{decided} if $a' = 1$: the fate is proved and the cycle set is
      finite and computable for each $b'$, but no formula in $b'$ is known;
\item \textsc{reducible} if $T_{a',b'}$ is open and $(a,b) \neq (a',b')$;
\item \textsc{open} if $T_{a',b'}$ is open and $(a,b) = (a',b')$.
\end{itemize}
\end{definition}

\begin{proposition}\label{prop:tiers}
The four tiers partition the family. \textsc{decided} is exactly the set of
$(a,b)$ with $\vtwo(a) = \vtwo(b)$ and $a$ a power of two; \textsc{open} is
exactly the set of primitive universes. Their natural densities are
\[
  \textsc{closed} \to \tfrac{2}{3}, \qquad
  \textsc{decided} \to 0, \qquad
  \textsc{reducible} \to \tfrac{1}{3} - \tfrac{2}{\pi^{2}}, \qquad
  \textsc{open} \to \tfrac{2}{\pi^{2}} .
\]
\end{proposition}

\begin{proof}
The tiers are mutually exclusive and exhaustive by construction, since
$T_{a',b'}$ is either open or proved, and in the proved case
Theorem~\ref{thm:master} leaves exactly the four normal-form shapes listed.
For \textsc{decided}: $a' = 1$ means $a/2^{s} = 1$ with
$s = \min(\vtwo(a),\vtwo(b))$, i.e.\ $a = 2^{s}$, and $b'$ odd forces
$\vtwo(b) = s$. The powers of two have density $0$, so the tier does.
\textsc{open} is the primitive set, of density $2/\pi^{2}$ by
Theorem~\ref{thm:density}; the open universes have density $1/3$, so
\textsc{reducible} takes the difference; and the proved universes make up the
remaining $2/3$, all but a density-zero part of which is \textsc{closed}.
\end{proof}

\begin{table}[!htbp]
\centering\small
\begin{tabular}{@{}rrrrr@{}}
\toprule
$N$ & \textsc{closed} & \textsc{decided} & \textsc{reducible} & \textsc{open} \\
\midrule
% >>>BEGIN-DATA tiers
$10$ & $0.640000$ & $0.100000$ & $0.120000$ & $0.140000$ \\
$50$ & $0.662400$ & $0.020000$ & $0.123200$ & $0.194400$ \\
$100$ & $0.665600$ & $0.010000$ & $0.128700$ & $0.195700$ \\
$250$ & $0.665664$ & $0.004000$ & $0.130208$ & $0.200128$ \\
$500$ & $0.666416$ & $0.002000$ & $0.129780$ & $0.201804$ \\
$1000$ & $0.666604$ & $0.001000$ & $0.130235$ & $0.202161$ \\
limit & $0.666667$ & $0.000000$ & $0.130691$ & $0.202642$ \\
% <<<END-DATA tiers
\bottomrule
\end{tabular}
\caption{Proportion of $(a,b) \in [1,N]^{2}$ in each tier of
Definition~\ref{def:tiers}, against the limits of
Proposition~\ref{prop:tiers}.}
\label{tab:tiers}
\end{table}

Table~\ref{tab:tiers} confirms the four limits numerically.

\begin{remark}
The \textsc{decided} tier is not empty, which is worth saying because it is
easy to assume that a proved fate comes with a known cycle structure. The
$a=1$ row is a counterexample: Theorem~\ref{thm:descent} settles the fate in
two lines, and Remark~\ref{rem:a1} shows the cycle counts
$1, 2, 2, 3, 3, 2, 2, 5, 3, 2, 6, \dots$ follow no pattern we can see. It is
also the only such row: every other proved family has its cycles in closed
form.
\end{remark}

% ================================================================ cycles ====
\section{Cycles}
\label{sec:cycles}

\subsection{The cycle equation}

\begin{theorem}[Cycle equation]\label{thm:cyceq}
Let $C$ be a cycle of $\Tab$ containing $k \geq 1$ odd terms
$n_{1}, \dots, n_{k}$ in cyclic order, where $n_{i+1}$ is reached from
$n_{i}$ by one odd step followed by $d_{i} \geq 1$ halvings, and let
$m = d_{1}+\cdots+d_{k}$ be the total number of halvings. Then
\[
  n_{1}\bigl(2^{m} - a^{k}\bigr)
  \;=\; b \sum_{i=1}^{k} a^{\,k-i}\, 2^{\,d_{1}+\cdots+d_{i-1}} .
\]
In particular $2^{m} > a^{k}$ for any cycle in $\Np$ with $b \geq 1$, so
$m/k > \log_{2} a$.
\end{theorem}

\begin{proof}
By construction $2^{d_{i}} n_{i+1} = a n_{i} + b$ for each $i$, indices modulo
$k$. Multiply the $i$-th relation by $a^{k-i} 2^{d_{1}+\cdots+d_{i-1}}$ and
sum over $i$; all terms $n_{2},\dots,n_{k}$ telescope away, leaving
$2^{m}n_{1} = a^{k}n_{1} + b\sum_{i} a^{k-i}2^{d_{1}+\cdots+d_{i-1}}$.
The right-hand sum is positive when $b \geq 1$, and $n_1 \geq 1$, so
$2^{m} - a^{k} > 0$.
\end{proof}

The inequality $m/k > \log_{2}a$ is the exact form of carykh's remark that his
two period-$8$ cycles in $3x+5$ have ``three rises and five falls, which
corresponds to $3^{3} = 27$ being close to $2^{5} = 32$'', and that the
period-$44$ cycles have $k=17$, $m=27$ with
$3^{17} = 129{,}140{,}163$ close to $2^{27} = 134{,}217{,}728$. Our census
confirms both $(k,m)$ pairs exactly. The mechanism is now visible: small
cycle elements require $2^{m} - a^{k}$ to be large relative to the right-hand
side, but $m/k$ must exceed $\log_{2}a$ only slightly for the sum to stay
small. This suggests --- heuristically; the theorem gives only
$m/k > \log_{2}a$, and nothing more --- that cycles should cluster where
$m/k$ is a good rational approximation to $\log_{2}a$, that is near the
convergents of its continued fraction. For $a=3$, $\log_{2}3 = 1.5849\ldots$ and $27/17 = 1.5882\ldots$

\subsection{Strata and the divisor decomposition}

\begin{theorem}[Stratification of the phase space]\label{thm:strata}
Let $\gcd(a,b)=1$ with $a, b \geq 1$ odd. For $n \geq 1$ put
\[
  \pi(n) \;=\; \prod_{p \mid b} p^{\min(\vp(n),\,\vp(b))} ,
\]
a divisor of $b$. Then $\pi$ is constant along orbits, and
\[
  \Np \;=\; \bigsqcup_{d \mid b} \Sigma_{d},
  \qquad \Sigma_{d} = \pi^{-1}(d),
\]
is a partition into $\tau(b)$ non-empty $\Tab$-invariant sets. Moreover
$n \mapsto n/d$ is an isomorphism from $\Tab|_{\Sigma_{d}}$ onto
$T_{a,b/d}|_{\Sigma_{1}}$, where $\Sigma_{1}$ is the corresponding
\emph{primitive stratum} of $T_{a,b/d}$.
\end{theorem}

\begin{proof}
Since $\gcd(a,b)=1$, every prime $p \mid b$ is odd and satisfies $p \nmid a$,
so Theorem~\ref{thm:strat}(ii) applies to each: $\min(\vp(\cdot),\vp(b))$ is
conserved. Hence so is $\pi$, and the fibres partition $\Np$ into invariant
sets. The value set is all of $\{d : d \mid b\}$ because $\pi(d) = d$ for each
divisor $d$ of $b$: indeed for $p \mid b$ we have
$\min(\vp(d),\vp(b)) = \vp(d)$.

Fix $d \mid b$ and $n \in \Sigma_{d}$. For each $p \mid b$, either
$\vp(d) < \vp(b)$, in which case $\vp(n) = \vp(d)$ so $\vp(n/d) = 0$; or
$\vp(d) = \vp(b)$, in which case $p \nmid b/d$. Either way
$\pi_{b/d}(n/d) = 1$, so $n/d \in \Sigma_{1}$ for $T_{a,b/d}$; and the map is
onto with inverse $m \mapsto dm$, which lands in $\Sigma_{d}$ by the same
computation. Since $d$ is odd, Proposition~\ref{prop:scale} makes it a
conjugacy.
\end{proof}

\begin{corollary}[Divisor decomposition of the cycle set]\label{cor:cycdec}
Under the hypotheses of Theorem~\ref{thm:strata}, writing $\PCyc$ for the set
of cycles lying in the primitive stratum,
\[
  \Cyc(\Tab) \;=\; \bigsqcup_{d \mid b} d \cdot \PCyc\bigl(T_{a,\,b/d}\bigr),
\]
and the bijections preserve periods and the pair $(k,m)$.
\end{corollary}

\begin{proof}
Each cycle lies in a single stratum, and Theorem~\ref{thm:strata} identifies
the $d$-th stratum with the primitive stratum of $T_{a,b/d}$ scaled by $d$.
Conjugacies preserve periods, and they preserve parities --- $d$ is odd ---
hence $k$ and $m$.
\end{proof}

Corollary~\ref{cor:cycdec} is the precise form of Artem's statement that ``a
composite universe contains a parallel universe of every single one of its
divisors''. Two things are added. First, the copies do not merely exist: they
\emph{exhaust} the cycle set, so the count is an exact sum, not a lower bound.
Second, the same holds for the whole phase space, not only the cycles
(Theorem~\ref{thm:strata}).

\begin{table}[!htbp]
\centering\small
\begin{tabular}{@{}llL@{}}
\toprule
universe & cycles & decomposition \\
\midrule
% >>>BEGIN-DATA decomposition
$3x+5$ & $6$ & $1{\cdot}\mathrm{P}(3x+5)^{5} \;\sqcup\; 5{\cdot}\mathrm{P}(3x+1)^{1}$ \\
$3x+7$ & $2$ & $1{\cdot}\mathrm{P}(3x+7)^{1} \;\sqcup\; 7{\cdot}\mathrm{P}(3x+1)^{1}$ \\
$3x+25$ & $8$ & $1{\cdot}\mathrm{P}(3x+25)^{2} \;\sqcup\; 5{\cdot}\mathrm{P}(3x+5)^{5} \;\sqcup\; 25{\cdot}\mathrm{P}(3x+1)^{1}$ \\
$3x+35$ & $9$ & $1{\cdot}\mathrm{P}(3x+35)^{2} \;\sqcup\; 5{\cdot}\mathrm{P}(3x+7)^{1} \;\sqcup\; 7{\cdot}\mathrm{P}(3x+5)^{5} \;\sqcup\; 35{\cdot}\mathrm{P}(3x+1)^{1}$ \\
$3x+49$ & $3$ & $1{\cdot}\mathrm{P}(3x+49)^{1} \;\sqcup\; 7{\cdot}\mathrm{P}(3x+7)^{1} \;\sqcup\; 49{\cdot}\mathrm{P}(3x+1)^{1}$ \\
$3x+65$ & $16$ & $1{\cdot}\mathrm{P}(3x+65)^{1} \;\sqcup\; 5{\cdot}\mathrm{P}(3x+13)^{9} \;\sqcup\; 13{\cdot}\mathrm{P}(3x+5)^{5} \;\sqcup\; 65{\cdot}\mathrm{P}(3x+1)^{1}$ \\
% <<<END-DATA decomposition
\bottomrule
\end{tabular}
\caption{Corollary~\ref{cor:cycdec} in action;
$\mathrm{P}(ax+b)^{j}$ denotes the $j$ primitive cycles of $ax+b$.
Compare Artem's counts of $9$, $8$ and $3$ for $b = 35, 25, 49$.}
\label{tab:decomp}
\end{table}

Table~\ref{tab:decomp} works the decomposition through for several offsets.

\begin{example}[$b = 35$]
$\Cyc(T_{3,35})$ has nine elements: two primitive, one equal to
$5 \cdot \PCyc(T_{3,7})$, five equal to $7 \cdot \PCyc(T_{3,5})$, and one
equal to $35 \cdot \PCyc(T_{3,1})$. Artem obtains the same nine by hand.
\end{example}

\begin{example}[$b = 65$]
Artem observes that multiplying two offsets with unusually many cycles gives
one with even more. Corollary~\ref{cor:cycdec} makes this exact: $5$ and $13$
contribute $5$ and $9$ primitive cycles respectively, and
$65 = 5 \cdot 13$ has $1 + 5 + 9 + 1 = 16$.
\end{example}

\subsection{Census}

Table~\ref{tab:rowa3} gives cycle counts for the row $a = 3$. The search
window is seeds $1 \leq n \leq 20{,}000$ with a step budget of $200{,}000$;
this window is \emph{empirically stable}: every count is identical at
$20{,}000$ and at $100{,}000$ seeds. Stability across the windows we tested is
evidence, not proof, that the counts are complete. Counts remain lower bounds in
principle --- a cycle whose basin of attraction misses the seed window is
invisible to any such search --- so they should be read as ``at least''.

\begin{table}[!htbp]
\centering\footnotesize
\setlength{\tabcolsep}{4pt}
\begin{tabular}{@{}rrr@{\hskip 1.6em}rrr@{\hskip 1.6em}rrr@{\hskip 1.6em}rrr@{}}
\toprule
$b$ & all & prim & $b$ & all & prim & $b$ & all & prim & $b$ & all & prim \\
\midrule
% >>>BEGIN-DATA row_a3
1 & 1 & 1 & 27 & 1 & 1 & 53 & 2 & 1 & 79 & 5 & 4 \\
3 & 1 & 1 & 29 & 5 & 4 & 55 & 11 & 3 & 81 & 1 & 1 \\
5 & 6 & 5 & 31 & 2 & 1 & 57 & 2 & 1 & 83 & 4 & 3 \\
7 & 2 & 1 & 33 & 3 & 2 & 59 & 8 & 7 & 85 & 9 & 1 \\
9 & 1 & 1 & 35 & 9 & 2 & 61 & 3 & 2 & 87 & 5 & 4 \\
11 & 3 & 2 & 37 & 4 & 3 & 63 & 2 & 1 & 89 & 2 & 1 \\
13 & 10 & 9 & 39 & 10 & 9 & 65 & 16 & 1 & 91 & 14 & 3 \\
15 & 6 & 5 & 41 & 2 & 1 & 67 & 2 & 1 & 93 & 2 & 1 \\
17 & 3 & 2 & 43 & 2 & 1 & 69 & 4 & 3 & 95 & 10 & 3 \\
19 & 2 & 1 & 45 & 6 & 5 & 71 & 8 & 7 & 97 & 3 & 2 \\
21 & 2 & 1 & 47 & 8 & 7 & 73 & 4 & 3 & 99 & 3 & 2 \\
23 & 4 & 3 & 49 & 3 & 1 & 75 & 8 & 2 &  &  &  \\
25 & 8 & 2 & 51 & 3 & 2 & 77 & 5 & 1 &  &  & \\
% <<<END-DATA row_a3
\bottomrule
\end{tabular}
\caption{Cycles of $3x+b$ for odd $b \leq 99$: total count, and the count of
\emph{primitive} cycles (those in the stratum $\Sigma_{1}$). By
Corollary~\ref{cor:cycdec} the total is the sum of the primitive counts of
$3x+(b/d)$ over the divisors $d \mid b$.}
\label{tab:rowa3}
\end{table}

Three of Artem's observations about this row are confirmed, and one is
explained.

\begin{corollary}\label{cor:pow3}
$3x+b$ has exactly one cycle whenever $b = 3^{j}$, and more generally
$\#\Cyc(3x+3^{j}b_{0}) = \#\Cyc(3x+b_{0})$ for every $j \geq 0$ and every
$b_{0}$ coprime to $3$.
\end{corollary}

\begin{proof}
$3 \mid a$, so rule \textbf{(A)} strips $3^{v_{3}(b)}$ entirely and
Theorem~\ref{thm:redfate} gives a period-preserving bijection of cycle sets.
For $b = 3^{j}$ the normal form is $3x+1$, whose only known cycle is
$\{1,4,2\}$.
\end{proof}

This is Artem's ``only the universes that are a power of three away from the
original have one loop, and the same for any prime universe''. Note the
statement is conditional on $3x+1$ having only one cycle, which is not known;
what \emph{is} unconditional is the equality of counts.

Artem also observes that no $b < 100$ yields exactly seven cycles. Our
wider search confirms this: the multiset of counts for odd $b \leq 99$
contains $1,2,3,4,5,6,8,9,10,11,14$ and $16$, but not $7$. We have no
explanation, and we note it is sensitive to the search window --- at
$3{,}000$ seeds $b=71$ appears to have seven, and only widening the window to
$20{,}000$ reveals its eighth cycle.

Table~\ref{tab:colb1} gives the complementary slice: the column $b=1$, where
$\gcd(a,1)=1$ forces every universe to be primitive and irreducible.

\begin{table}[!htbp]
\centering\small
\begin{tabular}{@{}rrrl@{}}
\toprule
$a$ & cycles & drift $a/4$ & periods \\
\midrule
% >>>BEGIN-DATA col_b1
$1$ & $1$ & $0.25$ & 2 \\
$3$ & $1$ & $0.75$ & 3 \\
$5$ & $3$ & $1.25$ & 7, 10, 10 \\
$7$ & $1$ & $1.75$ & 4 \\
$9$ & $0$ & $2.25$ & --- \\
$11$ & $0$ & $2.75$ & --- \\
$13$ & $0$ & $3.25$ & --- \\
$15$ & $1$ & $3.75$ & 5 \\
$17$ & $0$ & $4.25$ & --- \\
$19$ & $0$ & $4.75$ & --- \\
$21$ & $0$ & $5.25$ & --- \\
$23$ & $0$ & $5.75$ & --- \\
$25$ & $0$ & $6.25$ & --- \\
$27$ & $0$ & $6.75$ & --- \\
$29$ & $0$ & $7.25$ & --- \\
$31$ & $1$ & $7.75$ & 6 \\
$33$ & $0$ & $8.25$ & --- \\
$35$ & $0$ & $8.75$ & --- \\
$37$ & $0$ & $9.25$ & --- \\
$39$ & $0$ & $9.75$ & --- \\
% <<<END-DATA col_b1
\bottomrule
\end{tabular}
\caption{Cycles of $ax+1$ for odd $a \leq 39$. Every entry with $a \geq 3$ is
open. As the drift grows the cycles thin out and the seed window finds fewer
of them.}
\label{tab:colb1}
\end{table}

% =============================================================== open =======
\section{The open region}
\label{sec:open}

Nothing in Section~\ref{sec:solved} distinguishes $3x+1$ from $5x+1$; both are
case (4). The distinction is heuristic, and it is the one carykh draws when he
calls the $a=3$ row ``Goldilocks'' and the $a=5$ row ``Icarus''.

\begin{definition}[Drift]
For a primitive universe define
\[
  \delta_{\mathrm{step}} = \frac{\sqrt{a}}{2},
  \qquad
  \delta_{\mathrm{odd}} = \frac{a}{4} .
\]
\end{definition}

The first is the geometric mean of the two branch multipliers $a$ and $1/2$
weighted equally --- carykh's ``gravity''. The second is the expected ratio
between consecutive odd terms: an odd $n$ maps to $an+b \approx an$, which is
divided by $2^{v}$ where $v = \vtwo(an+b)$ has expectation $2$ under the
standard equidistribution heuristic. Both cross $1$ at $a = 4$, so:

\begin{itemize}\itemsep2pt
\item $a = 1$: $\delta_{\mathrm{odd}} = 1/4$. Strong contraction --- and here
      the heuristic is a theorem (Theorem~\ref{thm:descent}).
\item $a = 3$: $\delta_{\mathrm{odd}} = 3/4$. Weak contraction. Almost all
      orbits are expected to be eventually periodic, and the number of cycles
      is expected to be finite but is not known to be.
\item $a \geq 5$: $\delta_{\mathrm{odd}} \geq 5/4$. Expansion. Almost all
      orbits are expected to diverge, with finitely many cycles surviving as
      exceptional sets.
\end{itemize}

These are heuristics with no proof in either direction, and the reader should
treat the word ``expected'' strictly. What is worth recording is that the
threshold $a < 4$ admits only $a \in \{1,3\}$ among odd multipliers, and
$a=1$ is solved. If one grants the drift picture --- and it is only a
picture --- then the difficulty of the family is concentrated in the single
row $a = 3$, with the expanding rows $a \geq 5$ hard in the opposite
direction. We stress that this is a reading of a heuristic, not a corollary
of anything proved here.

\subsection{Empirical growth statistics}

Table~\ref{tab:growth} reproduces carykh's ``average max growth'' heat map:
the geometric mean over seeds $n \leq 5000$ of $\max(\mathrm{orbit}(n))/n$.

\begin{table}[!htbp]
\centering\small
\setlength{\tabcolsep}{4pt}
\begin{tabular}{@{}rccccccccc@{}}
\toprule
$a\backslash b$ & 0 & 1 & 2 & 3 & 4 & 5 & 6 & 7 & 8 \\
\midrule
% >>>BEGIN-DATA growth
$0$ & 1.00 & 1.00 & 1.00 & 1.00 & 1.00 & 1.00 & 1.00 & 1.00 & 1.00 \\
$1$ & 1.00 & 1.00 & \rn & 1.00 & \rn & 1.00 & \rn & 1.01 & \rn \\
$2$ & 1.41 & \rn & 1.42 & \rn & \rn & \rn & 1.42 & \rn & \rn \\
$3$ & \rn & 4.55 & \rn & 5.29 & \rn & 3.95 & \rn & 3.89 & \rn \\
$4$ & 2.38 & \rn & \rn & \rn & 2.38 & \rn & \rn & \rn & \rn \\
$5$ & \rn & 18.33 & \rn & 7.64 & \rn & 12.97 & \rn & 16.12 & \rn \\
$6$ & \rn & \rn & 8.61 & \rn & \rn & \rn & 9.89 & \rn & \rn \\
$7$ & \rn & 2.38 & \rn & 2.22 & \rn & 7.66 & \rn & 3.07 & \rn \\
$8$ & 4.36 & \rn & \rn & \rn & \rn & \rn & \rn & \rn & 4.38 \\
% <<<END-DATA growth
\bottomrule
\end{tabular}
\caption{Geometric mean of $\max(\mathrm{orbit}(n))/n$ over $n \leq 5000$;
\rn{} marks universes where every seed exceeded the value or step budget.}
\label{tab:growth}
\end{table}

The \rn{} pattern is the checkerboard of Theorem~\ref{thm:parity} together
with its halving images. The column $b=0$ is exactly computable, which pins
down the corner of the table carykh reasons about.

\begin{proposition}\label{prop:peak}
Let $b = 0$ and $a = 2^{t}$ with $t \geq 0$. For $n \geq 1$ the relative peak
is $\max(\mathrm{orbit}(n))/n = 2^{\max(\vtwo(n),\,t) - \vtwo(n)}$, and its
geometric mean over $n$ --- weighting $\vtwo(n) = j$ by its natural density
$2^{-j-1}$ --- is exactly
\[
  2^{\,t - 1 + 2^{-t}} .
\]
\end{proposition}

\begin{proof}
Write $n = 2^{j}m$ with $m$ odd. By Proposition~\ref{prop:shape} the orbit
descends to $m$ and then cycles through $2^{t}m, 2^{t-1}m, \dots, m$, so its
maximum is $\max(2^{j}, 2^{t})\,m$ and the relative peak is
$2^{\max(j,t)-j}$. The natural density of $\{n : \vtwo(n) = j\}$ is
$2^{-j-1}$, so the mean of $\log_{2}$ of the relative peak is
\[
  \sum_{j \geq 0} 2^{-j-1}\bigl(\max(j,t) - j\bigr)
  = \sum_{j=0}^{t-1} (t-j)\,2^{-j-1}
  = t - 1 + 2^{-t},
\]
the last equality by the standard evaluation of
$\sum_{j \geq 0} j x^{j}$ at $x = 1/2$.
\end{proof}

For $t = 1, 2, 3$ this gives $1.414214$, $2.378414$, $4.362031$; the measured
values over $n \leq 200{,}000$ agree to six decimal places. carykh reports
$\sqrt{2}$ for the $a = 2$ row, which is the case $t = 1$; his argument ---
that an odd $n$ peaks at $2n$ and an even $n$ does not grow at all --- is
exactly the $j = 0$ term of the sum, and it is complete for $t=1$ because no
other term contributes. It does not extend: for $t \geq 2$ an even seed with
$\vtwo(n) < t$ \emph{does} grow, which is why the higher rows sit at
$2^{\,t-1+2^{-t}}$ rather than the $2^{t/2}$ a naive reading would give.
Away from $b = 0$ the $a = 2$ row is only approximately $\sqrt{2}$ ---
Table~\ref{tab:growth} gives $1.414$, $1.417$, $1.421$ at $b = 0, 2, 6$ ---
and carykh notes this drift himself.

\begin{remark}[carykh's missing cell]
carykh reports a single anomalous blank at $6x+34$ in his heat map and
correctly diagnoses it: a seed in $3x+17$ peaks near $5.06 \times 10^{12}$, so
its image in the parallel universe $6x+34$ peaks just above his
$10^{13}$ cut-off and is misclassified as a runaway. This is an artefact of
the threshold, not of the mathematics; by Proposition~\ref{prop:halve} the two
cells must agree. It is a good illustration of why the tables here report
their budgets explicitly.
\end{remark}

% ============================================================ software =====
\section{Software}
\label{sec:software}

\paragraph{\texttt{Multiverse/multiverse.py}.}
A dependency-free library implementing the theory above. Its entry points map
one-to-one onto the results:

\smallskip
{\small
\begin{tabular}{@{}ll@{}}
\toprule
function & implements \\
\midrule
\texttt{classify} & Theorem~\ref{thm:master}: family, fate, drifts, reduction chain \\
\texttt{normal\_form} & Theorem~\ref{thm:nf} \\
\texttt{attracting\_primes} / \texttt{repelling\_primes}
  & Theorem~\ref{thm:strat} \\
\texttt{find\_cycles} & the memoised cycle census \\
\texttt{decompose\_cycles} & Corollary~\ref{cor:cycdec} \\
\texttt{cycle\_equation\_residual} & Theorem~\ref{thm:cyceq} \\
\texttt{growth\_statistics} & Table~\ref{tab:growth}, Prop.~\ref{prop:peak} \\
\bottomrule
\end{tabular}}
\smallskip

\noindent A command-line interface exposes \texttt{report}, \texttt{grid} and
\texttt{stats}.

\begin{lstlisting}[style=py,numbers=none]
$ python3 multiverse.py report 3 15
family      : attracting-reducible
fate        : open
normal form : 3x+5
   attraction  3x+15 -> 3x+5   (every orbit enters 3Z after at most 1 odd step)
attracting  : 3^1
repelling   : 5^1   (state space splits into mutually unreachable strata)
\end{lstlisting}

\paragraph{\texttt{Multiverse/qary.py}.}
The $q$-ary maps of Section~\ref{sec:qmulti} and the predicate maps of
Section~\ref{sec:beyond}: \texttt{QMap} with its two symmetries
(\texttt{scaled}, \texttt{offset\_scaled}), \texttt{residue\_graph} and
\texttt{escaping\_residues} for Theorem~\ref{thm:resobs},
\texttt{qary\_drift} for Heuristic~\ref{heur:qdrift}, and
\texttt{primality\_map} for the experiment of
Proposition~\ref{prop:primmap}.

\paragraph{\texttt{Multiverse/literature.py}.}
Classifies the ccchallenge.org snapshot into the domains of
Appendix~\ref{app:lit} and emits the appendix tables. Run with
\texttt{-{}-refresh} it re-downloads the catalogue.

\paragraph{\texttt{Multiverse/generate\_data.py} and \texttt{inline\_tables.py}.}
The first regenerates every table into \texttt{Multiverse/data/} as plain text
and LaTeX and runs the verification suite; the second splices those tables
into this paper between marker comments. The paper is a single self-contained
file with no \verb|\input| of external resources, so it compiles anywhere as a
one-file upload, and no number in it is transcribed by hand.

\paragraph{\texttt{Collatz Program/Collatz\_Program.py}.}
The existing orbital analyser gains a \emph{Multiverse} workspace
(\texttt{multiverse\_ui.py}, reached from the \textsf{Multiverse} button) with
five panels: an interactive classification grid --- Table~\ref{tab:grid} made
clickable, each cell opening its own report; a per-universe report showing
family, fate, reduction chain and strata; a cycle census displaying the
divisor decomposition and checking the cycle equation on every cycle it finds;
a symmetry explorer that places the base and transformed trajectories side by
side, marking the terms the halving symmetry inserts; and heat maps of
relative peak, path length, cycle count and drift.

% =========================================================== verification ===
\section{Verification}
\label{sec:verif}

Every theorem above is accompanied by a numerical check.
Table~\ref{tab:verif} reports them. These are checks, not proofs: they
establish that the statements are not false on a large finite range, and they
guard the implementation against the proofs drifting apart from the code.

\begin{table}[!htbp]
\centering\small
\begin{tabular}{@{}l>{\raggedright\arraybackslash}p{0.50\textwidth}rl@{}}
\toprule
check & claim & cases & result \\
\midrule
% >>>BEGIN-DATA verification
\texttt{T-halve} & $T_{2a,2b}$ refines $T_{a,b}$: one extra even term after each odd step & $341941$ & \ok \\
\texttt{T-scale} & $T_{a,db}(dn) = d\,T_{a,b}(n)$ for odd $d$ & $278400$ & \ok \\
\texttt{T-attr} & $p \mid \gcd(a,b)$ forces every orbit into $p^{v_p(b)}\mathbb{Z}$ & $16059$ & \ok \\
\texttt{T-rep} & $\min(v_p(n), v_p(b))$ is an orbit invariant when $p \mid b$, $p \nmid a$ & $73632$ & \ok \\
\texttt{T-norm} & the normal form is coprime and idempotent for $a \geq 1$; $a = 0$ needs no reduction & $39800$ & \ok \\
\texttt{T-parity} & $a+b$ odd with $a,b \geq 1$: no orbit is ever periodic & $29640$ & \ok \\
\texttt{T-desc} & $a=1$, $b$ odd: every orbit cycles, least cycle element $\leq b$ & $23940$ & \ok \\
\texttt{T-const} & $a=0$: exactly one cycle, of period $v_2(b)+1$ & $199$ & \ok \\
\texttt{T-shape} & $b=0$: every orbit cycles iff $a$ is a power of two, period $\log_2 a + 1$ & $39$ & \ok \\
\texttt{T-cyceq} & cycle equation $n_1(2^m - a^k) = b\sum_i a^{k-i}2^{d_1+\cdots+d_{i-1}}$ & $301$ & \ok \\
\texttt{T-decomp} & the $d$-graded part of $\mathrm{Cyc}(T_{a,b})$ is $d \cdot \mathrm{Cyc}^{\ast}(T_{a,b/d})$ & $33$ & \ok \\
\texttt{T-attred} & the attracting reduction is a period-preserving bijection on cycles & $30$ & \ok \\
\texttt{T-tier} & knowledge ladder: \textsc{open} $=$ minimal with an open normal form, \textsc{decided} $=$ the $a=1$ row, \textsc{closed} $=$ proved with closed-form cycles & $14400$ & \ok \\
\texttt{T-Zhalf} & $\Z_{>0}$ invariant iff $a+b \geq 1$; $\Z_{<0}$ invariant iff $b \leq a-1$ & $3660$ & \ok \\
\texttt{T-Zcross} & over $\Z$ the sign changes only at odd $n$ with $|n| \leq |b|/a$ & $1415200$ & \ok \\
\texttt{T-one} & $1$ is periodic iff its orbit meets a power of two & $3481$ & \ok \\
\texttt{T-onestep} & $a+b = 2^k$ gives a cycle through $1$ of period $k+1$ & $673$ & \ok \\
\texttt{T-desc2b} & every cycle of $x+b$ lies in $[1, 2b]$ & $381$ & \ok \\
\texttt{T-seedscale} & $n$ on a cycle of $ax+b$ $\Rightarrow$ $mn$ on one of $ax+mb$, for \emph{odd} $m$ & $104$ & \ok \\
\texttt{T-noconj} & $cC$ is never a cycle of $T_{ca,cb}$ -- the refuted conjugacy & $104$ & \ok \\
\texttt{T-radical} & $\mathrm{nf}(a,b) = (a,1)$ iff $\operatorname{rad}(b) \mid a$ & $7140$ & \ok \\
\texttt{T-qscale} & $q$-ary scaling: $F_{q\mathbf a, q\mathbf b}$ refines $F_{\mathbf a,\mathbf b}$ & $80808$ & \ok \\
\texttt{T-qtwist} & $q$-ary twisted scaling: $n \mapsto dn$ permutes the branches & $510900$ & \ok \\
\texttt{T-resobs} & escaping residues never divide, and diverge under the growth condition & $92510$ & \ok \\
\texttt{T-qattract} & attracting sublattice: the induced coefficients are $a'_j = a_{mj \bmod q}$, $b'_j = b_{mj \bmod q}/m$ -- \emph{permuted} & $91260$ & \ok \\
\texttt{T-primmap} & primality map: every seed tested is eventually periodic, with exactly two cycles & $60000$ & \ok \\
\texttt{T-neg} & $T_{a,-b}(n) = -T_{a,b}(-n)$ & $9251$ & \ok \\
% <<<END-DATA verification
\bottomrule
\end{tabular}
\caption{Numerical verification, run by \texttt{generate\_data.py}.}
\label{tab:verif}
\end{table}

% ============================================================== q-ary ======
\section{The \texorpdfstring{$q$}{q}-ary multiverse}
\label{sec:qmulti}

Everything so far concerns $q = 2$. This section asks which of it survives
when the modulus is raised, and the answer is clean: \emph{the two symmetries
survive verbatim, the parity obstruction survives in a stronger form, and the
drift threshold survives with $4$ replaced by $q^{q}$}. What does not survive
is rigidity, and with it the uniqueness of the normal form.

Recall the map of Section~\ref{sec:qary}: fix $q \geq 2$ and integers
$a_{1},\dots,a_{q-1}$, $b_{1},\dots,b_{q-1}$, and set
\[
  F(n) \;=\;
  \begin{cases}
    a_{i}n + b_{i} & n \equiv i \pmod q, \quad 1 \leq i \leq q-1,\\[1mm]
    n/q            & n \equiv 0 \pmod q .
  \end{cases}
\]
We write $F_{\mathbf a, \mathbf b}$ when the coefficients matter. This is the
shape studied by Matthews and Watts~\cite{mw1984,mw1985} following Hasse, and
by Venturini~\cite{venturini1992} under the name \emph{periodically linear}
maps; their concern is the ergodic theory, ours the elementary reductions.

\subsection{The two symmetries}

\begin{proposition}[$q$-ary scaling]\label{prop:qscale}
For every $n$,
\[
  n \not\equiv 0 \!\!\pmod q \;\Longrightarrow\;
    F_{q\mathbf a,\, q\mathbf b}(n) = q \cdot F_{\mathbf a,\mathbf b}(n)
    \ \text{and} \
    F_{q\mathbf a,\, q\mathbf b}^{2}(n) = F_{\mathbf a,\mathbf b}(n),
\]
while $F_{q\mathbf a, q\mathbf b}(n) = F_{\mathbf a,\mathbf b}(n) = n/q$ for
$n \equiv 0$. So $F_{q\mathbf a, q\mathbf b}$ traverses the same orbits as
$F_{\mathbf a,\mathbf b}$ with one extra dividing step inserted after each
affine step.
\end{proposition}

\begin{proof}
For $n \equiv i \neq 0$ the value $q(a_{i}n+b_{i})$ is divisible by $q$, so
the next step divides it back. The $n \equiv 0$ case is identical on both
sides.
\end{proof}

\begin{proposition}[Twisted scaling]\label{prop:qtwist}
Let $\gcd(d,q) = 1$ and let $\sigma(i) = di \bmod q$, a permutation of
$(\Z/q)^{\times}$-orbits on $\{1,\dots,q-1\}$. Define
$a'_{\sigma(i)} = a_{i}$ and $b'_{\sigma(i)} = d\, b_{i}$. Then for every $n$,
\[
  F_{\mathbf a', \mathbf b'}(dn) \;=\; d \cdot F_{\mathbf a, \mathbf b}(n).
\]
\end{proposition}

\begin{proof}
If $n \equiv 0 \pmod q$ then $dn \equiv 0$ and both sides are $dn/q$. If
$n \equiv i \neq 0$ then $dn \equiv \sigma(i) \neq 0$, because $d$ is
invertible mod $q$, and
$F_{\mathbf a',\mathbf b'}(dn) = a'_{\sigma(i)}dn + b'_{\sigma(i)}
= d(a_{i}n + b_{i})$.
\end{proof}

For $q = 2$ the only choice is $d$ odd, $\sigma$ is the identity, and
Proposition~\ref{prop:qtwist} is exactly Proposition~\ref{prop:scale}. For
$q \geq 3$ something genuinely new appears: the scaling \emph{permutes the
branches}, so a universe is equivalent to one with its coefficient vector
rearranged. Rigidity (Theorem~\ref{thm:rigid}) therefore fails for
$q \geq 3$: two different coefficient vectors can describe the same dynamics.

\begin{proposition}[Attracting sublattice]\label{prop:qattract}
Let $m \geq 2$ with $\gcd(m,q) = 1$, and suppose $m \mid a_{i}$ and
$m \mid b_{i}$ for every $i$. Then $m\Z$ is invariant, every orbit enters it
after its first affine step, and $n \mapsto n/m$ carries $F$ restricted to
$m\Z$ to the map $F'$ with coefficients
\[
  a'_{j} \;=\; a_{\,mj \bmod q}, \qquad
  b'_{j} \;=\; \frac{b_{\,mj \bmod q}}{m},
  \qquad 1 \leq j \leq q-1 .
\]
\end{proposition}

\begin{proof}
\emph{Invariance and absorption.} If $m \mid n$ and $q \mid n$ then
$m \mid n/q$, because $\gcd(m,q)=1$; and $a_{i}n + b_{i} \equiv 0 \pmod m$ for
every $n$ and every $i$, by hypothesis. The second statement also gives the
absorption: the value after the first affine step already lies in $m\Z$.

\emph{The induced map.} Let $k \in \Z$ and put $n = mk$. Since
$\gcd(m,q) = 1$, $n \equiv 0 \pmod q$ if and only if $k \equiv 0 \pmod q$,
and then $F(n)/m = (n/q)/m = k/q$. Otherwise $k \equiv j \not\equiv 0$ and
$n \equiv mj \not\equiv 0$, so $F$ applies the branch indexed by
$mj \bmod q$:
\[
  \frac{F(mk)}{m} \;=\; \frac{a_{\,mj}\,mk + b_{\,mj}}{m}
  \;=\; a_{\,mj}\,k + \frac{b_{\,mj}}{m},
\]
which is the displayed coefficient vector.
\end{proof}

\begin{remark}[Why the coefficients must be permuted]\label{rem:qattract}
It is tempting to write the induced coefficients as $(a_{j}, b_{j}/m)$ ---
scaling the offsets and leaving the multipliers where they are. That is
\emph{false} for $q \geq 3$, and an earlier draft of this paper asserted it.
Multiplication by $m$ permutes the nonzero residue classes, so the branch $F$
applies to $mk$ is the one indexed by $mj$, not by $j$. For $q = 3$, $m = 2$,
$\mathbf a = (2,4)$, $\mathbf b = (2,2)$ and $k = 1$ one gets
$F(2)/2 = 10/2 = 5$, whereas the naive vector
$(a_{j}, b_{j}/2) = ((2,4),(1,1))$ gives $2 \cdot 1 + 1 = 3$.

The error is invisible at $q = 2$, where there is a single nonzero class and
every permutation of it is the identity --- which is why nothing in
Sections~\ref{sec:sym}--\ref{sec:open} is affected. It is the same phenomenon
as the failure of rigidity noted after Proposition~\ref{prop:qtwist}: at
$q \geq 3$ a scaling does not merely rescale a universe, it relabels its
branches.
\end{remark}

\subsection{The residue obstruction}

For $q = 2$, ``$a+b$ odd'' says an odd number can never become even. The right
generalisation is a finite directed graph.

\begin{definition}\label{def:resgraph}
The \emph{residue graph} of $F$ is the map
$\tau : \{1,\dots,q-1\} \to \Z/q$, $\tau(i) = a_{i}i + b_{i} \bmod q$. Call
$i$ \emph{escaping} if the forward $\tau$-orbit of $i$ never reaches $0$.
\end{definition}

The graph is well defined because $n \equiv i$ forces
$a_{i}n + b_{i} \equiv a_{i}i + b_{i} \pmod q$. Note it captures only the
affine branches: $n/q$ is \emph{not} determined by $n \bmod q$, since it
depends on $n \bmod q^{2}$. That asymmetry is exactly why the obstruction is
one-directional.

\begin{theorem}[Residue obstruction]\label{thm:resobs}
Let $i$ be escaping and let $n \equiv i \pmod q$. Then:
\begin{enumerate}\itemsep2pt
\item[(i)] the orbit of $n$ never meets the dividing branch;
\item[(ii)] if in addition every residue $j$ on the forward $\tau$-orbit of
      $i$ satisfies $a_{j} \geq 2$, or $a_{j} = 1$ and $b_{j} \geq 1$, then
      the orbit of every $n \geq 1$ in the class of $i$ is strictly
      increasing, hence divergent.
\end{enumerate}
\end{theorem}

\begin{proof}
(i) By induction: the residue after one step is $\tau$ of the current one, and
the $\tau$-orbit of $i$ avoids $0$ by hypothesis, so every term is in a
non-zero class and the dividing branch is never selected.

(ii) By (i) every step is affine, using the branch $j$ of the current residue.
The stated condition is exactly $a_{j}n + b_{j} > n$ for all $n \geq 1$: for
$a_{j} \geq 2$ it gives $a_{j}n + b_{j} \geq 2n > n$, and for $a_{j} = 1$ it
gives $n + b_{j} > n$. So the orbit strictly increases at every step and is
unbounded.
\end{proof}

\begin{corollary}
For $q = 2$ the only affine residue is $i = 1$, $\tau(1) = a + b \bmod 2$, and
$1$ is escaping iff $a+b$ is odd. Theorem~\ref{thm:resobs} is then
Theorem~\ref{thm:parity}.
\end{corollary}

The two conditions in Theorem~\ref{thm:resobs} are both needed. Dropping the
growth condition admits the identity branch $a_{j}=1$, $b_{j}=0$, whose orbits
are fixed points rather than divergent, and admits $a_{j}=0$, whose branch is
constant. A machine check over $q \leq 5$ and coefficients in
$\{0,\dots,3\}$ found $239{,}530$ orbits from escaping residues, none of which
ever divided, confirming (i); and $125{,}765$ orbits satisfying the growth
condition, all strictly increasing, confirming (ii).

\subsection{Drift}

\begin{heuristic}\label{heur:qdrift}
Suppose the residues equidistribute along orbits and that an affine step
whose image is divisible by $q$ is followed by $q/(q-1)$ divisions on average.
Then the geometric drift per affine step is
\[
  \delta \;=\;
  \frac{\bigl(\prod_{i=1}^{q-1} a_{i}\bigr)^{1/(q-1)}}{q^{\,q/(q-1)}},
\]
so $F$ contracts on average precisely when
\[
  \prod_{i=1}^{q-1} a_{i} \;<\; q^{\,q}.
\]
\end{heuristic}

For $q = 2$ this is $a < 4$, the threshold of Section~\ref{sec:open}, and it
recovers $\delta = a/4$. For $q = 3$ the condition is $a_{1}a_{2} < 27$: the
map with $a = (3,3)$ has $\delta \approx 0.577$ and should collapse, while
$a = (5,7)$ has $\delta \approx 1.139$ and should escape. As at $q = 2$ this
is a heuristic and nothing more; the equidistribution it assumes is not known
for any single map.

\begin{remark}
The threshold $q^{q}$ grows so fast that the ``interesting'' region --- where
the drift is near $1$ --- is a vanishing fraction of coefficient space as $q$
grows, but it is never empty: taking $a_{i} \approx q^{q/(q-1)}$ for all $i$
puts $\delta$ near $1$. The $q$-ary multiverse therefore has a Goldilocks
region at every modulus, and $3x+1$ is merely its smallest member.
\end{remark}

% =========================================================== beyond mod ===
\section{Beyond modularity}
\label{sec:beyond}

Modularity is doing an enormous amount of work in everything above. This
section asks what happens without it. The answer is that none of the
machinery above survives, though the maps themselves need not be
structureless; the reason is interesting, and it suggests a principle we
state as a conjecture.

\begin{definition}[Predicate map]\label{def:predmap}
Let $P_{1},\dots,P_{k}$ be predicates on $\Z$ with $P_{k}$ always true, and
let $f_{1},\dots,f_{k}$ be functions $\Z \to \Z$. The associated
\emph{predicate map} sends $n$ to $f_{j}(n)$ for the least $j$ with
$P_{j}(n)$.
\end{definition}

Every map $\Z \to \Z$ is a predicate map, so there is no classification
theorem to be had; and by Conway~\cite{conway} even the affine, congruence-driven
case is universal once enough branches are allowed. The question worth asking
is therefore not ``what are the predicate maps like'' but \emph{which
predicates preserve the theory}.

\begin{proposition}[Congruence predicates change nothing]\label{prop:congpred}
If every $P_{j}$ is a union of residue classes modulo some $M_{j}$, and every
$f_{j}$ is affine, then the predicate map is a periodically linear map with
modulus $M = \operatorname{lcm}(M_{1},\dots,M_{k})$, and if one branch is
$n \mapsto n/q$ on the class $0 \bmod q$ with $q \mid M$ it is a $q$-ary map
after refining the classes. All of Section~\ref{sec:qmulti} applies.
\end{proposition}

\begin{proof}
Each $P_{j}$ is determined by $n \bmod M$, so the branch taken is a function
of $n \bmod M$; grouping the classes gives the stated shape.
\end{proof}

Proposition~\ref{prop:congpred} is one direction only: \emph{periodic}
predicates --- unions of residue classes --- land inside the theory. An
\emph{eventually} periodic predicate, one that is periodic beyond some
threshold, differs from a periodic one on a finite set, so it gives the same
map outside a finite region; the finitely many exceptional integers must then
be handled by inspection, and they can create or destroy cycles, though they
cannot affect the asymptotics. The converse --- that nothing outside this
class admits a reduction theory --- is Conjecture~\ref{conj:boundary}, which
we do not prove.

What we can say is which of \emph{our} tools break. Once a predicate is not a
congruence condition --- primality, squarefreeness, being a perfect power, a
digit condition in a base coprime to the branch moduli --- the residue graph
does not exist, Terras-style class refinement has nothing to refine, and the
affine bookkeeping behind the cycle equation has no periodic structure to
lean on. That is a statement about this paper's machinery, not about the
maps: a non-congruential map may well have arithmetic structure of some other
kind, and the primality map below has a good deal of it.

\subsection{A worked example: the primality map}

The programme that produced this paper experimented with
\[
  f(n) \;=\;
  \begin{cases}
    3n+1 & n \text{ prime},\\
    n/2  & n \text{ even},\\
    n+1  & \text{otherwise},
  \end{cases}
\]
and observed that every orbit it tried was eventually periodic. We can say
rather more.

\begin{proposition}\label{prop:primmap}
For the map above: every odd prime $p$ satisfies $f^{2}(p) = (3p+1)/2
\approx 1.5\,p$, and every odd composite $m$ satisfies $f^{2}(m) = (m+1)/2
\approx 0.5\,m$. Hence an orbit can diverge only if the proportion $r$ of
\emph{prime} values among the odd numbers it visits satisfies
\[
  \Bigl(\tfrac32\Bigr)^{r}\Bigl(\tfrac12\Bigr)^{1-r} > 1,
  \qquad\text{i.e.}\qquad
  r \;>\; \frac{\ln 2}{\ln 3} \;=\; 0.63092\ldots
\]
\end{proposition}

\begin{proof}
An odd $n$ is sent to $3n+1$ or $n+1$, both even, and the next step halves.
Taking logarithms of the product of the per-odd-step multipliers gives the
stated condition.
\end{proof}

A search over all seeds $n \leq 200{,}000$ finds that \emph{every} orbit is
eventually periodic, and that there are exactly two cycles:
\begin{align*}
  C_{1} &= \{2,4,5,7,8,10,11,13,16,17,20,22,26,34,40,52\},
    &&\text{period } 16, &&86.9\% \text{ of seeds},\\
  C_{2} &= \{31,41,47,62,71,81,82,94,107,124,142,161,162,214,322\},
    &&\text{period } 15, &&13.1\% .
\end{align*}

The tempting argument --- ``primes have density $0$, so the growing branch is
negligible'' --- is too quick, and measuring carelessly makes it look false.
Counting prime values over whole orbits gives densities of $55$--$71\%$,
alarmingly close to the threshold, because the two cycles are short and
prime-rich and orbits spend most of their time inside them. Restricting to
the transient part and to values above a floor gives the honest picture:

\begin{center}\small
\begin{tabular}{@{}rrrr@{}}
\toprule
values above & odd values seen & prime & density \\
\midrule
$10^{3}$   & $3958$ & $1328$ & $33.6\%$ \\
$10^{4}$   & $3500$ & $723$  & $20.7\%$ \\
$10^{6}$   & $2481$ & $344$  & $13.9\%$ \\
$10^{8}$   & $3177$ & $331$  & $10.4\%$ \\
$10^{11}$  & $1270$ & $107$  & $8.4\%$ \\
\bottomrule
\end{tabular}
\end{center}

\noindent
The density decays, but it sits consistently around $2$--$3$ times the
$1/\ln n$ one would predict for random integers of that size ($3.6\%$ at
$10^{12}$) --- the map's odd values are biased towards primes, because
$(m+1)/2$ and $(3p+1)/2$ appear to avoid small factors more often than a
random integer of the same size does. We offer that as a description of what
we measured, not as an explanation: we have not analysed the distribution of
small factors along these orbits, and the bias may have another cause. The bias is real and worth noting; it is also nowhere near enough. At
$10^{12}$ the observed $8.4\%$ must reach $63.1\%$ for divergence, and it is
falling.

\begin{conjecture}[Primality map]\label{conj:primmap}
Every orbit of $f$ is eventually periodic, and the two cycles listed above are
the only ones.
\end{conjecture}

\subsection{The density--drift principle}

The computation above is an instance of something general, which we state as
the organising conjecture of this section.

\begin{conjecture}[Density--drift principle]\label{conj:density}
Let $T$ be a predicate map whose branches are asymptotically multiplicative,
$f_{j}(n) \sim c_{j}n$, and suppose that along orbits the branch $j$ is taken
with asymptotic frequency $d_{j}$, $\sum_j d_{j} = 1$. Then almost every orbit
is bounded if $\prod_{j} c_{j}^{\,d_{j}} < 1$ and almost every orbit diverges
if $\prod_{j} c_{j}^{\,d_{j}} > 1$.
\end{conjecture}

For congruence predicates the frequencies $d_{j}$ are the densities of the
classes and the principle specialises to Heuristic~\ref{heur:qdrift} and, at
$q=2$, to the classical $a/4$. Its content in the non-modular case is the
assertion that the branch frequencies along orbits exist at all --- which for
the primality map is the statement that the orbit does not conspire to visit
primes more often than a density argument allows, and which the table above
supports without settling.

\begin{conjecture}[Sparse growth]\label{conj:sparse}
If the branches with $c_{j} > 1$ are selected by predicates whose density
along orbits tends to $0$, and some branch with $c_{j} < 1$ is taken with
frequency bounded below, then every orbit is bounded.
\end{conjecture}

Conjecture~\ref{conj:sparse} would give Conjecture~\ref{conj:primmap}
immediately, since the primes have density $0$. We state it separately
because it is the weakest statement we can see that does the job, and because
it isolates the difficulty: it is a statement about the \emph{invariant
measure} of the map, not about the density of primes in $\N$.

\begin{conjecture}[Reduction boundary]\label{conj:boundary}
A predicate map admits a nontrivial reduction theory --- symmetries of the
kind in Section~\ref{sec:sym}, a normal form, and a finite obstruction graph
--- if and only if its predicates are eventually periodic.
\end{conjecture}

The forward direction is Proposition~\ref{prop:congpred}. The converse is the
substantial half, and Conway's universality~\cite{conway} together with the
$\Pi^{0}_{2}$-completeness of Kurtz and Simon~\cite{kurtzsimon} is evidence
for it: a family rich enough to encode arbitrary computation cannot admit a
classification, and predicate maps are certainly that rich. What would make
Conjecture~\ref{conj:boundary} a theorem is a precise definition of
``reduction theory'', which we do not have.

\begin{remark}
This section is deliberately short. The point we want to record is
structural: the reason $3x+1$ is attackable at all --- the reason there is a
Terras argument, a cycle equation, a normal form --- is that its branch
predicate is a congruence. Change that, and the problem does not become
harder in degree; it changes kind, from number theory to computability.
\end{remark}

% ========================================================== conjectures ====
\section{Conjectures, and which of them survive}
\label{sec:conj}

The programme that produced this paper accumulated a list of conjectures. It
is worth publishing the list together with its casualties: three of them are
false, and the counterexamples are small enough to check by hand.

\begin{enumerate}\itemsep6pt

\item[\textbf{C1}] \emph{(Generalised Crandall conjecture.)} For every
      $b \geq 1$ and every odd $a > 3$, the map $\Tab$ has an orbit that
      never reaches a cycle through $1$. \textbf{Open.} For $b = 1$ this is
      Crandall's conjecture~\cite{crandall}, known for almost all $a$ by
      Franco and Pomerance~\cite{fp} and for $a = 2^{k}-1$ by
      Santos~\cite{santos}. By Theorem~\ref{thm:redfate} it suffices to
      consider $(a,b)$ in normal form, and by Corollary~\ref{cor:cycdec} the
      cycle structure for composite $b$ is determined by that of its divisors
      --- so the conjecture really only concerns primitive universes.

\item[\textbf{C2}] \emph{Every orbit of every map with $a < 3$ is bounded.}
      \textbf{False as stated, and completely settled.} It fails for
      $a = 2$, $b$ odd --- $2x+1$ sends $1 \to 3 \to 7 \to 15 \to \cdots$ ---
      which is Theorem~\ref{thm:parity}. Restricted to the cases where it is
      not refuted by parity it is a theorem:
      $a = 0$ is Proposition~\ref{prop:const}, $a = 1$ with $b$ odd is
      Theorem~\ref{thm:descent}, and $a = 2$ with $b$ even reduces to
      $a = 1$ by Proposition~\ref{prop:halve}.

\item[\textbf{C3}] \emph{The $q$-ary map of Section~\ref{sec:qary} is not
      Turing-complete.} \textbf{False for the unrestricted family.}
      Conway~\cite{conway} constructed universal generalized Collatz
      functions, and Kurtz and Simon~\cite{kurtzsimon} showed the natural
      reachability problem is $\Pi^{0}_{2}$-complete. We note honestly that
      Conway's construction uses rational multipliers chosen so that the
      image is an integer, which is more freedom than the shape in
      Section~\ref{sec:qary} allows; whether that exact restricted shape is
      universal we do not know, and we regard \textbf{C3} as open only in
      that narrow reading.

\item[\textbf{C4}] \emph{If $1$ lies on a cycle $C$ containing another odd
      number, and no scaling relation applies, then no odd number outside $C$
      enters $C$.} \textbf{False.} In $5x+1$ the cycle through $1$ is
      $\{1,6,3,16,8,4,2\}$, which contains the odd number $3$; the offset
      $b=1$ admits no scaling; and the odd number $15 \notin C$ enters it:
      \[
        15 \to 76 \to 38 \to 19 \to 96 \to 48 \to 24 \to 12 \to 6 \in C .
      \]
      A search over odd $a \leq 39$, odd $b \leq 39$ and odd seeds below
      $600$ finds eight such violations, the smallest in the most-studied
      universe of the whole family.

\item[\textbf{C5}] \emph{$1$ lies on a cycle of every map $3x+b$.}
      \textbf{False}, as the programme already recorded: in $3x+3$ the orbit
      of $1$ is $1 \to 6 \to 3 \to 12 \to 6$, which cycles without returning
      to $1$.

\item[\textbf{C5$'$}] \emph{$1$ lies on a cycle of $3x+b$ whenever $b$ is not
      a power of $3$.} \textbf{False.} The smallest counterexample is
      $b = 7$: the orbit of $1$ is
      $1 \to 10 \to 5 \to 22 \to 11 \to 40 \to 20 \to 10$, which enters the
      cycle $\{10,5,22,11,40,20\}$ and never returns to $1$. Further
      counterexamples are $b = 15, 19, 21, 23, 25, 31, 33, 35, 37, 39, 45,
      \dots$ Theorem~\ref{thm:one} replaces the conjecture with an exact
      criterion: $1$ lies on a cycle of $3x+b$ if and only if the orbit of
      $1$ meets a power of two, and $3 + b = 2^{k}$
      \textup{(}Corollary~\ref{cor:onestep}\textup{)} is the one-step case
      that accounts for $93$ of the $147$ instances we found.

\item[\textbf{Q1}] \emph{Does $3511x+1$ have an orbit avoiding $1$?} $3511$
      is the second known Wieferich prime; Crandall's conjecture is known for
      the first, $1093$~\cite{crandall}. \textbf{Open}, and outside the reach
      of anything here: our reductions never change the multiplier
      \textup{(}Theorem~\ref{thm:rigid}\textup{)}, so the row $a = 3511$ is
      untouched by them.
\end{enumerate}

\begin{remark}
The pattern in \textbf{C4} and \textbf{C5$'$} is worth naming. Both
conjectures generalise from the universes where a scaling relation is doing
the work, and both break exactly where it stops. This is the practical value
of Theorem~\ref{thm:nf}: once a universe is in normal form, one knows that no
further structure is hiding, and a pattern that persists there is a pattern
about the universe rather than about its presentation.
\end{remark}

% ============================================================== limits =====
\section{What is \emph{not} proved}
\label{sec:limits}

\begin{enumerate}\itemsep4pt
\item \textbf{No primitive universe is resolved.} Theorem~\ref{thm:master}
      case (4) is a statement of ignorance. Nothing here bears on whether
      $3x+1$ has other cycles, whether any $3x+b$ orbit diverges, or whether
      $5x+1$ has a divergent orbit. The reductions move problems around; they
      do not solve any.

\item \textbf{Cycle counts are lower bounds.} Every count in
      Section~\ref{sec:cycles} comes from a search over a finite seed window
      with a finite step budget. A cycle whose basin misses the window is not
      found. We report the window with every table and we checked stability by
      widening it, but stability under widening is not a proof of
      completeness. Artem makes the same caveat; it is a real one, and
      $b = 71$ in Table~\ref{tab:rowa3} is a concrete example of a count that
      changes when the window grows.

\item \textbf{The drifts are heuristics.} $\delta_{\mathrm{odd}} = a/4$ rests
      on $\mathbb{E}[\vtwo(an+b)] = 2$, which assumes an equidistribution that
      is not proved for any particular universe. It predicts nothing about
      any individual orbit. Tao's result~\cite{tao} is the strongest rigorous
      statement we know of in this direction, on
      making such heuristics rigorous for $3x+1$, and it falls well short of
      the conjecture.

\item \textbf{Densities are natural densities.} Theorem~\ref{thm:density}
      says nothing about any particular $(a,b)$, and in particular does not
      say that a positive proportion of universes will remain open --- only
      that a positive proportion is open \emph{now}.

\item \textbf{The $q$-ary chapter is elementary.} Section~\ref{sec:qmulti}
      proves only the two symmetries and the residue obstruction. It does not
      construct a normal form for $q \geq 3$, and it cannot: rigidity fails
      there, because twisted scaling permutes the branches
      (Proposition~\ref{prop:qtwist}), so a universe has several equally
      valid coefficient vectors. Building the analogue of
      Theorem~\ref{thm:nf} modulo that permutation action is the obvious next
      problem, and we have not solved it.

\item \textbf{Section~\ref{sec:beyond} proves almost nothing.} Its content is
      one proposition, one computation, and three conjectures. In particular
      Conjecture~\ref{conj:density} is not even precisely stated --- ``the
      branch frequency along orbits'' is not known to exist --- and
      Conjecture~\ref{conj:boundary} awaits a definition of ``reduction
      theory''. We include the chapter because the boundary it points at
      seems to us the right one, not because we can defend it.

\item \textbf{Only affine equivalences are classified.}
      Theorem~\ref{thm:rigid} rules out affine intertwinings. It does not rule
      out more exotic conjugacies, and it does not rule out that two rows are
      related by an argument that is not a change of variable at all.

\item \textbf{$b < 0$ is only transported, not studied.}
      Proposition~\ref{prop:neg} reduces negative offsets to the negative
      integers of a positive-offset universe, but we do not classify the
      dynamics there; $3x-1$ has at least three cycles and is as open as
      $3x+1$.
\end{enumerate}

% ============================================================= attribution ==
\section{Attribution}
\label{sec:attrib}

The multiverse framing and a large fraction of the phenomena in this paper
were found, experimentally and independently of the literature, by
carykh~\cite{carykh,lazykh} and by Artem~\cite{artem}. This paper's
contribution is to prove them, to unify them, and to draw the line between
solved and open. Table~\ref{tab:attrib} attributes each result.

\begin{table}[!htbp]
\centering\footnotesize
\begin{tabular}{@{}R{0.44\textwidth}R{0.14\textwidth}R{0.34\textwidth}@{}}
\toprule
observation & first stated & status here \\
\midrule
The multiverse framing; the $(a,b)$ grid & carykh~\cite{carykh}
  & adopted \\
Hub-and-spoke family, $a=0$ & carykh
  & Prop.~\ref{prop:const}: proof, exact period $\vtwo(b)+1$, uniqueness \\
Shape family $b=0$ with $\log_{2}(a)+1$ sides & carykh
  & Prop.~\ref{prop:shape}: proof, and the $u \geq 3$ runaway \\
Runaway checkerboard, $a+b$ odd; linear vs exponential subtypes & carykh
  & Thm.~\ref{thm:parity}: proof \\
Even parallel universes $(2a,2b) \sim (a,b)$ & carykh
  & Prop.~\ref{prop:halve}: exact orbit refinement, period formula, cycle
    bijection \\
Odd parallel universes, offset scaled by odd $d$ & carykh
  & Prop.~\ref{prop:scale}: proof \\
Connected vs disconnected odd parallels & carykh (named), Artem (criterion)
  & Thm.~\ref{thm:strat}, Cor.~\ref{cor:conn}: the criterion is $p \mid a$,
    per odd prime, with absorption time $\vp(b)$ \\
``Factor loyalty'' on the odd diagonal $a=b$ & carykh
  & special case of Thm.~\ref{thm:strat}(i) \\
$3x+9$ attracts in two odd steps; $3x+3^{j}$ in $j$ & Artem
  & Thm.~\ref{thm:strat}(i) \\
$3x+15$ is a copy of $3x+5$, not of $3x+1$ & Artem
  & Thm.~\ref{thm:nf}, Thm.~\ref{thm:redfate} \\
Composite offsets contain a copy of every divisor & Artem
  & Thm.~\ref{thm:strata}, Cor.~\ref{cor:cycdec}: strengthened to an exact
    partition of the phase space and of the cycle set \\
Cycle counts for $3x+b$, $b < 100$ & Artem
  & Table~\ref{tab:rowa3}: recomputed at a wider, empirically stable
    window \\
Only powers of $3$ away from a base have one cycle & Artem
  & Cor.~\ref{cor:pow3} \\
No $b < 100$ gives exactly seven cycles & Artem
  & confirmed at a wider window (\S\ref{sec:cycles}) \\
Drift $\sqrt{a}/2$; Goldilocks / Icarus / black-hole rows & carykh
  & \S\ref{sec:open}, with the equivalent form $a/4$ per odd step and the
    threshold $a<4$ \\
$(k,m) = (3,5)$ and $(17,27)$ in $3x+5$, from $3^{k} \approx 2^{m}$ & carykh
  & Thm.~\ref{thm:cyceq} \\
$5x+1$ and $5x+5$ share the fingerprint $7,10,10$ & carykh
  & Prop.~\ref{prop:scale} + Thm.~\ref{thm:strat}(i) \\
Average max growth heat map; $\sqrt{2}$ on the $a=2$ row & carykh
  & Table~\ref{tab:growth} \\
The $6x+34$ anomaly is a threshold artefact & carykh
  & confirmed, \S\ref{sec:open} \\
Nine ``unique'' universes in the $6\times6$ grid & carykh
  & Thm.~\ref{thm:nine}: three solved, two reducible, four open \\
\midrule
\multicolumn{3}{@{}l@{}}{\emph{From the authors' own earlier programme
(research notes and the repository's \texttt{proofs.txt}):}} \\
Hub-and-spoke ($a=0$), shapes ($b=0$), parity runaways, and the $a=1$ row
  & the authors
  & Props.~\ref{prop:const},~\ref{prop:shape}, Thms.~\ref{thm:parity},
    \ref{thm:descent}: proofs, plus two corrections noted below \\
$n$ seeds a cycle of $ax+b$ $\Rightarrow$ $mn$ seeds one of $ax+mb$
  & the authors
  & Prop.~\ref{prop:scale}; \textbf{needs $m$ odd} --- see below \\
$a = 2^{k}-b$ gives a cycle through $1$
  & the authors
  & Cor.~\ref{cor:onestep}, and strengthened to an exact criterion in
    Thm.~\ref{thm:one} \\
Every cycle of $x+b$ lies in $[1,2b]$
  & the authors
  & Thm.~\ref{thm:descent}(iii); this is what makes Table~\ref{tab:rowa1}
    exact \\
$2ax+2b$ is equivalent to $ax+b$
  & the authors, and carykh independently
  & Prop.~\ref{prop:halve} \\
$n$ in $ax+b$ corresponds to $n/b$ in $ax+1$; equivalence iff $b \mid a$
  & the authors
  & Prop.~\ref{prop:scale}, Prop.~\ref{prop:onestep}; the correct criterion
    for equivalence is $\operatorname{rad}(b) \mid a$
    (Thm.~\ref{thm:radical}) \\
Post-composing the odd branch: $T_{g}(n) = g(T(n))$ for odd $n$
  & the authors
  & Prop.~\ref{prop:postcomp} \\
$g$ conjugates $T$ to $T_{g}$ when $\vtwo$ is preserved
  & the authors
  & \textbf{refuted}, Prop.~\ref{prop:noconj} \\
The conjecture list C1--C5$'$, Q1
  & the authors
  & \S\ref{sec:conj}: C2, C4 and C5$'$ refuted, C3 refuted in the
    unrestricted reading \\
\midrule
The $q$-ary framing (mod $q$ rather than mod $2$), and the programme of
pushing the theory beyond modularity & the authors
  & \S\ref{sec:qmulti}, \S\ref{sec:beyond}; the $q$-ary shape itself is
    Matthews--Watts~\cite{mw1984} \\
The primality-driven map $3n+1$ / $n/2$ / $n+1$ & the authors
  & Prop.~\ref{prop:primmap}, Conj.~\ref{conj:primmap}: the threshold
    $\ln 2/\ln 3$ and the measured densities \\
\midrule
Rigidity of affine equivalences & ---
  & Thm.~\ref{thm:rigid} \\
Normal form; termination and confluence & ---
  & Thm.~\ref{thm:nf}, Thm.~\ref{thm:redfate} \\
Master classification & ---
  & Thm.~\ref{thm:master} \\
Densities $1/3$ and $2/\pi^{2}$ & ---
  & Thm.~\ref{thm:density} \\
$q$-ary scaling, twisted scaling, residue obstruction, threshold $q^{q}$ & ---
  & Props.~\ref{prop:qscale},~\ref{prop:qtwist}, Thm.~\ref{thm:resobs},
    Heur.~\ref{heur:qdrift} \\
The density--drift principle and the reduction boundary & ---
  & Conjs.~\ref{conj:density},~\ref{conj:sparse},~\ref{conj:boundary} \\
$a=1$ solved with the bound $\min C \leq b$ & ---
  & Thm.~\ref{thm:descent} \\
\bottomrule
\end{tabular}
\caption{Attribution. Entries with ``---'' in the middle column are, to our
knowledge, first stated here; we make no claim of priority against the
research literature on generalised $3x+1$ maps, for which see
Lagarias~\cite{lagarias,lagariasbook}.}
\label{tab:attrib}
\end{table}

\begin{remark}[Corrections to the earlier notes]
Five items needed amending, and all five are recorded above and in
\S\ref{sec:conj}.
\begin{enumerate}\itemsep2pt
\item The $b=0$ cycle length is $\log_{2}(a)+1$, not $\log_{2}(a)+2$
      (Proposition~\ref{prop:shape}); the discrepancy is the convention of
      listing the starting term twice.
\item The $a=1$ family was recorded as having ``two loops''. The two named
      always exist, but there are frequently more, first at $b=7$
      (Remark~\ref{rem:a1}).
\item The cycle-seed scaling $n \mapsto mn$ requires $m$ \emph{odd}. For even
      $m$ the image $mn$ is even and simply halves; a search over
      $a \in \{3,5,7\}$, $b \leq 7$ found $140$ failures for even $m$ and none
      for odd $m$.
\item The claim that a $\vtwo$-preserving $g$ conjugates $T$ to $T_{g}$ is
      false (Proposition~\ref{prop:noconj}); $0$ of $104$ test cases lift.
\item Conjectures C2, C4 and C5$'$ are false, with counterexamples
      $2x+1$, $5x+1$ and $3x+7$ respectively (\S\ref{sec:conj}).
\end{enumerate}
None of these affects the reductions or the classification.
\end{remark}

\begin{remark}
Several of the results here are certainly known in some form to specialists.
The cycle equation (Theorem~\ref{thm:cyceq}) is classical, appearing in
Lagarias~\cite{lagarias} for $3x+1$. The scaling relation
(Proposition~\ref{prop:scale}) is folklore. What we have not found stated
anywhere is the combination: the normal form, the resulting master
classification, and the two densities.
\end{remark}

% ============================================================ conclusion ====
\section{Conclusion}

The generalised Collatz multiverse is more structured than it looks and
smaller than it looks. Two symmetries --- halving and odd scaling --- and one
dichotomy --- whether an odd prime divides the multiplier or only the offset
--- reduce every universe to a coprime normal form. Four of the five resulting
cases are solved completely and elementarily. The fifth, $a \geq 3$ odd with
$b$ odd and coprime, is open, has density $2/\pi^{2}$, and contains
$3x+1$.

Pushing outward sharpens rather than dilutes that picture. Raising the
modulus keeps both symmetries and replaces the parity obstruction by a finite
directed graph on $\Z/q$, but costs rigidity: scaling permutes the branches,
so there is no normal form at $q \geq 3$ and we do not know what should
replace it. Dropping modularity altogether costs everything. The reduction
theory, the cycle equation and the Terras-style refinement all rest on the
branch being selected by a congruence, and once it is not --- primality, say
--- the family is rich enough to encode arbitrary computation, so no
classification can exist. On that reading the boundary of the subject is not
a difficulty of degree but a change of kind, and it runs exactly along
periodicity of the branch predicate.

The videos that prompted this paper end with carykh's speculation that Collatz
chose $3x+1$ ``not because that one was particularly interesting, but because
it was the only rule set that wasn't obvious in some way''. The correct
version is narrower and, we think, more striking: the drift threshold
$a < 4$ admits exactly two odd multipliers, one of which is solved by a
two-line descent argument. $3x+1$ is not the only interesting universe, but
it is the smallest member of the only row where the difficulty is of the
Collatz kind.

% ============================================================= AI note ======
\section*{Note on the use of artificial intelligence}
\label{sec:ai}

The research programme, the earlier hand computations, and the existing
software in the repository are the author's. The theorems and proofs in this
paper, the library \texttt{multiverse.py}, the verification suite, the
Multiverse workspace in \texttt{Collatz\_Program.py}, and the text of this
paper were produced in collaboration with an artificial-intelligence system
(Claude, Anthropic), at the author's direction.

The proofs here are hand-checkable and short by design, and each is
accompanied by an independent numerical check reported in
Table~\ref{tab:verif}; the cycle censuses were cross-checked against the
counts published by Artem~\cite{artem} and carykh~\cite{carykh,lazykh}, which
agree in every case where the search windows overlap. Section~\ref{sec:limits}
states explicitly what is not proved. Unlike the companion
paper~\cite{sievepaper} in this repository, this development is \emph{not}
machine-checked in a proof assistant; the reader should hold its theorems to
the ordinary standard of a hand-written proof.

% ============================================================= appendix =====
\FloatBarrier
\appendix
\section{The ccchallenge corpus, classified}
\label{app:lit}

ccchallenge.org~\cite{ccc} is a collaborative project to formalise the
Collatz literature in proof assistants; it maintains a catalogue with a public
API. We took a snapshot of $375$ papers and classified them by keyword into
the domains below, first match winning, so that each paper is counted once.
The classification is produced by \texttt{Multiverse/literature.py} from the
cached snapshot \texttt{data/ccchallenge\_papers.json}, and the ordering of
the domain list is significant: generalised maps are tested first, because
those are the papers this paper engages with.

Table~\ref{tab:litdomains} gives the resulting counts.

\begin{table}[!htbp]
\centering\small
\begin{tabular}{@{}lr@{}}
\toprule
domain & papers \\
\midrule
% >>>BEGIN-DATA lit_domains
generalised maps & 49 \\
cycles and lower bounds & 32 \\
stopping time and density & 15 \\
computation and verification & 53 \\
logic, automata, formalisation & 8 \\
p-adic, ergodic, analytic & 13 \\
surveys and bibliographies & 8 \\
other $3x+1$ & 197 \\
\midrule \textbf{total} & \textbf{375} \\ \\
% <<<END-DATA lit_domains
\bottomrule
\end{tabular}
\caption{The corpus by domain.}
\label{tab:litdomains}
\end{table}

\paragraph{A caution.}
The corpus is a catalogue, not a curated bibliography: it includes preprints
and self-published items alongside refereed work. Two entries claim a complete
proof of the conjecture in a venue with no evidence of refereeing
(Halemane, 2025; Novgorodtsev, 2026); neither is formalised on ccchallenge,
and we list them for completeness without endorsement. A third, Bruckman
(2008), appeared in a refereed venue and its argument is generally regarded
as flawed. Nothing in this paper depends on any of them. Readers wanting a
curated entry point should use Lagarias' annotated
bibliographies~\cite{lagbib1,lagbib2,lagbib3}.

\subsection{The generalised-map class}

These are the $49$ papers our classifier assigns to generalised maps --- the
part of the corpus relevant to the multiverse. Those we engage with
individually are cited in Section~\ref{sec:priorgen}.

\begin{center}\footnotesize
\begin{longtable}{@{}R{0.26\textwidth}R{0.68\textwidth}@{}}
\toprule
author (year) & title \\
\midrule
\endhead
% >>>BEGIN-DATA lit_generalised
\textit{Heppner} (1978) & Eine Bemerkung zum Hasse-Syracuse Algorithmus \\
\textit{Steiner} (1981) & On the ``$Qx+1$'' Problem, $Q$ odd \\
\textit{Steiner} (1981) & On the ``$Qx+1$'' Problem, $Q$ odd II \\
\textit{Matthews  \& Watts} (1984) & A generalization of Hasse's generalization of the Syracuse algorithm \\
\textit{Lagarias} (1985) & The $3x  +  1$ problem and its generalizations \\
\textit{Matthews  \& Watts} (1985) & A Markov approach to the generalized Syracuse algorithm \\
\textit{Leigh} (1986) & A Markov process underlying the generalized Syracuse algorithm \\
\textit{Dolan et al.} (1987) & A generalization of Everett's result on the Collatz $3x  +  1$ problem \\
\textit{Matthews  \& Leigh} (1987) & A generalization of the Syracuse algorithm to $F_q [x]$ \\
\textit{Buttsworth  \& Matthews} (1990) & On some Markov matrices arising from the generalized Collatz mapping \\
\textit{Franco} (1990) & Diophantine Approximation and the $qx  +  1$ Problem \\
\textit{Zhang  \& Yu} (1990) & A generalization of Kakutani's conjecture \\
\textit{Venturini} (1992) & Iterates of Number Theoretic Functions with Periodic Rational Coefficients (Generalization of the $3x + 1$ problem) \\
\textit{Korec} (1992) & The $3x  +  1$ Problem, Generalized Pascal Triangles, and Cellular Automata \\
\textit{Matthews} (1992) & Some Borel measures associated with the generalized Collatz mapping \\
\textit{Wirsching} (1994) & A Markov chain underlying the backward Syracuse algorithm \\
\textit{Mignosi} (1995) & On a Generalization of the $3x  +  1$ Problem \\
\textit{Brocco} (1995) & A Note on Mignosi's Generalization of the $3x+1$ Problem \\
\textit{Franco  \& Pomerance} (1995) & On a Conjecture of Crandall Concerning the $qx  +  1$ Problem \\
\textit{Bernstein  \& Lagarias} (1996) & The $3x+1$ Conjugacy Map \\
\textit{Venturini} (1997) & On a Generalization of the $3x + 1$ problem \\
\textit{Belaga  \& Mignotte} (1999) & Embedding the $3x+1$ Conjecture in a $3x+d$ Context \\
\textit{Garcia  \& Tal} (1999) & A note on the generalized $3n + 1$ problem \\
\textit{Belaga  \& Mignotte} (2000) & Cyclic Structure of Dynamical Systems Associated with $3x+d$ Extensions of Collatz Problem \\
\textit{Dumont  \& Reiter} (2001) & Visualizing Generalized 3x+1 Function Dynamics \\
\textit{Ke  \& Wei} (2001) & Generalization Concerning the $3x+1$ Problem \\
\textit{Zarnowski  \& E.} (2001) & Generalized inverses and the total stopping time of Collatz sequences \\
\textit{Rouet  \& Feix} (2002) & A generalization of the Collatz problem. Building cycles and a stochastic approach \\
\textit{Belaga} (2003) & Effective polynomial upper bounds to perigees and numbers of $(3x+d)$-cycles of a given oddlength \\
\textit{Wang et al.} (2003) & The distribution of the fixed points on the real axis of a generalized $3x+1$ function and some related fractal images \\
\textit{Michel  \& Pascal} (2004) & Small Turing machines and generalized busy beaver competition \\
\textit{Ohira  \& Yamashita} (2004) & A Generalization of the Collatz problem \\
\textit{Farkas} (2005) & Variants of the $3N+1$ problem and multiplicative semigroups \\
\textit{Matthews  \& R.} (2005) & The generalized $3x+1$ mapping \\
\textit{Belaga  \& Mignotte} (2006) & The Collatz problem and its generalizations: Experimental Data. Table 1. Primitive cycles of $3x+d$ mappings \\
\textit{Belaga  \& Mignotte} (2006) & The Collatz problem and its generalizations: Experimental Data. Table 2. Factorization of Collatz numbers $2^l-3^k$ \\
\textit{Simons} (2007) & A simple (inductive) proof for  the non-existence of $2$-cycles for the $3x+1$ problem \\
\textit{Kurtz  \& Simon} (2007) & The undecidability of the generalized Collatz problem \\
\textit{Wang  \& Yu} (2007) & Visualizing generalized $3x+1$ function dynamics based on fractal \\
\textit{Snapp  \& Tracy} (2008) & The Collatz problem and Analogues \\
\textit{Simons} (2008) & On the (non)-existence of $m$-cycles for generalized Syracuse sequences \\
\textit{Hicks et al.} (2008) & A Polynomial Analogue of the $3N+1$ Problem? \\
\textit{Wang et al.} (2008) & Dynamics of the generalized $3x+1$ function determined by its fractal images \\
\textit{Muller  \& Helmut} (2009) & Uber Periodenlangen and die Vermutungen von Collatz und Crandall \\
\textit{Capco} (2015) & Odd Collatz Sequence and Binary Representations \\
\textit{Santos} (2020) & On the Collatz general problem qn+1 \\
\textit{Yolcu et al.} (2021) & An automated approach to the Collatz conjecture. \\
\textit{Siegel} (2024) & {p, q}-adic Analysis and the Collatz Conjecture \\
\textit{Novgorodtsev} (2026) & Algebraic Proof of the Collatz Conjecture and Universal Classification of 2-Adic Dynamical Systems \\ \\
% <<<END-DATA lit_generalised
\bottomrule
\end{longtable}
\end{center}

\subsection{The full snapshot}

For completeness, the whole corpus, tagged by domain:
\textbf{G} generalised maps,
\textbf{C} cycles and lower bounds,
\textbf{D} stopping time and density,
\textbf{V} computation and verification,
\textbf{L} logic, automata, formalisation,
\textbf{P} $p$-adic, ergodic, analytic,
\textbf{S} surveys,
\textbf{--} other $3x+1$.

\begin{center}\scriptsize
\begin{longtable}{@{}R{0.02\textwidth}R{0.17\textwidth}R{0.05\textwidth}R{0.66\textwidth}@{}}
\toprule
 & author & year & title \\
\midrule
\endhead
% >>>BEGIN-DATA lit_corpus
- & Slakmon  \& Macot & 2006 & On the almost convergence of Syracuse sequences \\
G & Novgorodtsev & 2026 & Algebraic Proof of the Collatz Conjecture and Universal Classifica ... \\
D & Novgorodtsev & 2026 & Exact Synchronized Coalescence Clocks and Peak-Envelope Transfer f ... \\
- & Kontorovich  \& Miller & 2005 & Benford's law, values of $L$-functions, and the $3x+1$ problem \\
- & Kontorovich  \& Sinai & 2002 & Structure Theorem for $(d, g, h)$-maps \\
V & Stokes & 2021 & The search for self-contained numbers: k-special 3-smooth represen ... \\
C & Amleh et al. & 1998 & On some difference equations with eventually periodic solutions \\
- & Atkin & 1966 & Comment on Problem $63-13^*$ \\
- & Baker  \& A. & 1966 & Linear forms in the logarithms of algebraic numbers \\
- & Brent & 2002 & $3X+1$ dynamics on rationals with fixed denominator \\
G & Snapp  \& Tracy & 2008 & The Collatz problem and Analogues \\
C & Seifert & 1988 & On the arithmetic of cycles for the Collatz-Hasse (`Syracuse') con ... \\
P & Kraft  \& Monks & 2010 & On conjugacies for the $3x+1$ map induced by continuous endomorphi ... \\
- & Wiggin & 1988 & Wondrous Numbers -- Conjecture about the $3n  +  1$ family \\
- & Hong & 1986 & About $3X+1$ problem \\
V & Hayes & 1984 & Computer recreations: The ups and downs of hailstone numbers \\
- & Thwaites & 1985 & My conjecture \\
- & Thwaites & 1996 & Two Conjectures, or how to win $pounds$ 1000 \\
C & Chisala & 1994 & Cycles in Collatz Sequences \\
- & Everett & 1977 & Iteration of the number theoretic function $f(2n)=n, f(2n+1)=3n+2$ \\
S & Ogilvy & 1972 & Tomorrow's Math: unsolved problems for the amateur \\
- & Viola & 1983 & Un Problema di Aritmetica (A problem of arithmetic) \\
- & Inclan & 2026 & Analysis of Collatz Trajectories via Modulo 6 Modular Families and ... \\
- & Inclan & 2026 & Inverse Collatz map study, version 3 \\
C & Farruggia et al. & 1996 & The elimination of a family of periodic parity vectors in the $3x+$ ... \\
V & Cadogan & 1996 & Exploring the $3x+1$ problem I \\
- & Ashbacher & 1992 & Further Investigations of the Wondrous Numbers \\
- & Cadogan & 1984 & A note on the $3x +  1$ problem \\
V & Cadogan & 1991 & Some observations on the $3x+1$ problem \\
V & Cadogan & 2000 & The $3x+1$ problem: towards a solution \\
V & Cadogan & 2003 & Trajectories in the $3x+1$ problem \\
- & Cadogan & 2006 & A Solution to the $3x+1$ Problem \\
- & Trigg et al. & 1976 & Comments on Problem 133 \\
C & Hercher & 2023 & There are no Collatz $m$-Cycles with $m leq91$ \\
- & Baiocchi & 1998 & 3N+1, UTM e Tag-Systems \\
- & Pickover & 1989 & Hailstone $3n  +  1$ Number graphs \\
C & Bohm  \& Sontacchi & 1978 & On the existence of cycles of given length in integer sequences li ... \\
- & Levy & 2004 & Injectivity and Surjectivity of Collatz Functions \\
- & Rawsthorne & 1985 & Imitation of an iteration \\
P & Bernstein & 1994 & A Non-Iterative 2-adic Statement of the $3x +  1$ Conjecture \\
G & Bernstein  \& Lagarias & 1996 & The $3x+1$ Conjugacy Map \\
- & Shanks & 1965 & Comments on Problem $63-13^*$ \\
- & Shanks & 1975 & Problem $\#5$ \\
- & Merlini  \& Sala & 1999 & On the Fibonacci's Attractor and the Long Orbits in the $3n+1$ Pro ... \\
D & Klarner & 1981 & An algorithm to determine when certain sets have $0$ density \\
- & Klarner & 1982 & A sufficient condition for certain semigroups to be free \\
- & Klarner & 1988 & $m$-recognizability of sets closed under certain affine functions \\
- & Klarner  \& Post & 1992 & Some fascinating integer sequences \\
- & Klarner  \& Rado & 1973 & Linear combinations of sets of consecutive integers \\
- & Klarner  \& Rado & 1974 & Arithmetic properties of certain recursively defined sets \\
D & Applegate  \& Lagarias & 1995 & Density Bounds for the $3x+1$ Problem I. Tree-Search Method \\
D & Applegate  \& Lagarias & 1995 & Density Bounds for the $3x+1$ Problem II. Krasikov Inequalities \\
V & Applegate  \& Lagarias & 1995 & The distribution of $3x+1$ trees \\
D & Applegate  \& Lagarias & 2003 & Lower bounds for the total stopping time  of $3x+1$ iterates \\
- & Applegate  \& Lagarias & 2006 & The $3x+1$ semigroup \\
V & Barina & 2025 & Improved verification limit for the convergence of the Collatz con ... \\
- & Gale & 1991 & Mathematical Entertainments: More Mysteries \\
- & Gluck  \& Taylor & 2002 & A new statistic for the $3x+1$ problem \\
V & Kay & 1972 & An Algorithm for Reducing the Size of an Integer \\
- & Boyd & 1985 & Which rationals are ratios of Pisot sequences? \\
- & Clark & 1995 & Second-Order Difference Equations Related to the Collatz $3n+1$ Co ... \\
C & Clark & 2006 & Periodic solutions of arbitrary length in a simple integer iteration \\
V & Clark  \& Lewis & 1995 & A Collatz-Type Difference Equation \\
V & Clark  \& Lewis & 1998 & Symmetric solutions to a Collatz-like system of  Difference Equati ... \\
- & Hoffman  \& Klarner & 1978 & Sets of integers closed under affine operators- the closure of fin ... \\
- & Hoffman  \& Klarner & 1979 & Sets of integers closed under affine operators- the finite basis t ... \\
- & Domenici & 2009 & A few observations on the Collatz problem \\
- & Coppersmith & 1975 & The complement of certain recursively defined sets \\
- & Shaw & 2006 & The pure numbers generated by the Collatz sequence \\
- & Marcu & 1991 & The powers of two and three \\
- & Belaga & 1998 & Reflecting on the $3x+1$ Mystery: Outline of a Scenario- Improbabl ... \\
G & Belaga & 2003 & Effective polynomial upper bounds to perigees and numbers of $(3x+$ ... \\
G & Belaga  \& Mignotte & 1999 & Embedding the $3x+1$ Conjecture in a $3x+d$ Context \\
G & Belaga  \& Mignotte & 2000 & Cyclic Structure of Dynamical Systems Associated with $3x+d$ Exten ... \\
V & Belaga  \& Mignotte & 2006 & Walking Cautiously into the Collatz Wilderness: Algorithmically, N ... \\
G & Belaga  \& Mignotte & 2006 & The Collatz problem and its generalizations: Experimental Data. Ta ... \\
G & Belaga  \& Mignotte & 2006 & The Collatz problem and its generalizations: Experimental Data. Ta ... \\
L & Lehtonen & 2008 & Two undecidable versions of Collatz's problems \\
- & Eliahou et al. & 2022 & Is the Syracuse Falling Time Bounded by 12? \\
P & Jesus & 2026 & A Global Occupation Conjecture for the Accelerated 3x + 1 Map with ... \\
- & Oliveri  \& Vella & 1998 & Alcune Questioni Correlate al Problema Del ``3D+1'' \\
G & Yolcu et al. & 2021 & An automated approach to the Collatz conjecture. \\
- & Barone & 1999 & A heuristic probabilistic argument for the Collatz sequence \\
- & Roosendaal & 2004 & On the $3x+1$ problem \\
G & Heppner & 1978 & Eine Bemerkung zum Hasse-Syracuse Algorithmus \\
P & Akin & 2004 & Why is the $3x+1$ Problem Hard? \\
G & Mignosi & 1995 & On a Generalization of the $3x  +  1$ Problem \\
V & Kascak & 1992 & Small universal one-state linear operator algorithms \\
- & Jarvis & 1989 & 13, 31 and the $3x  +  1$ problem \\
V & Leavens  \& Vermeulen & 1992 & $3x +1$ Search Programs \\
V & Gonnet & 1991 & Computations on the $3n+1$ Conjecture \\
G & Leigh & 1986 & A Markov process underlying the generalized Syracuse algorithm \\
- & Rhin & 1987 & Approximants de Pade et mesures effectives d'irrationalite \\
- & Rhin  \& Toffin & 1986 & Approximants de Pade simultanes de logarithmes \\
- & Venturini & 1982 & Sul Comportamento delle Iterazioni di Alcune Funzioni Numeriche \\
- & Venturini & 1989 & On the $3x + 1$ Problem \\
G & Venturini & 1992 & Iterates of Number Theoretic Functions with Periodic Rational Coef ... \\
G & Venturini & 1997 & On a Generalization of the $3x + 1$ problem \\
- & Hasenjager & 1990 & Hasse's Syracuse-Problem und die Rolle der Basen \\
L & Scollo & 2005 & $omega$-rewriting the Collatz problem \\
- & Gao & 1993 & On consecutive numbers of the same height in the Collatz problem \\
- & Targonski & 1991 & Open questions about KW-orbits and iterative roots \\
- & Wirsching & 1993 & An Improved Estimate Concerning $3N  +  1$ Predecessor Sets \\
G & Wirsching & 1994 & A Markov chain underlying the backward Syracuse algorithm \\
- & Wirsching & 1996 & On the Combinatorial Structure of $3N+1$ Predecessor Sets \\
P & Wirsching & 1996 & $3n+1$ Predecessor Densities and Uniform Distribution in $mathbb$ ... \\
P & Wirsching & 1998 & The Dynamical System Generated by the $3n + 1$ Function \\
- & Wirsching & 1998 & Balls in Constrained Urns and Cantor-Like Sets \\
- & Wirsching & 2000 & Uber das $3n + 1$ Problem \\
V & Glaser  \& Weigand & 1989 & Das-ULAM Problem-Computergestutze Entdeckungen \\
- & Riele & 1983 & Problem 669 \\
- & Riele & 1983 & Iteration of Number-theoretic functions \\
- & Coxeter & 1971 & Cyclic Sequences and Frieze Patterns, (The Fourth Felix Behrend Me ... \\
- & He  \& Wang & 2003 & $3n+1$ problem's simplifying and structure property of $T(n)$ ... \\
- & Lunkenheimer & 1988 & Eine kleine Untersuchung zu einem zahlentheoretischen Problem \\
- & Muller & 1991 & Das `$3n+1$' Problem \\
- & Muller & 1994 & Uber eine Klasse 2-adischer Funktionen im Zusammenhang mit dem `` ... \\
S & Hasse & 1975 & Unsolved Problems in Elementary Number Theory \\
V & Moller & 1978 & Uber Hasses Verallgemeinerung des Syracuse-Algorithmus (Kakutanis ... \\
G & Farkas & 2005 & Variants of the $3N+1$ problem and multiplicative semigroups \\
- & Huang et al. & 2000 & One-to-one correspondence between the natural numbers and the pari ... \\
- & Krasikov & 1989 & How many numbers satisfy the $3x  +  1$ Conjecture? \\
V & Krasikov  \& Lagarias & 2003 & Bounds for the $3x+1$ problem using difference inequalities \\
- & Kerner & 2000 & Die Collatz-Ulam-Kombination CUK \\
- & Korec  \& Znam & 1987 & A Note on the $3x  +  1$ Problem \\
G & Korec & 1992 & The $3x  +  1$ Problem, Generalized Pascal Triangles, and Cellular ... \\
D & Korec & 1994 & A Density Estimate for the $3x+1$ Problem \\
V & Davison & 1977 & Some Comments on an Iteration Problem \\
- & azewitz  \& Pettorossi & 1983 & Some properties of binary sequences useful for proving Collatz's c ... \\
V & Arsac & 1986 & Algorithmes pour verifier la conjecture de Syracuse \\
- & Roubaud & 1969 & Un probleme combinatoire pose par la poesie lyrique des troubadours \\
- & Roubaud & 1993 & $N$-ine, autrement dit quenine (encore), Reflexions historiques et ... \\
G & Dolan et al. & 1987 & A generalization of Everett's result on the Collatz $3x  +  1$ pro ... \\
- & Kuttler & 1994 & On the $3x+1$ Problem \\
L & Kari  \& Kopra & 2017 & Cellular automata and powers of p/q \\
- & Bendegem & 2005 & The Collatz Conjecture: A Case Study in Mathematical Problem Solving \\
- & Dumas & 2008 & Caracterisation des quenines et leur representation spirale \\
- & Daudin & 2020 & $3n+1$ problem: an heuristic lower bound for the number of integer ... \\
G & Rouet  \& Feix & 2002 & A generalization of the Collatz problem. Building cycles and a sto ... \\
- & Allouche & 1979 & Sur la conjecture de ``Syracuse-Kakutani-Collatz'' \\
C & Lagarias et al. & 1994 & Asymmetric Tent Map Expansions II. Purely Periodic Points \\
V & Shallit  \& Wilson & 1991 & The ``$3x+1$'' Problem and Finite Automata \\
G & Lagarias & 1985 & The $3x  +  1$ problem and its generalizations \\
C & Lagarias & 1990 & The set of rational cycles for the $3x  +  1$ problem \\
- & Lagarias & 1999 & How random are 3x + 1 function iterates? \\
S & Lagarias & 2003 & The $3x+1$ Problem: An Annotated Bibliography (1963--1999) \\
- & Lagarias & 2006 & Wild and Wooley Numbers \\
S & Lagarias & 2010 & The Ultimate Challenge: The $3x+1$ Problem \\
S & Lagarias & 2011 & The 3x+1 problem: An annotated bibliography (1963--1999) (sorted b ... \\
S & Lagarias & 2012 & The 3x+1 Problem: An Annotated Bibliography, II (2000-2009) \\
S & Lagarias & 2021 & The $3x+1$ Problem: An Overview \\
- & Lagarias  \& Weiss & 1992 & The $3x  +  1$ Problem: Two Stochastic Models \\
- & Lagarias  \& Soundararajan & 2006 & Benford's Law for the $3x+1$ Function \\
V & Lagarias  \& Sloane & 2004 & Approximate squaring \\
C & Lagarias et al. & 1993 & Asymmetric Tent Map Expansions I. Eventually Periodic Points \\
G & Dumont  \& Reiter & 2001 & Visualizing Generalized 3x+1 Function Dynamics \\
- & Dumont  \& Reiter & 2003 & Real dynamics of a $3$-power extension of the $3x+1$ function \\
- & Goodwin & 2003 & Results on the Collatz conjecture \\
V & Marcinkowski & 1999 & Achilles, Turtle, and Undecidable Boundedness Problems for Small D ... \\
- & Joseph & 1998 & A chaotic extension of the Collatz function to $Z_2 [i]$ \\
- & Conway & 1972 & Unpredicatable Iterations \\
V & Conway & 1987 & FRACTRAN- A Simple Universal Computing Language for Arithmetic \\
C & Simons & 2005 & On the non-existence of $2$-cycles for the $3x+1$ problem \\
G & Simons & 2007 & A simple (inductive) proof for  the non-existence of $2$-cycles fo ... \\
G & Simons & 2008 & On the (non)-existence of $m$-cycles for generalized Syracuse sequ ... \\
C & Simons & 2008 & Exotic Collatz Cycles \\
C & Simons & 2008 & Post-transcendence conditions for the existence of $m$-cycles for ... \\
- & Simons & 2008 & On isomorphism between Farkas sequences and Collatz sequences \\
- & Simons  \& Weger & 2004 & Mersenne en het Syracuseprobleem [Mersenne and the Syracuse problem] \\
C & Simons  \& Weger & 2005 & Theoretical and computational bounds for $m$-cycles of the $3n+1$ ... \\
G & Capco & 2015 & Odd Collatz Sequence and Binary Representations \\
P & Lopez  \& Stoll & 2012 & The 2-Adic, Binary and Decimal Periods of 1/3k Approach Full Compl ... \\
V & Pe & 2004 & The $3x+1$ Fractal \\
- & Alves et al. & 2005 & A linear algebra approach to the conjecture of Collatz \\
V & Nivergelt & 1975 & Computers and Mathematics Education \\
- & Sander & 1990 & On the $(3N  +  1)$-conjecture \\
L & Kari  \& Jarkko & 2012 & Cellular Automata, the Collatz Conjecture and Powers of 3/2 \\
V & Metzger & 1995 & Untersuchungen zum $(3n+1)$-Algorithmus. Teil I: Das Problem der Z ... \\
V & Metzger & 1999 & Zyklenbestimmung beim $(bn+1)$-Algorithmus \\
G & Ke  \& Wei & 2001 & Generalization Concerning the $3x+1$ Problem \\
- & Ke  \& Yong-Sheng & 2005 & The proof of the hypothesis about $"3x+1"$ \\
G & Matthews & 1992 & Some Borel measures associated with the generalized Collatz mapping \\
G & Matthews  \& Watts & 1984 & A generalization of Hasse's generalization of the Syracuse algorithm \\
G & Matthews  \& Watts & 1985 & A Markov approach to the generalized Syracuse algorithm \\
G & Matthews  \& Leigh & 1987 & A generalization of the Syracuse algorithm to $F_q [x]$ \\
G & Hicks et al. & 2008 & A Polynomial Analogue of the $3N+1$ Problem? \\
- & Stolarsky & 1998 & A prelude to the $3x+1$ problem \\
C & Halemane & 2025-12-10 & Collatz-Conjecture Proved and Halemane-Conjecture Proposed \\
C & Knight & 2026 & Collatz high cycles do not exist \\
- & Knight & 2026 & A Small Collatz Rule without the Plus One \\
- & Borovkov  \& Pfeifer & 2000 & Estimates for the Syracuse problem via a probabilistic model \\
- & Mahler & 1968 & An unsolved problem on the powers of $3/2$ \\
- & Michel et al. & 1995 & Formes lineaires en deux logarithmes et determinants d'interpolation \\
V & Ratzan & 1973 & Some work on an Unsolved Palindromic Algorithm \\
- & Flatto & 1992 & $Z$-numbers and $beta$-transformations \\
- & Li  \& Chun & 2002 & Compressive Iteration for the $3N+1$ conjecture \\
- & Li  \& Chun & 2003 & Contractible iteration for the $3n+1$ conjecture \\
- & Li  \& Chun & 2004 & Same-flowing numbers and super contraction iteration in $3N+1$ con ... \\
C & Li  \& Chun & 2005 & Necessary condition for periodic numbers in the $3N+1$ conjecture \\
- & Li  \& Chun & 2006 & Some properties of super contraction iteration in the $3N+1$ conje ... \\
V & Li et al. & 2002 & Research on the Syracuse Operator and its properties \\
- & Li et al. & 2006 & Equivalence of the $3N+1$ and $3N+3k$ conjecture and some related ... \\
C & Li et al. & 2004 & Study of periodic numbers in the $3N+1$ conjecture \\
V & Mol & 2008 & Tag Systems and Collatz-like functions \\
C & Halbeisen  \& Hungerbuhler & 1997 & Optimal bounds for the length of rational Collatz cycles \\
- & Berg  \& Meinardus & 1994 & Functional equations connected with the Collatz problem \\
- & Berg  \& Meinardus & 1995 & The 3n+1 Collatz Problem and Functional Equations \\
- & Collatz & 1986 & On the Motivation and Origin of the $(3n +  1)$-Problem \\
- & Kauffman & 1995 & Arithmetic in the form \\
C & Luca  \& Florian & 2005 & On the nontrivial cycles in Collatz's problem \\
V & Garner & 1981 & On the Collatz $3n+1$ algorithm \\
V & Garner & 1985 & On heights in the Collatz $3n  +  1$ problem \\
- & Lopez et al. & 2009 & The $3x+1$ conjugacy map over a Sturmian word \\
C & Macindoe  \& James & 2026 & Reduced coordinates for the Collatz map: exact per-step laws, anch ... \\
G & Garcia  \& Tal & 1999 & A note on the generalized $3n + 1$ problem \\
P & Chamberland & 1996 & A Continuous Extension of the $3x+1$ Problem to the Real Line \\
- & Feix et al. & 1995 & The $(3x +  1)/2$ Problem: A Statistical Approach \\
- & Feix et al. & 1994 & Statistical Properties of an Iterated Arithmetic Mapping \\
V & Margenstern  \& Maurice & 2000 & Frontier between decidability and undecidability: a survey \\
- & Pierantoni  \& Curcic & 1996 & A transformation of iterated Collatz mappings \\
- & Chamberland & 2003 & Una actualizachio del problema $3x+1$ \\
- & Gardner & 1972 & Mathematical Games \\
- & Yamada & 1980 & A convergence proof about an integral sequence \\
V & Cook et al. & 2021 & Small Tile Sets That Compute While Solving Mazes \\
G & Matthews  \& R. & 2005 & The generalized $3x+1$ mapping \\
C & Sinisalo & 2003 & On the minimal cycle lengths of the Collatz sequences \\
- & Margenstern  \& Matiyasevich & 1999 & A binomial representation of the $3X +1$ problem \\
G & Siegel & 2024 & {p, q}-adic Analysis and the Collatz Conjecture \\
V & Metzger  \& Heinz & 2000 & Untersuchungen zum $(3n+1)$-Algorithmus. Teil II: Die Konstruktion ... \\
V & Metzger  \& Heinz & 2003 & Untersuchungen zum $(3n+1)$-Algorithmus, Teil III: Gesetzmassigkei ... \\
- & Beeler et al. & 1972 & HAKMEM \\
- & Fredman & 1972 & Growth properties of a class of recursively defined functions \\
- & Fredman  \& Knuth & 1974 & Recurrence relations based on minimization \\
- & Avidon & 1997 & On primitive $3$-smooth partitions of $n$ \\
- & Williams & 2024 & Collatz conjecture: an order isomorphic recursive machine \\
G & Michel  \& Pascal & 2004 & Small Turing machines and generalized busy beaver competition \\
C & Mimuro  \& Tomoaki & 2001 & On certain simple cycles of the Collatz conjecture \\
- & Misiurewicz et al. & 2005 & Real $3X+1$ \\
- & Bringer & 1969 & Sur un probleme de R. Queneau \\
V & Monks  \& G. & 2002 & $3X+1$ Minus the $+$ \\
- & Monks et al. & 2004 & The Autoconjugacy of the $3x+1$ function \\
C & Monks  \& M. & 2006 & The sufficiency of arithmetic progressions for the $3x+1$ conjecture \\
P & Monks  \& Maria & 2009 & Endomorphisms of the shift dynamical systems, discrete derivatives ... \\
- & Klamkin & 1963 & Problem $63-13^*$ \\
V & Chacon  \& Vielhaber & 2004 & Structural and Computational Complexity of Isometries and their Sh ... \\
G & Muller  \& Helmut & 2009 & Uber Periodenlangen and die Vermutungen von Collatz und Crandall \\
P & Pippenger & 1993 & An elementary approach to some analytic asymptotics \\
C & Rozier & 1990 & Demonstraton de l'absence de cycles d'une certain forme pour le pr ... \\
C & Rozier & 1991 & Probleme de Syracuse: ``Majorations'' Elementaires des Cycles \\
- & Rozier & 2017 & The 3x+1 problem: a lower bound hypothesis \\
P & Rozier & 2018 & Parity Sequences of the $3x+1$ Map on the 2-adic Integers and Eucl ... \\
C & Rozier & 2022 & Upper bound on the elements of a Collatz cycle \\
- & Rozier  \& Terracol & 2025 & Paradoxical behavior in Collatz sequences \\
- & Pan  \& Hong-liang & 2000 & Notes on the $3x+1$ problem \\
V & Park  \& Jungwon & 2026 & A Machine-Verified Conditional Theory of the Collatz Conjecture. \\
L & Michel & 1993 & Busy Beaver Competition and Collatz-like Problems \\
L & Michel & 2014 & Simulation of the Collatz 3x+1 function by Turing machines \\
D & Andaloro & 2000 & On total stopping times under $3X+1$ iteration \\
- & Andaloro & 2002 & The $3X+1$ problem and directed graphs \\
- & Erdos  \& Graham & 1979 & Old and new problems and results in combinatorial number theory: v ... \\
- & Erdos  \& Graham & 1980 & Old and new problems and results in combinatorial number theory \\
- & Bruckman & 2008 & A proof of the Collatz conjecture \\
- & Konstadinidis & 2006 & The real $3x+1$ problem \\
- & Elliott & 1985 & Arithmetic Functions and Integer Products \\
P & Eisele  \& Hadeler & 1990 & Game of Cards, Dynamical Systems, and a Characterization of the Fl ... \\
C & Kosobutskyy  \& Rebot & 2023 & Collatz Conjecture 3n  1 as a Newton Binomial Problem \\
L & Devienne et al. & 1993 & Halting Problem of One Binary Horn Clause is Undecidable \\
- & Filipponi & 1991 & On the $3n+1$ Problem: Something Old, Something New \\
- & Qu  \& Jing-hua & 2002 & The $3x+1$ problem and its necessary and sufficient condition \\
- & Banerji & 1996 & Some Properties of the $3n+1$ Function \\
V & Steiner & 1978 & A Theorem on the Syracuse Problem \\
G & Steiner & 1981 & On the ``$Qx+1$'' Problem, $Q$ odd \\
G & Steiner & 1981 & On the ``$Qx+1$'' Problem, $Q$ odd II \\
- & Queneau & 1963 & Note complementaire sur la Sextine \\
- & Queneau & 1972 & Sur les suites $s$-additives \\
G & Ohira  \& Yamashita & 2004 & A Generalization of the Collatz problem \\
V & Ohira  \& Yamashita & 2004 & On the $p$-Collatz problem \\
- & Blecksmith et al. & 1998 & 3-Smooth Representations of Integers \\
V & Dunn & 1973 & On Ulam's Problem \\
- & Crandall & 1978 & On the ``$3x + 1$'' problem \\
S & Guy & 1981 & Unsolved Problems in Number Theory \\
- & Guy & 1983 & Don't try to solve these problems! \\
- & Guy & 1983 & Conway's prime producing machine \\
- & Guy & 1986 & John Isbell's Game of Beanstalk and John Conway's Game of Beans Do ... \\
D & Terras & 1976 & A stopping time problem on the positive integers \\
D & Terras & 1979 & On the existence of a density \\
G & Buttsworth  \& Matthews & 1990 & On some Markov matrices arising from the generalized Collatz mapping \\
G & Santos & 2020 & On the Collatz general problem qn+1 \\
- & Tijdeman & 1972 & Note on Mahler's $3/2$-problem \\
- & Stemmler & 1964 & The ideal Waring problem for exponents $401-200,000$ \\
- & Anderson & 1987 & Struggling with the $3x  +  1$ problem \\
V & Burckel & 1994 & Functional equations associated with congruential functions \\
- & Albeverio et al. & 1989 & Una breve introduzione a diffusioni su insiemi frattali e ad alcun ... \\
C & Eliahou & 1993 & The $3x +  1$ problem: New Lower Bounds on Nontrivial Cycle Lengths \\
- & Eliahou  \& Verger-Gaugry & 2025 & The number system in rational base $3/2$ and the $3x+1$ problem \\
V & Wang & 1988 & Some researches for transformation over recursive programs \\
C & Shi  \& Li & 2005 & Properties of cycle sets \\
- & Shi  \& Li & 2005 & Properties of the $3x+1$ problem \\
V & Letherman et al. & 1999 & On the $3X+1$ problem and holomorphic dynamics \\
- & Volkov & 2006 & A probabilistic model for the $5k+1$ problem and related problems \\
- & Wagon & 1985 & The Collatz problem \\
- & Andrei  \& Masalagiu & 1998 & About the Collatz Conjecture \\
- & Andrei et al. & 2000 & Some results on the Collatz problem \\
- & Andrei et al. & 2007 & A Functional View over the Collatz's Problem \\
- & Kohl & 2005 & Restklassenweise affine Gruppen \\
D & Kohl & 2007 & Wildness of iteration of certain residue class-wise affine mappings \\
- & Kohl & 2008 & On conjugates of Collatz-type mappings \\
V & Kohl & 2008 & Algorithms for a class of infinite permutation groups \\
D & Kohl & 2010 & A simple group generated by involutions interchanging residue clas ... \\
- & Beltraminelli et al. & 1994 & Orbite inverse nel problema del $3n+1$ \\
G & Brocco & 1995 & A Note on Mignosi's Generalization of the $3x+1$ Problem \\
- & Isard  \& Zwicky & 1970 & Three open questions in the theory of one-symbol Smullyan systems \\
G & Kurtz  \& Simon & 2007 & The undecidability of the generalized Collatz problem \\
- & Sterin  \& Tristan & 2020 & Binary expression of ancestors in the Collatz graph \\
V & Sterin et al. & 2020 & The Collatz Process Embeds a Base Conversion Algorithm \\
- & Puddu & 1986 & The Syracuse problem (Spanish), 5$^rmth$ \\
- & Tang  \& Xian & 2006 & About recursion relation of $3x+1$ Conjecture (Chinese) \\
V & Urvoy & 2000 & Regularity of congruential graphs \\
D & Tao  \& Terence & 2022 & Almost all orbits of the Collatz map attain almost bounded values \\
- & Amdeberhan et al. & 2008 & The $2$-adic valuation of a sequence arising from a rational integ ... \\
C & Brox & 2000 & Collatz cycles with few descents \\
L & Cloney et al. & 1987 & The $3x  +  1$ Problem: a Quasi-Cellular Automaton \\
C & Keller & 1999 & Finite cycles of certain periodically linear functions \\
V & Silva & 2004 & Computational verification of $3x+1$ conjecture \\
D & Silva & 1999 & Maximum Excursion and Stopping Time Record-Holders for the $3x+1$ ... \\
- & Urata & 1999 & Collatz Problem IV \\
- & Urata & 2002 & Some holomorphic functions connected with the Collatz problem \\
- & Urata & 2003 & The Collatz problem over $2$-adic integers \\
- & Urata & 2007 & Collatz's problem \\
- & Urata  \& Kajita & 1998 & Collatz Problem III \\
- & Urata  \& Suzuki & 1996 & Collatz Problem \\
- & Urata et al. & 1997 & Collatz Problem II \\
- & Urata  \& Hamada & 2005 & Positive values of a holomorphic function connected with the Colla ... \\
V & Sterin & 2023 & Six Tiles : From Collatz Sequences to Algorithmic DNA Origami \\
V & Chvatal et al. & 1972 & Selected combinatorial research problems \\
P & Bergelson et al. & 2006 & Affine Actions of a Free Semigroup on the Real Line \\
- & Ellison & 1971 & On a Theorem of S. Sivasankaranarayana Pillai \\
- & Wang  \& Nian-liang & 2002 & The way to prove Collatz problem \\
G & Wang et al. & 2003 & The distribution of the fixed points on the real axis of a general ... \\
G & Wang et al. & 2008 & Dynamics of the generalized $3x+1$ function determined by its frac ... \\
- & Qiu & 1997 & Study on $``3x+1"$ problem \\
- & Hotzel & 2003 & Beitage zum $3n+1$-Problem \\
- & Wirsching  \& J. & 2001 & A functional differential equation and $3n + 1$ dynamics \\
D & Wirsching  \& J. & 2003 & On the problem of positive predecessor density in $3N+1$ dynamics \\
- & Gaschutz & 1982 & Linear abgeschlossene Zahlenmengen I. [Linearly closed number sets ... \\
- & Wu  \& Bang & 1992 & A tendency of the proportion  of consecutive numbers of 
the same ... \\
- & Wu  \& Bang & 1993 & On the consecutive positive integers of the same height in the 
Co ... \\
- & Wu  \& Bang & 1995 & The monotonicity of pairs of coalescence numbers in the Collatz
pr ... \\
D & Wu et al. & 2003 & On the equality of the adequate stopping time and the coefficient ... \\
- & Wu et al. & 2001 & On the traditional iteration and the elongate iteration of the $3N$ ... \\
- & Wu et al. & 2001 & Family of consecutive integer pairs of the same height in the Coll ... \\
- & Wu et al. & 2001 & Elongate iteration for the $3N+1$ Conjecture \\
- & Wu  \& Quan & 2003 & An equivalent proposition of Collatz conjecture \\
- & Narkiewicz & 1980 & A note on a paper of H. Gupta concerning powers of 2 and 3 \\
- & Xia  \& Ye & 2003 & Some thoughts about $3X+1$ Conjecture \\
- & Xie  \& Guo & 2006 & The interesting $3x+1$ problem \\
G & Wang  \& Yu & 2007 & Visualizing generalized $3x+1$ function dynamics based on fractal \\
- & Sinai & 2003 & Statistical $(3X+1)$-Problem \\
- & Sinai & 2003 & Uniform distribution in the $(3x+1)$ problem \\
- & Sinai & 2004 & A theorem about uniform distribution \\
- & Yamashita  \& Michinori & 2002 & $3x+1$ problem from the $(e,k)$-perspective \\
- & Yamashita  \& Michinori & 2006 & Note on the $3x pm1$ Problem \\
- & Yang  \& Hua & 1998 & An equivalent set for the $3x+1$ conjecture \\
- & Yang et al. & 1988 & Kakutani conjecture and graph theory representation of the black h ... \\
V & Ito  \& Nakano & 2009 & A Hardware-Software Cooperative Approach to the Exhaustive verific ... \\
- & Peres et al. & 2006 & Absolute continuity for random iterated function systems with over ... \\
G & Franco  \& Pomerance & 1995 & On a Conjecture of Crandall Concerning the $qx  +  1$ Problem \\
G & Franco & 1990 & Diophantine Approximation and the $qx  +  1$ Problem \\
G & Zarnowski  \& E. & 2001 & Generalized inverses and the total stopping time of Collatz sequen ... \\
- & Zarnowski  \& E. & 2008 & The congruence structure of the $3x+1$ map \\
G & Zhang  \& Yu & 1990 & A generalization of Kakutani's conjecture \\
- & Zhang et al. & 1998 & Problems on mapping sequences \\
- & Zhou  \& Zhong & 1995 & Some discussion on the $3x+1$ problem \\
- & Zhou  \& Zhong & 1997 & Some recurrence relations connected with the $3x+1$ problem \\
- & Zhou et al. & 2007 & On the Collatz Conjecture \\ \\
% <<<END-DATA lit_corpus
\bottomrule
\end{longtable}
\end{center}

% ============================================================= references ===
\FloatBarrier
\begin{thebibliography}{99}

\bibitem{carykh}
carykh.
\newblock The Collatz Multiverse.
\newblock YouTube, 22 November 2025.
\newblock \url{https://youtu.be/n63FBYqj98E}

\bibitem{lazykh}
lazykh (carykh's second channel).
\newblock Collatz Multiverse Ramble, Part 2!
\newblock YouTube, 22 November 2025.
\newblock \url{https://youtu.be/3JO-8oZ-IlQ}

\bibitem{artem}
Artem.
\newblock The Collatz Multiverse Multiverse.
\newblock YouTube, 14 December 2025.
\newblock \url{https://youtu.be/MNFWrB8iUbM}

\bibitem{crandall}
R.~E.~Crandall.
\newblock On the ``$3x+1$'' problem.
\newblock \emph{Mathematics of Computation}, 32(144):1281--1292, 1978.

\bibitem{fp}
Z.~Franco and C.~Pomerance.
\newblock On a conjecture of Crandall concerning the $qx+1$ problem.
\newblock \emph{Mathematics of Computation}, 64(211):1333--1336, 1995.

\bibitem{santos}
R.~Santos.
\newblock On the Collatz general problem $qn+1$.
\newblock arXiv:2005.00346, 2020.
\newblock \url{https://arxiv.org/abs/2005.00346}

\bibitem{bm}
E.~G.~Belaga and M.~Mignotte.
\newblock Embedding the $3x+1$ conjecture in a $3x+d$ context.
\newblock \emph{Experimental Mathematics}, 7(2):145--151, 1998.

\bibitem{matthews}
K.~R.~Matthews.
\newblock Generalized $3x+1$ mappings: Markov chains and ergodic theory.
\newblock In J.~C.~Lagarias, editor, \emph{The Ultimate Challenge: The
  $3x+1$ Problem}, pages 79--103. American Mathematical Society, 2010.

\bibitem{conway}
J.~H.~Conway.
\newblock Unpredictable iterations.
\newblock In \emph{Proceedings of the 1972 Number Theory Conference}, pages
  49--52. University of Colorado, Boulder, 1972.

\bibitem{kurtzsimon}
S.~A.~Kurtz and J.~Simon.
\newblock The undecidability of the generalized Collatz problem.
\newblock In \emph{Theory and Applications of Models of Computation (TAMC
  2007)}, volume 4484 of \emph{Lecture Notes in Computer Science}, pages
  542--553. Springer, 2007.

\bibitem{eliahou}
S.~Eliahou.
\newblock The $3x+1$ problem: new lower bounds on nontrivial cycle lengths.
\newblock \emph{Discrete Mathematics}, 118(1--3):45--56, 1993.

\bibitem{hercher}
C.~Hercher.
\newblock There are no Collatz $m$-cycles with $m \leq 91$.
\newblock \emph{Journal of Integer Sequences}, 26, Article 23.3.5, 2023.

\bibitem{barina}
D.~Barina.
\newblock Improved verification limit for the convergence of the Collatz
  conjecture.
\newblock \emph{The Journal of Supercomputing}, 2025.

\bibitem{ggm}
F.~Gon\c{c}alves, R.~Greenfeld, and J.~Madrid.
\newblock Generalized Collatz maps with almost bounded orbits.
\newblock arXiv:2111.06170, 2021.
\newblock \url{https://arxiv.org/abs/2111.06170}

\bibitem{lagbib1}
J.~C.~Lagarias.
\newblock The $3x+1$ problem: an annotated bibliography (1963--1999).
\newblock arXiv:math/0309224.

\bibitem{lagbib2}
J.~C.~Lagarias.
\newblock The $3x+1$ problem: an annotated bibliography, II (2000--2009).
\newblock arXiv:math/0608208.

\bibitem{lagbib3}
J.~C.~Lagarias.
\newblock The $3x+1$ problem: an overview.
\newblock arXiv:2111.02635.

\bibitem{cresearch}
Collatz Research.
\newblock \url{https://www.collatzresearch.org/}

\bibitem{ccc}
The Collatz Conjecture Challenge.
\newblock Formalising the Collatz literature, one paper at a time.
\newblock \url{https://ccchallenge.org/} (snapshot of 375 papers, 2026).

\bibitem{steiner1981a}
R.~P.~Steiner.
\newblock On the ``$Qx+1$'' problem, $Q$ odd.
\newblock \emph{Fibonacci Quarterly}, 19:285--288, 1981.

\bibitem{steiner1981b}
R.~P.~Steiner.
\newblock On the ``$Qx+1$'' problem, $Q$ odd II.
\newblock \emph{Fibonacci Quarterly}, 19:293--296, 1981.

\bibitem{franco1990}
Z.~M.~Franco.
\newblock Diophantine approximation and the $qx+1$ problem.
\newblock Ph.D. thesis / report, 1990.

\bibitem{belaga2003}
E.~Belaga.
\newblock Effective polynomial upper bounds to perigees and numbers of
  $(3x+d)$-cycles of a given oddlength.
\newblock \emph{Acta Arithmetica}, 106(2):197--206, 2003.

\bibitem{bm2006}
E.~Belaga and M.~Mignotte.
\newblock The Collatz problem and its generalizations: experimental data.
  Table 1: primitive cycles of $3x+d$ mappings.
\newblock Universit\'e Louis Pasteur, Strasbourg, 2006.

\bibitem{simons2008}
J.~L.~Simons.
\newblock On the (non-)existence of $m$-cycles for generalized Syracuse
  sequences.
\newblock \emph{Acta Arithmetica}, 131(3):217--254, 2008.

\bibitem{mw1984}
K.~R.~Matthews and A.~M.~Watts.
\newblock A generalization of Hasse's generalization of the Syracuse
  algorithm.
\newblock \emph{Acta Arithmetica}, 43:167--175, 1984.

\bibitem{mw1985}
K.~R.~Matthews and A.~M.~Watts.
\newblock A Markov approach to the generalized Syracuse algorithm.
\newblock \emph{Acta Arithmetica}, 45:29--42, 1985.

\bibitem{leigh1986}
G.~M.~Leigh.
\newblock A Markov process underlying the generalized Syracuse algorithm.
\newblock \emph{Acta Arithmetica}, 46:125--143, 1986.

\bibitem{bmatthews1990}
R.~N.~Buttsworth and K.~R.~Matthews.
\newblock On some Markov matrices arising from the generalized Collatz
  mapping.
\newblock \emph{Acta Arithmetica}, 55:43--57, 1990.

\bibitem{matthews2005}
K.~R.~Matthews.
\newblock The generalized $3x+1$ mapping.
\newblock \url{http://www.numbertheory.org/}, 2005.

\bibitem{venturini1992}
G.~Venturini.
\newblock Iterates of number theoretic functions with periodic rational
  coefficients (generalization of the $3x+1$ problem).
\newblock \emph{Studies in Applied Mathematics}, 86:185--218, 1992.

\bibitem{venturini1997}
G.~Venturini.
\newblock On a generalization of the $3x+1$ problem.
\newblock \emph{Advances in Applied Mathematics}, 19:295--305, 1997.

\bibitem{mignosi1995}
F.~Mignosi.
\newblock On a generalization of the $3x+1$ problem.
\newblock \emph{Journal of Number Theory}, 55:28--45, 1995.

\bibitem{brocco1995}
S.~Brocco.
\newblock A note on Mignosi's generalization of the $3x+1$ problem.
\newblock \emph{Journal of Number Theory}, 52:173--178, 1995.

\bibitem{rouetfeix}
J.-L.~Rouet and M.~R.~Feix.
\newblock A generalization of the Collatz problem. Building cycles and a
  stochastic approach.
\newblock \emph{Journal of Statistical Physics}, 107(5/6):1283--1298, 2002.

\bibitem{mleigh1987}
K.~R.~Matthews and G.~M.~Leigh.
\newblock A generalization of the Syracuse algorithm to $\mathbb{F}_{q}[x]$.
\newblock \emph{Journal of Number Theory}, 25:274--278, 1987.

\bibitem{hicks2008}
K.~Hicks, G.~L.~Mullen, J.~L.~Yucas, and R.~Zavislak.
\newblock A polynomial analogue of the $3N+1$ problem?
\newblock \emph{American Mathematical Monthly}, 115:615--622, 2008.

\bibitem{lagarias}
J.~C.~Lagarias.
\newblock The $3x+1$ problem and its generalizations.
\newblock \emph{The American Mathematical Monthly}, 92(1):3--23, 1985.

\bibitem{lagariasbook}
J.~C.~Lagarias, editor.
\newblock \emph{The Ultimate Challenge: The $3x+1$ Problem}.
\newblock American Mathematical Society, 2010.

\bibitem{terras}
R.~Terras.
\newblock A stopping time problem on the positive integers.
\newblock \emph{Acta Arithmetica}, 30(3):241--252, 1976.

\bibitem{tao}
T.~Tao.
\newblock Almost all orbits of the Collatz map attain almost bounded values.
\newblock \emph{Forum of Mathematics, Pi}, 10:e12, 2022.

\bibitem{sievepaper}
RobinCodes.
\newblock Formal verification of a Collatz residue sieve.
\newblock In the repository below, \texttt{Sieve (Residual Analysis)/paper/}.

\bibitem{repo}
RobinCodes.
\newblock Collatz --- sieves, generalised maps and orbital graphs.
\newblock \url{https://github.com/RobinCodes/Collatz}

\end{thebibliography}

\end{document}
